\documentclass[12pt,a4paper]{article}

%% --- Encoding and fonts ---
\usepackage[utf8]{inputenc}
\usepackage[T1]{fontenc}
%\usepackage{lmodern}  % disabled for review copy: newpx loaded by paper_style_preamble

%% --- Language ---
\usepackage[english]{babel}

%% --- Mathematics ---
\usepackage{amsmath, amssymb, amsthm}
\usepackage{mathtools}

%% --- Layout ---
\usepackage[margin=1in]{geometry}
\usepackage{setspace}
\onehalfspacing

%% --- Figures and tables ---
\usepackage{graphicx}
\usepackage{booktabs}
\usepackage{caption}
\usepackage{subcaption}
\usepackage{float}

%% --- Appendix ---
\usepackage[title,titletoc]{appendix}

%% --- Cross-references and links ---
\usepackage{hyperref}

%% --- Bibliography ---
\usepackage{natbib}
\bibliographystyle{apalike}

%% --- Draft tools ---
\usepackage{xcolor}
\newcommand{\todo}[1]{\textcolor{red}{\textbf{[TODO: #1]}}}
\newcommand{\note}[1]{\textcolor{blue}{\textit{[Note: #1]}}}

%% --- Paper style (splice review copy): Carrasco-2026 look ---
\input{MainDraft/paper_style_preamble}
\newcommand{\TODO}[1]{\textbf{[TBD#1]}}
%% REVIEW-COPY ONLY: four figure files and references.bib are not in this
%% repo (Overleaf assets). Render a placeholder instead of aborting.
\makeatletter
\let\spliced@incgfx\includegraphics
\renewcommand{\includegraphics}[2][]{%
  \IfFileExists{#2}{\spliced@incgfx[#1]{#2}}{%
    \IfFileExists{figures/#2}%
      {\spliced@incgfx[#1]{figures/#2}}%
      {\fbox{\parbox{0.85\linewidth}{\centering\ttfamily [figure file not in repo: \detokenize{#2}]}}}}}
\makeatother
% --- Custom Math Commands ---
\newcommand{\E}{\mathbb{E}}
\newcommand{\R}{\mathbb{R}}
\newcommand{\cA}{\mathcal{A}}          % asset space
\newcommand{\cH}{\mathcal{H}}          % productivity space
\newcommand{\cM}{\mathcal{M}}          % health space
\newcommand{\cF}{\mathcal{F}}          % formality space
\newcommand{\cQ}{\mathcal{Q}}          % pension-counter space
\newcommand{\ind}{\mathbf{1}}           % indicator function
\newcommand{\phij}{\phi_j}             % cohort weight
\newcommand{\jR}{j_R}                  % retirement age
% Type-specific shortcuts
\newcommand{\tL}{\theta_L}
\newcommand{\tH}{\theta_H}

%% --- Theorem environments ---
\newtheorem{assumption}{Assumption}
\newtheorem{proposition}{Proposition}
\newtheorem{definition}{Definition}

%% ====================================================================
\begin{document}

%% --- Title block ---
\title{\textbf{Spending Smarter under Demographic Pressure:}\\[6pt]
Fiscal Efficiency, Health, and Gender Dynamics\\
in Latin America\thanks{The work was financed with the support of the Latin America and the Caribbean Research Network of the Inter-American Development Bank. The opinions expressed in this publication are those of the authors and do not necessarily reflect the views of the Inter-American Development Bank, its Board of Directors, or the countries they represent.}}

\author{
Diego Ascarza-Mendoza\textsuperscript{1} \quad Hermilo Cort\'es\textsuperscript{1}\\[3pt]
Judith M\'endez\textsuperscript{1} \quad H\'ector J.\ Villarreal\textsuperscript{1}
\\[6pt]
\textsuperscript{1}Iniciativa para la Transici\'on Econ\'omica y Demogr\'afica (ITED),\\
Escuela de Gobierno y Transformaci\'on P\'ublica,\\
Tecnol\'ogico de Monterrey
}

\date{%
Prepared for the Inter-American Development Bank,\\
Fiscal Division\\[12pt]
\textit{Preliminary draft -- August 2026}\\
\textit{Please do not cite or circulate without permission}
}

\maketitle

%% --- Abstract ---
\begin{abstract}
Latin America is aging before it has grown rich, and the fiscal pressure this places on health and pension systems is sharpened by persistent labor informality and limited fiscal space. We develop a stochastic overlapping-generations model with endogenous health investment, agents heterogeneous in sex and skill, and a pay-as-you-go pension system whose contribution rate clears endogenously. Health capital depreciates with age, raises survival, productivity, and the amenity value of consumption, and is restored through medical spending, so it links longevity to the contribution rate through a fiscal amplification loop whose net sign is ambiguous. The model is designed for M\'exico, Costa Rica, and Panam\'a; this version reports the Mexican calibration, estimating the labor-efficiency profile and skill premium from ENIGH household surveys, the health-depreciation schedule and the health--earnings gradients from the Mexican health-and-aging study, and the remaining preference and health-technology parameters by classical minimum distance. The estimation block is not yet re-run under the paper's current preference specification, and the quantitative results reported below are the last complete solve of the model. Three findings follow. At M\'exico's projected 2050 demographics, the 2020 contribution rate and benefit formula open a pension financing gap that no finite stationary stock of public debt can close, narrowing the menu of feasible long-run responses below what the policy debate assumes. The calibrated capital stock is governed first by the persistence of idiosyncratic income risk and only second by skill composition. And a replacement-rate cut is regressive, while under unchanged policy the equilibrium contribution rate rises from 9.3\% to 15.9\% as the population ages.
\end{abstract}

\noindent\textbf{Keywords:} demographic transition; fiscal sustainability; health expenditure; overlapping generations; labor informality; Latin America

\noindent\textbf{JEL codes:} E62, H51, I13, J11, O17

%% ====================================================================
%% SECTION 1 — INTRODUCTION
%% ====================================================================
\section{Introduction}
\label{sec:intro}

Latin America is aging before it has grown rich. Its fertility decline is rapid, its formal labor markets remain incomplete, its fiscal space is tight, and its public health systems are chronically underfunded; the aging that took Europe roughly seventy years to accomplish will take the region approximately twenty-five, leaving far less time to legislate and implement the fiscal adjustments that demographic change demands \citep{WorldBankInformalizing2025}. The conjunction of these forces places the joint finance of health and pensions at the center of the region's fiscal prospects, and it does so under institutional conditions---informality, thin contributory bases, fragmented health systems---to which high-income-country experience does not speak directly. This paper builds, estimates, and runs experiments on a general-equilibrium model made for that environment; the extension to Costa Rica and Panam\'a lies ahead, and this version reports the Mexican calibration.

What the data say about health spending disciplines everything that follows. Mexican out-of-pocket medical spending is predominantly a response to realized illness rather than forward-looking prevention: classified item by item, unambiguously preventive care accounts for 2.0 percent of household medical spending, unambiguously curative care for 53.6 percent, and the remaining 44.5 percent is ambiguous, leaving measured spending majority curative under any assignment of the ambiguous block. The capacity to finance that response is, in addition, deeply unequal across households. The evidence for both statements is presented in the health-block section of the paper. % source: health-block section, pending
These facts matter at the outset because they fix what a structural model of health investment must confront when it is taken to Mexican data: measured spending is not the forward-looking investment flow that theory writes down, and the distance between the two turns out to be where the paper's central measurement result lives.

We build a stochastic overlapping-generations economy for M\'exico in which health is a continuous, endogenous capital stock that simultaneously raises survival, labor productivity, and the amenity value of consumption. Agents are heterogeneous in sex and skill; factor prices are endogenous; and the pay-as-you-go pension system's contribution rate clears endogenously, so demographic structure translates into the tax wedge inside the equilibrium rather than by assumption outside it. The model is estimated by classical minimum distance on ten moments, nine of them the paper's own estimates from Mexican micro-data: the labor-efficiency profile and the skill premium come from the ENIGH household surveys, and the health-depreciation schedule and the education- and sex-specific survival schedules from ENASEM, with the remaining preference, health-technology, and health-risk parameters disciplined by the level and age gradient of medical spending, its income elasticities, the skill premium, and survival anchors. Every experiment that follows is computed at the estimated baseline---the model's own general-equilibrium solve at the estimates---so each counterfactual perturbs an economy the data have disciplined, not an illustrative parameterization.

The first finding concerns the cost of demographic inaction. M\'exico's old-age dependency ratio is projected to more than double between 2020 and 2050, from 0.141 to 0.307 \citep{CONAPO2023}. Imposing the 2050 demographics with policy unchanged, the equilibrium contribution rate more than doubles, from 5.66 to 12.26 percent---and newborn welfare is essentially unchanged, a consumption-equivalent gain of 0.14 percent. The resolution of the apparent paradox is capital deepening: slower population growth tilts the stationary age distribution toward older, asset-rich cohorts, the annual return falls from 8.58 to 6.82 percent, the wage rises 12.4 percent, and the wage-indexed pension rises almost 20 percent, so for an agent who arrives holding nothing the induced wage and pension gains nearly offset the doubled wedge. The cost falls instead on cohorts already holding assets, whose remaining lifetimes are repriced at the lower return, and it rises steeply with age: one to two percent of remaining consumption at ages 25--29, six to eight percent at 40--44, and twelve to twenty-one percent at 65--69. Aging under unchanged policy is a redistribution across cohorts, not an aggregate welfare event.

The identity that closes a pay-as-you-go system leaves exactly three margins---the contribution rate, the replacement rate, or the boundary between contributors and beneficiaries---and the paper prices all three in general equilibrium under the same 2050 demographics. The fiscal arithmetic on each margin is exact and predictable before any equilibrium is solved: inaction takes the contribution rate to 12.26 percent; holding it at its 2020 value requires cutting the replacement rate from 0.40 to 0.185, a 54 percent benefit cut; and raising the retirement age from 65 to 70 lands the contribution rate at 8.78 percent. What general equilibrium adds is the welfare and incidence content the identity cannot supply. For newborns the ranking is strict and holds for every type: the benefit cut dominates the retirement-age rise, worth 8.3 and 3.1 percent of lifetime consumption against inaction respectively, and both dominate inaction. Where the reforms differ most is in whose old age pays. Inaction taxes wealth-holders through the return; the benefit cut taxes pension claims; the retirement rise taxes the extended work band. A planner reading this menu is choosing an incidence rather than a welfare level. One qualification governs all of these comparisons: they price destinations, not journeys---the equilibria are steady states, and the transition generations of a benefit cut, who contributed at the old rate and would retire under the new one, are not modeled.

The health block is the paper's methodological contribution, and its performance under estimation is the paper's central measurement result. The apparatus succeeds where it was expected to fail: once the felicity carries an estimated flow value of being alive, the model's willingness to pay for survival comes in close to an external benefit-transfer target---179 times annual output per capita against a reference of 152---having been off by orders of magnitude without it. It fails where it was expected to succeed: at parameter values that fit the level and age gradient of medical spending, both income elasticities, the skill premium, and the survival anchors, the endogenous investment channel generates almost none of the observed education--survival gradient---0.06 percentage points of a 2.09 percentage-point gap---while the estimator drives the auxiliary channels to corners and holds the health stock in a narrow band at its ceiling, spanning $[0.990, 1.000]$ across all types and ages. The diagnosis is a mismatch of objects, not a defect of implementation. Medical spending in the model is committed before the health shock is realized and is therefore prevention by construction; measured Mexican spending is predominantly restoration conditioned on a realized draw. A block that contains only the first, disciplined on moments generated by the second, cannot hold health away from its ceiling and cannot open a spending-driven survival gradient. We state this as a measurement result about the Grossman-descended investment block taken to emerging-economy micro-data.

The paper's position in the literature is defined by a combination rather than by any single ingredient. Quantitative macroeconomic work on health over the life cycle descends from the Grossman tradition through \citet{DalgaardStrulik2014}, who derive aging and mortality endogenously from health investment, and continues in the health-deficit framework of Hosseini, Kopecky, and Zhao, estimated on United States panel data; Jung and Tran build health investment and insurance into overlapping-generations settings of comparable ambition, again for the United States. What is distinctive here is the combination of an emerging economy; a health stock that carries survival, labor productivity, and the fiscal loop at once, so that health outcomes feed the pension system's contribution rate; and a negative result that only a fully disciplined calibration could have produced---an illustrative parameterization would never have been forced to confront the education--survival gradient it cannot generate. One further gap frames that result: the empirical literature on health spending in developing economies measures it as a response to illness, while the macroeconomic literature models it as forward-looking investment, and the two have not been reconciled; this paper's measurement result is what their collision looks like inside an estimated model.

The results suggest four extensions, and one paragraph suffices for them. The transition path between steady states would price the journeys that the steady-state comparisons cannot, and is the natural next computation. A formalization program---incorporating a fraction of the informal population into the contributory system and pricing the competing fiscal effects of a broader base against higher expenditure claims---is designed and awaits that apparatus. The extension to Costa Rica and Panam\'a would discipline the model against institutional architectures that differ instructively from M\'exico's. And the respecification of the health block, adding the restoration margin that the measurement result shows the data demand, is the binding item, because it converts the paper's negative finding into its constructive sequel. The remainder of the paper proceeds as follows. Section~\ref{sec:lit} reviews the related literature. Section~\ref{sec:model} presents the model and defines the stationary equilibrium. Section~\ref{sec:calibration} reports the two-step Mexican calibration, the estimation, and the calibrated baseline. Section~\ref{sec:healthblock} presents the evidence on the composition and distribution of Mexican medical spending.
Section~\ref{sec:experiments} develops the pension experiments at the estimated baseline, and Section~\ref{sec:conclusion} concludes.

%% ====================================================================
%% SECTION 2 — LITERATURE REVIEW
%% ====================================================================
\section{Literature Review}
\label{sec:lit}

This paper draws on five strands of literature: the economics of health financing and social insurance, the quantitative OLG literature on health and aging, the literature that models health as a capital stock built by forward-looking investment, the empirical work on informality and social security in Latin America, and the incorporation of gender heterogeneity into lifecycle fiscal models. They are not treated at equal depth. The third strand supplies the investment block whose capacity this paper measures against Mexican data, and the review is weighted toward establishing what that block has been asked to carry, and on what evidence.

%% --------------------------------------------------------------------
\subsection{Health Financing, Social Insurance, and the Tax Wedge}
\label{subsec:lit_health}

The economics of health insurance in developing economies builds on the insight that health shocks are both large and poorly insurable through private markets, generating welfare losses that pooling arrangements can substantially mitigate \citep{Arrow1963}. In Latin America the design question is not whether to pool health risk but through which institutional channel. Systems in the Bismarckian tradition link coverage to formal employment; systems in the Beveridgean tradition link it to residency or citizenship; and, as \citet{Levy2008} emphasizes, most countries in the region operate hybrids of the two, often with fragmentation, duplication, and unequal benefit packages.

\citet{Summers1989} supplies the valuation logic that organizes the modern literature: mandated contributions are priced into wages and generate no net distortion when workers value the benefits they purchase, and they function as a tax on formal labor precisely when workers do not---because benefits are not credible, or because their arrival is too distant. \citet{WorldBankInformalizing2025} formalize this logic as the Formalization Tax Rate, the net fiscal cost of formal employment inclusive of forgone non-contributory benefits, and show that it is highest for low-income workers, the group with the weakest attachment to the formal sector. Within this valuation literature, \citet{BaezaPackard2006} document the asymmetry most relevant to this paper's design: health coverage delivers immediate, tangible protection against catastrophic expenditure and is the social insurance benefit workers learn most about and value most highly, while pension benefits arrive far in the future and are discounted for institutional credibility, so the channel through which health coverage is financed has larger formalization consequences per unit of fiscal cost than pension design does.

The model in this paper takes a deliberate position in this landscape. Its pension benefit is flat and indexed to contemporaneous average earnings, so the model contains no benefit--contribution link at the individual level---a Beveridgean benefit inside a Bismarckian financing shell. The simplification buys tractability in the fiscal block, and it bounds what the paper can say about formalization incentives, which respond to exactly the individual-level link the model omits.

%% --------------------------------------------------------------------
\subsection{OLG Models of Fiscal Sustainability, Health, and Aging}
\label{subsec:lit_olg}

The quantitative OLG literature provides the methodological foundation for this paper. Since \citet{AuerbachKotlikoff1987}, the workhorse framework for analyzing the intergenerational fiscal effects of demographic change has been the multi-period OLG model with a government sector. The health-augmented extensions most relevant here include \citet{DeNardi2010}, who show that uncertain medical expenses are a primary driver of late-life savings among older households in the United States, and \citet{KuhnFeichtinger2015}, who analyze optimal health and retirement choices over the life cycle in a framework where health capital depreciates with age and can be partially restored through investment.

\citet{DalgaardStrulik2014} extend the Grossman framework of health as a depreciable capital stock to show how optimal aging and mortality patterns emerge endogenously from the interaction of health investment decisions with biological depreciation, yielding a structural interpretation of the Preston curve. \citet{Fioroni2010} analyzes optimal savings and health spending jointly over the life cycle, demonstrating that the interaction between precautionary saving motives and health investment generates consumption-savings profiles that differ substantially from models without endogenous health. \citet{FrankovicWrzaczek2020} study how medical progress interacts with overlapping-generations dynamics, showing that improvements in health technology alter not only individual life-cycle profiles but also aggregate capital accumulation and factor prices through general-equilibrium channels. \citet{YewZhang2018} introduce longevity externalities into a lifecycle-dynastic model, underscoring the importance of modeling health as an endogenous state variable rather than an exogenous process.

For the Latin American context, \citet{CortesEtAl2024} develop a partial-equilibrium OLG model for M\'exico in which agents differ by skill and face stochastic health transitions, health-dependent productivity, and age-varying out-of-pocket expenditure profiles. Their results---consumption gaps of 40--70 percent between good and poor health groups, and a 20-year gap in healthy longevity by skill level---are the empirical targets against which the present paper's model is disciplined and tested. The primary extension in this work is the move to general equilibrium with endogenous wages, an explicit pension system with an endogenous contribution rate, and a continuous health capital stock with diminishing-returns investment technology. \citet{Lim2020} analyzes the case of South Korea, showing that public provision of health insurance alters aggregate saving behavior in an OLG framework with endogenous health risk---a mechanism that resonates with the fiscal loop studied here, in which health outcomes and the pension system's finances are jointly determined.

%% --------------------------------------------------------------------
\subsection{Health as Capital: What the Investment Block Is Asked to Carry}
\label{subsec:lit_longevity}

Every model in the lineage just described rests on a common assumption, inherited from Grossman, and it should be stated as an assumption rather than absorbed as a convention. Health is a depreciable capital stock; medical spending is forward-looking investment in that stock, committed before uncertainty about illness resolves; and the stock so built governs survival---in the richer versions, labor productivity as well---so that spending today purchases longevity and earnings tomorrow. \citet{DalgaardStrulik2014} and the health-deficit line derive aging itself from this accumulation problem. \citet{Carrasco2026} shows the same block operating as a macroeconomic transmission channel: households that invest in health live longer, longer lives strengthen the life-cycle saving motive, and the aggregate capital stock and the equilibrium interest rate move in consequence. In all of these frameworks the investment assumption is not an approximation at the margin; it is the engine.

The evidentiary base on which the block has been estimated is narrower than its use. These models are calibrated or estimated overwhelmingly on high-income data---chiefly United States panels, as in the health-deficit estimations of Hosseini, Kopecky, and Zhao---in economies where insurance coverage is broad, out-of-pocket payments are a modest share of total health spending, and the data do not force a distinction between spending that builds the stock and spending that repairs it after a shock. Whether the same block survives contact with an economy in which health spending is largely out of pocket, largely triggered by realized illness, and deeply unequal in its financing is a question the literature has not settled, because no estimated general-equilibrium model has been taken to such an economy with moments disciplined enough to force it.

The question has a precise form, and it is open. The empirical literature on health expenditure in developing economies measures spending as a response to illness---care sought and financed after a health draw is realized. The macroeconomic literature models spending as forward-looking investment committed in advance. These are not variants of one description: they imply different timing, different exposure to expenditure risk, and different capacity to generate gradients in survival from gradients in spending. The two literatures have not been reconciled, and the reconciliation is not a theoretical exercise---it is a measurement question about what observed spending is. The estimation in this paper is where the collision becomes measurable: a fully disciplined calibration of the investment block to Mexican micro-data confronts, moment by moment, spending data generated largely by the object the block does not contain.

%% --------------------------------------------------------------------
\subsection{Informality, Formalization Incentives, and Social Security}
\label{subsec:lit_informality}

The empirical literature documents a robust negative relationship between social security contribution rates and formal employment, particularly at the lower end of the wage distribution: \citet{Kugler2017} estimate payroll-tax elasticities for Colombia between 0.4 and 1.0, and the expected return on contributions matters as much as their cost---an increase in the vesting period required for the minimum pension reduced Colombian formal employment by approximately 3 percentage points. The non-contributory channel operates in the opposite direction: benefits available outside formal employment raise the opportunity cost of formalizing, and the estimated informality effects of the non-contributory health coverage expansions in M\'exico and Colombia range from substantial to negligible depending on program design and institutional context \citep{BoschCampos2014, AteridoEtAl2011, AzuaraMarinescu2013, CamachoConoverHoyos2013, CalderonMarinescu2012}. Behind both findings sits what \citet{PerryEtAl2007} read as a mass opting out of state institutions. \citet{WorldBankInformalizing2025} organize the resulting reform problem along three dimensions that must move together---the tax wedge firms can absorb, the social security wealth workers perceive, and the credibility of enforcement---on the compressed timeline that frames this paper: the aging that took Europe roughly 70 years will take Latin America approximately 25.

This paper's baseline is calibrated to the formal-sector population, and the choice is material: the model's dependency ratio is demographic rather than contributor-based, so its PAYG arithmetic abstracts from the contribution-base margin this literature studies. A formalization program---incorporating a fraction of the informal population into the contributory system and pricing the fiscal effect of a broader base against the higher expenditure claims of newly formalized workers---is designed and left as an extension.

%% --------------------------------------------------------------------
\subsection{Gender Heterogeneity in Lifecycle Fiscal Models}
\label{subsec:lit_gender}

The incorporation of gender into quantitative OLG models remains at an early stage. \citet{FehrKindermann2018} provide the reference framework, a two-person household whose female member's labor supply responds to childcare, parental leave, and joint-versus-individual taxation---but in an environment without health as a state variable, without survival risk, and without a fiscal loop. Embedding a bargaining household of that kind in a model with continuous health capital, endogenous survival, and endogenous factor prices would compound the curse of dimensionality with the fixed-point problem in prices, and to our knowledge it has not been attempted in the published literature. The approach here is different, and it is implemented rather than proposed: the agent space carries four types throughout---two skill levels for each sex---calibrated with sex-specific earnings trajectories and survival schedules, so that the sex-differentiated profiles embed in reduced form the cumulative effects of labor market discrimination, career interruptions, and gendered sorting that a structural household model would generate endogenously, and every experiment in the paper reports welfare and incidence for all four types. The structural household model remains a natural extension; it is not a prerequisite for the fiscal decomposition by sex that the paper's questions require.

Section~\ref{sec:model} presents the framework that integrates these strands. It extends the ITED baseline \citep{CortesEtAl2024} to general-equilibrium factor pricing, a continuous health capital stock with age-dependent depreciation, and a contribution-rate rule that links demographic aging to the fiscal wedge on formal labor; Section~\ref{sec:healthblock} then presents the evidence on Mexican medical spending against which the investment assumption of Section~\ref{subsec:lit_longevity} is put to the test.

\bigskip
%% --------------------------------------------------------------------

\section{Model}
\label{sec:model}

\subsection{Environment and Demographics}
\label{subsec:environment}

Time is discrete. The economy is populated by $J$ overlapping generations indexed by $j = 1, \ldots, J$, where $j = 1$ corresponds to age~20 and each model period spans five years. Households enter economic life at $j = 1$, retire exogenously at age $j_R$, and may live until $j = J$ (age~100).

\paragraph{Agent types.}
Each individual is indexed by sex $g \in \{m, f\}$ (male, female) and by a permanent skill type $\theta \in \{\theta_L, \theta_H\}$ (low-skilled, high-skilled), where $\theta_L$ and $\theta_H$ denote the log-productivity intercepts of the two skill groups (see equation~\eqref{eq:labor_productivity}). The two skill groups correspond to education groups, and we use the two terms interchangeably. Sex and skill are both drawn at birth and remain fixed over the life cycle. Sex enters the model through two channels: the age-efficiency profile $e^{g,\theta}_j$ that shapes earnings, and the baseline survival schedule $\psi^{\circ\, g,\theta}_j$ that shapes longevity. Skill enters through three: the log-intercept $\theta$; the skill-specific component of the age-efficiency profile $e^{g,\theta}_j$, which lets the life-cycle growth of earnings differ across education groups; and the skill-specific component of the baseline survival schedule $\psi^{\circ\, g,\theta}_j$ in the survival modulation~\eqref{eq:survival_modulation}, which carries mortality differences across education groups that do not operate through health investment. The health-productivity slope $\varrho$ in~\eqref{eq:labor_productivity} is common to both groups. Let $\pi^{g,\theta}_1$ denote the birth share of agents of sex $g$ and skill $\theta$, with $\sum_{g,\theta} \pi^{g,\theta}_1 = 1$. The population grows at gross rate $1 + n_p$, so the absolute mass of newborns of type $(g,\theta)$ at calendar time $t$ is $\pi^{g,\theta}_1 \cdot N_t$.

\paragraph{Idiosyncratic risk.}
Households face three sources of individual uncertainty: mortality risk, which depends on health; idiosyncratic labor-productivity risk, summarized by an AR(1) shock $\eta$; and health risk, an age-dependent probability of losing a fixed amount of health (Section~\ref{subsubsec:health}). Health $h$ is therefore a controlled stochastic state: endogenous through chosen medical spending, risky through the shock. The productivity and health shocks are independent, and there is no aggregate uncertainty.

\paragraph{Individual state.}
The individual state of a household at the beginning of a period is
\begin{equation}
\label{eq:state}
s = (a,\, h,\, \eta),
\end{equation}
where $a \in \mathcal{A} = [0, \bar{a}]$ are beginning-of-period assets, $h \in \mathcal{H}_h = (0, \bar{h}]$ is health capital, and $\eta \in \mathcal{H}_\eta = \{\eta_1, \ldots, \eta_{N_\eta}\}$ is the idiosyncratic productivity component. Sex $g$, skill type $\theta$ and age $j$ are deterministic indices that label the value and policy functions.

\paragraph{Cohort weights.}
Let $\phi_j^{g}(\theta)$ denote the stationary mass of age-$j$ households of sex $g$ and skill $\theta$, normalized so that $\sum_g \sum_\theta \sum_{j=1}^{J} \phi_{j}^{g}(\theta) = 1$ with newborn weights $\phi_{1}^{g}(\theta) = \pi^{g,\theta}_1 / Z$, where
\[
Z \;=\; \sum_{g}\sum_{\theta}\sum_{j=1}^{J} \pi^{g,\theta}_1 \prod_{i=2}^{j} \Bigl[\bar\psi_i^{g}(\theta) \big/ (1+n_p)\Bigr]
\]
is the normalizing constant, with the convention that the empty product at $j=1$ equals one. Birth shares are sex- and skill-specific and need not coincide across types. The weights satisfy the forward recursion
\begin{equation}
\label{eq:cohort_weights}
\phi_{j+1}^{g}(\theta) \;=\; \frac{\bar{\psi}_{j+1}^{g}(\theta)}{1 + n_p}\, \phi_{j}^{g}(\theta), \qquad j = 1, \ldots, J-1,
\end{equation}
where the population-weighted average survival is
\begin{equation}
\label{eq:avg_survival}
\bar{\psi}_{j+1}^{g}(\theta)
\;=\;
\frac{1}{\phi_{j}^{g}(\theta)} \int \sum_{k \in \{0,1\}} \Pr(k \mid j)\;
\psi_{j+1}^{g,\theta}\bigl(h'_k(s;\, j, g, \theta)\bigr)\, d\mu_j^{g,\theta}(s),
\end{equation}
with $h'_k(s;\, j, g, \theta)$, $k \in \{0,1\}$, the two nodes of the health lottery~\eqref{eq:health_lom_shock} evaluated at the optimal medical spending $m_j(s;\, g, \theta)$, $\Pr(1 \mid j) = \pi^h_j$ the shock arrival probability of~\eqref{eq:shock_prob}, and $\mu_{j}^{g,\theta}$ the unconditional age-$j$ population measure for sex $g$ and skill $\theta$, so that $\int d\mu_{j}^{g,\theta} = \phi_{j}^{g}(\theta)$. Terminal-age survival is zero: $\psi_{J+1}^{g,\theta}(\cdot) \equiv 0$.

\paragraph{Within-period timing.}
The timing within each period is:
\begin{enumerate}
\item The household observes its current state $(a, h, \eta)$.
\item Working-age households choose $(c, \ell, a', m)$. Retired households choose $(c, a', m)$. Medical spending is committed before the health shock is realized, so investment cannot condition on the shock's arrival: medical spending in the model is a preventive object by construction, building the stock in advance of illness rather than restoring it afterwards.
\item End-of-period assets are $a'$. Post-investment health is $\tilde h = (1-\delta_j^h)\, h + H_j(m, h)$, with $H_j$ the deficit-scaled production function~\eqref{eq:health_production_deficit}. The health shock is then realized: with probability $\pi^h_j$ the household loses $\varsigma$ units of health, and next-period health $h'$ is the realized node of the lottery~\eqref{eq:health_lom_shock}. Survival to next period is realized with probability $\psi_{j+1}^{g,\theta}(h')$, evaluated at the realized $h'$. Assets of agents who die are voided: they neither enter aggregate capital next period nor flow to survivors. This rules out an accidental-bequest channel and isolates the life-cycle savings motive; the aggregate resource implication is tracked by the void term $\Lambda_{\mathrm{void}}$ introduced in equation~\eqref{eq:leakage_voided}.
\end{enumerate}


\subsection{Preferences}
\label{subsec:preferences}

Individuals rank streams of consumption, leisure, and health by expected discounted lifetime utility. The lifetime objective of a newborn household of sex $g$ and type~$\theta$ is
\begin{equation}
\label{eq:lifetime_utility}
\E_1 \left[\sum_{j=1}^{J} \beta^{j-1} \left(\prod_{i=2}^{j} \psi_{i}^{g,\theta}(h_{i-1}')\right) u(c_j,\, \ell_j,\, h_j)\right],
\end{equation}
where $\beta \in (0,1)$, the survival probability $\psi_{i}^{g,\theta}(h_{i-1}')$ is the probability of being alive at age $i$ given that realized next-period health entering age $i$ is $h_{i-1}' = h_i$, and the expectation is taken over the random path $\{\eta_j\}_{j=1}^{J}$ of the idiosyncratic shock generated by~\eqref{eq:ar1} and over the realizations of the health lottery~\eqref{eq:health_lom_shock}.

The period utility function is a GHH aggregator scaled by a health amenity:
\begin{equation}
\label{eq:utility}
u(c, \ell, h) \;\equiv\; \Omega(h) \left[ \frac{\bigl(c - v(\ell)\bigr)^{1-\gamma} - 1}{1 - \gamma} \;+\; \chi \right],
\qquad
v(\ell) \equiv \Psi \frac{\ell^{1+\nu_\ell}}{1+\nu_\ell},
\end{equation}
\begin{equation}
\label{eq:amenity}
    \Omega(h) \equiv
\begin{cases}
\exp\!\left(\Xi \dfrac{h^{1-\xi} - 1}{1-\xi}\right) & \xi \neq 1,\\[8pt]
h^{\Xi} & \xi = 1,
\end{cases}
\end{equation}

with $\Psi > 0$, $\Xi \geq 0$, $\nu_\ell > 0$ (the inverse Frisch elasticity, written with an $\ell$ subscript to avoid collision with the productivity function $\nu_j$), $\gamma > 0$, and $\chi > -1$. Here $c \geq 0$ is consumption, $\ell \in [0, 1]$ is hours worked, $1-\ell$ is leisure, and $h$ is the contemporaneous stock of health capital, which simultaneously affects survival, productivity, and the amenity value of consumption. The amenity $\Omega$ is nondecreasing in $h$ and strictly positive for every admissible $(\Xi,\xi)$, and strictly increasing whenever $\Xi > 0$, so health is unambiguously a good wherever the bracket in~\eqref{eq:utility} is positive.

Two conventions fix the health block. The stock is measured as one minus a health deficit, so it is bounded above, and we place its upper endpoint at $\bar h = 1$: $h \in \mathcal{H}_h = (0, \bar h]$, with $h = \bar h$ denoting full health. This choice of units is not free. Health enters labor productivity through $\varrho\, h$ in~\eqref{eq:labor_productivity}, survival through $\psi^{g,\theta}_{j+1}(h)$ in~\eqref{eq:survival_modulation}, and its own law of motion~\eqref{eq:health_lom_shock}, so a common rescaling of $h$ would have to be undone simultaneously in all three objects, and the measurement of the stock is what pins the scale down. (The domain is open at the origin because $\Omega(h) \downarrow 0$ as $h \downarrow 0$ whenever $\xi \geq 1$, so $h = 0$ is excluded rather than assigned a utility level.) The normalization $\Omega(\bar h) = 1$, by contrast, is free: the amenity multiplies period utility and death carries a continuation value of zero, so replacing $\Omega$ by $\lambda \Omega$ for any $\lambda > 0$ leaves preferences unchanged and $\lambda$ is not identified. The $-1$ in the exponent of~\eqref{eq:amenity} is what imposes $\lambda = 1$; since $(h^{1-\xi} - 1)/(1-\xi) \to \log h$ as $\xi \to 1$, the same subtraction makes the second branch of~\eqref{eq:amenity} the continuous limit of the first rather than a separate assumption. With the scale fixed, $\Xi$ governs the gradient of the amenity in health and $\xi$ its curvature.

The multiplicative amenity is a substantive specification choice, and it is where the utility side of the paper's mechanism lives. Scaling the GHH aggregator by $\Omega(h)$ makes the marginal utility of consumption increasing in health, $u_{ch} = \Omega'(h)(c - v(\ell))^{-\gamma} \geq 0$, with strict inequality whenever $\Xi > 0$, so health and consumption are complements in period utility. The property that matters for behavior, however, is that the marginal rate of substitution between health and consumption, $u_h/u_c$, is increasing in consumption at every state the household ranks above not being alive: the willingness to pay for health rises with income. Had the amenity entered the composite additively, that ratio would be independent of consumption---health and consumption would be perfect substitutes inside the composite, and the willingness to pay for health would carry no income gradient at all. The multiplicative form therefore does not merely tilt a gradient; it creates one where the additive form has none, and that gradient is the utility-side channel through which income inequality is transmitted into health inequality. Appendix~\ref{app:amenity} derives both marginal rates of substitution and records the comparison. Two further properties come with the specification. Because the amenity multiplies the aggregator, $u_c \to \infty$ as $c \downarrow v(\ell)$, so consumption satisfies an Inada condition that an additive amenity can forfeit; and because $\Omega(h)$ scales $u$ uniformly, it cancels from the intratemporal labor margin, so the wealth-effect-free property of GHH survives on the labor-supply decision.

The constant $\chi$ is not a free normalization either. Death carries a continuation value of zero in~\eqref{eq:lifetime_utility}, so the level of $u$ is economically meaningful; and because the $-1$ in the numerator of~\eqref{eq:utility} itself contributes the constant $1/(\gamma-1)$ to the bracket, the object that is identified---and that constitutes the flow value of being alive---is the sum $\chi + 1/(\gamma-1)$ rather than $\chi$ alone. That sum determines the household's willingness to pay to reduce mortality risk; states at which the bracket in~\eqref{eq:utility} is negative are states the household ranks below not being alive. Separating $\chi$ from $\Omega(h)$ assigns to distinct parameters the two roles the amenity would otherwise play at once---$\Xi$ and $\xi$ govern how utility varies with health, while $\chi$ governs the level of utility from being alive at all---and because $\chi$ is absent from $u_c$, it shifts the value function, and hence the value of a statistical life and the demand for survival-improving medical spending, without disturbing the consumption Euler equation or the labor margin.

For retirees ($j \geq j_R$), $\ell_j = 0$ and period utility reduces to $u(c, 0, h) = \Omega(h)\left[\frac{c^{1-\gamma} - 1}{1-\gamma} + \chi\right]$.


\subsection{Technology}
\label{subsec:technology}

A representative, competitive firm produces the single consumption good using capital $K$ and aggregate effective labor $L$ according to the Cobb--Douglas technology
\begin{equation}
\label{eq:production}
Y = A\, K^\alpha\, L^{1-\alpha}, \qquad \alpha \in (0,1),
\end{equation}
with $A > 0$ total factor productivity. Capital depreciates at rate $\delta \in (0,1)$. Constant returns to scale imply that the firm's first-order conditions depend only on the ratio $K/L$ and are therefore invariant to the choice of units, so we write $K$ and $L$ in the stationary per-cohort-mass units of Section~\ref{subsec:markets} throughout.

\paragraph{Firm's problem.}
The firm takes factor prices $(r, w)$ as given and maximizes static profits:
\begin{equation}
\label{eq:firm_problem}
\max_{K, L} \bigl\{A K^\alpha L^{1-\alpha} - (r+\delta)\,K - w\,L\bigr\}.
\end{equation}
The first-order conditions yield
\begin{align}
r &= \alpha A \left(\frac{K}{L}\right)^{\alpha-1} - \delta, \label{eq:foc_r} \\
w &= (1-\alpha)\, A \left(\frac{K}{L}\right)^{\alpha}. \label{eq:foc_w}
\end{align}

\paragraph{Factor prices.}
Both $r$ and $w$ are determined endogenously: given aggregate capital $K$ and aggregate effective labor $L$, the capital--labor ratio pins down both prices simultaneously through~\eqref{eq:foc_r}--\eqref{eq:foc_w}. The interest rate is the object that adjusts to clear the capital market.


\subsection{Exogenous Processes}
\label{subsec:exogenous}

\subsubsection{Mortality}
\label{subsubsec:mortality}

Conditional survival probabilities depend on age, sex, skill, and end-of-period health. The probability of surviving from age $j$ to age $j+1$, realized at the end of age $j$, is
\begin{equation}
\label{eq:survival_modulation}
\psi_{j+1}^{g,\theta}(h')
\;=\;
\psi^{\circ\, g,\theta}_{j+1}\,
\min\Bigl\{1,\;\; \varpi \;+\; (1-\varpi)\, h'\Bigr\},
\end{equation}
where $h'$ is the realized next-period health stock determined by the health lottery~\eqref{eq:health_lom_shock}, $\psi^{\circ\, g,\theta}_{j+1} \in (0,1]$ is a sex- and education-specific baseline survival schedule---a primitive of the environment, exactly as the sex split is---and $\varpi \in (0,1)$ is the survival floor, the fraction of baseline survival retained as health approaches its lower bound. Under the normalization $\bar h = 1$ of Section~\ref{subsec:preferences}, the modulation equals one at full health, so $\psi^{\circ\, g,\theta}_{j+1}$ is survival at full health; the $\min$ caps the modulation at one, so full health cannot raise survival above the baseline schedule. The modulation is piecewise linear in $h'$: the survival return to health, $(\psi^{g,\theta}_{j+1})'(h') = \psi^{\circ\, g,\theta}_{j+1}(1-\varpi)$, is constant wherever the cap is slack and zero where it binds, and this derivative is the object that prices survival in the medical-spending condition~\eqref{eq:foc_m} of Appendix~\ref{app:focs}. Terminal-age survival is zero, $\psi_{J+1}^{g,\theta}(\cdot) \equiv 0$, so death by age $j = J$ is certain. Section~\ref{sec:calibration} carries the measurement of the schedules and the floor.

\paragraph{Two mortality channels, one decomposition.}
Health enters~\eqref{eq:survival_modulation} through the term $(1-\varpi)h'$, and that term is the investment channel: medical spending raises next-period health and with it survival, and it is the only survival margin a household controls. The education split of the baseline schedule is a non-investment channel. Education groups differ in mortality for reasons that do not run through the purchase of medical care---differential exposure to occupational hazards, health behaviors, and access to care---and the schedule carries those differences as a measured input, so that the investment channel is not forced to absorb mortality variation it cannot generate. The model's education gap in survival at any age therefore decomposes into a component generated by the education gap in health investment---the baseline schedule multiplied by $(1-\varpi)\,\Delta_j$ wherever the cap is slack, where $\Delta_j$ denotes the education gap in mean next-period health averaged over the population distribution and the shock lottery---and a component inherited from the education gap in the baseline schedules themselves, measured rather than generated. Every statement the paper makes about the income-to-health mechanism refers to the first component; a matched total gap is not by itself evidence on the mechanism, since the baseline schedules carry whatever the investment channel does not produce. Section~\ref{sec:calibration} reports the split.

\subsubsection{Labor Productivity}
\label{subsubsec:productivity}

Effective labor productivity of a working-age household ($j < j_R$) is
\begin{equation}
\label{eq:labor_productivity}
\nu_j(h, \eta;\, g,\theta) \;=\; e_{j}^{g,\theta} \exp\bigl(\theta + \eta + \varrho\, h\bigr),
\end{equation}
where $e_{j}^{g,\theta} > 0$ is a deterministic sex- and skill-specific age-efficiency profile, $\theta \in \{\theta_L, \theta_H\}$ is the log-productivity intercept of the skill group, and $\varrho \geq 0$ is the elasticity of productivity with respect to health capital, common to both skill groups. Skill therefore shifts both the level of earnings, through $\theta$, and their life-cycle shape, through $e^{g,\theta}_j$: education groups differ in their returns to experience, and the profile is the object that carries that difference. Health enters symmetrically across groups, through the single slope $\varrho$; any skill difference in the level or growth of earnings is attributed to the intercept and the profile, not to a differential return to health. Productivity is zero for retirees ($j \geq j_R$).

The idiosyncratic component $\eta$ follows the AR(1) process
\begin{equation}
\label{eq:ar1}
\eta_{j+1} = \rho\, \eta_j + \varepsilon_{j+1}, \qquad \varepsilon_{j+1} \sim \mathcal{N}(0, \sigma_\varepsilon^2),
\end{equation}
with $|\rho| < 1$.

\subsubsection{Health Formation, Health Investment, and Health Risk}
\label{subsubsec:health}
\label{subsubsec:health_risk}

Health evolves through three forces: deterministic depreciation at the age-dependent rate $\delta^h_j \in (0,1)$, repair through chosen medical spending $m$, and a stochastic health shock. We take the three in turn and close with the deterministic special case that the specification nests.

\paragraph{Deficit-scaled health production.}
Medical spending produces health according to
\begin{equation}
\label{eq:health_production_deficit}
H_j(m, h) \;=\; \overline{H}_j\, m^{\zeta_h}\,
   \underbrace{\frac{\bar h - (1-\delta^h_j)\, h}{\bar h}}_{\text{remaining deficit}},
\qquad \zeta_h \in (0,1),
\end{equation}
where $\overline{H}_j > 0$ is an age-dependent scale and the subscript on the curvature $\zeta_h$ keeps it distinct from the curvature $\xi$ of the amenity in~\eqref{eq:amenity}. The return to medical spending is proportional to the health deficit remaining after within-period depreciation: a household already at $\bar h$ gains nothing from further spending and one far below it gains a great deal. The deficit is measured after depreciation, so it is the gap actually available to be repaired. Deficit scaling is the clinically natural property of care---treatment restores a loss, and there is less to restore the healthier the patient---and it disciplines the model in a way constant-returns repair cannot; the deterministic limit discussed at the end of this section makes the comparison precise.

Write $\tilde h \equiv (1-\delta^h_j)\,h + H_j(m,h)$ for post-investment health. Two consequences of~\eqref{eq:health_production_deficit} follow. First, $\bar h - \tilde h = (\bar h - \overline{H}_j m^{\zeta_h})(\bar h - (1-\delta^h_j)h)/\bar h$, so $\tilde h < \bar h$ for every admissible $h$ if and only if $\overline{H}_j m^{\zeta_h} < \bar h$. On that region the upper bound is an asymptote, approached but not attained, and the cap at $\bar h$ is retained but slack. The condition restricts the chosen $m$ rather than the parameters, and $\tilde h$ is unbounded in $m$, so the cap can still bind at sufficiently large spending. Second, health production depends on the health stock, so
\begin{equation}
\label{eq:dh_dh}
\frac{\partial \tilde h}{\partial h} \;=\; (1-\delta^h_j)\left[1 - \frac{\overline{H}_j m^{\zeta_h}}{\bar h}\right],
\qquad
\frac{\partial H_j}{\partial m} \;=\; \zeta_h\, \overline{H}_j\, m^{\zeta_h - 1}\,\frac{\bar h - (1-\delta^h_j)h}{\bar h},
\end{equation}
and these are the derivatives that enter the envelope conditions~\eqref{eq:env_h} and~\eqref{eq:env_h_ret} and the medical-spending condition~\eqref{eq:foc_m} of Appendix~\ref{app:focs}. The marginal product of medical spending falls as health rises.

\paragraph{Stochastic health.}
Health is risky. With probability $\pi^h_j$ the household suffers a health shock at age $j$ and loses $\varsigma > 0$ units of health, so next-period health is the two-point random variable
\begin{equation}
\label{eq:health_lom_shock}
h' =
\begin{cases}
\min\{\tilde h,\ \bar h\}
  & \text{with probability } 1 - \pi^h_j \quad \text{(no shock)},\\[4pt]
\min\bigl\{\max\{\tilde h - \varsigma,\ \underline h\},\ \bar h\bigr\}
  & \text{with probability } \pi^h_j \quad \text{(shock)},
\end{cases}
\end{equation}
with $\underline h > 0$ the floor of the shock lottery: an acute loss can
displace health below any level reached by depreciation alone, but not below
$\underline h$. (Health thus lives on $[\underline h, \bar h]$ once shocks
operate; the open lower endpoint of $\mathcal{H}_h$ in~\eqref{eq:state} is the
domain on which preferences and technology are defined, and $\underline h$ is
the smallest value the law of motion visits.) We index the two nodes by $k \in \{0,1\}$, writing $h'_k$ for the realization and $\Pr(k \mid j)$ for its probability, with $\Pr(1 \mid j) = \pi^h_j$. The arrival probability inherits its age profile from the depreciation schedule rather than introducing new age-varying parameters,
\begin{equation}
\label{eq:shock_prob}
\pi^h_j \;=\; \min\!\left\{\,\bar\pi^h\,
   \frac{\delta^h_j}{\overline{\delta^h}}, \;\; \overline{\pi}^{\,\max} \right\},
\qquad
\overline{\delta^h} \equiv \frac{1}{j_R-1}\sum_{i=1}^{j_R-1} \delta^h_i ,
\end{equation}
so that only the scalar $\bar\pi^h$, the mean arrival rate over working ages, is free: the ages at which health depreciates fastest are also the ages at which acute losses arrive most often. The cap $\overline{\pi}^{\,\max}$ prevents the oldest bands, where $\delta^h_j$ is largest, from degenerating into certain arrival. It comes at a price that should be stated rather than discovered: wherever the cap binds, incidence is a constant rather than a profile inherited from $\delta^h_j$, and once every working-age band has reached the cap, the working-age mean incidence ceases to respond to $\bar\pi^h$ at all (with $\delta^h_j$ increasing in age, the youngest band is the last to cap). The health block carries exactly two risk parameters, $\bar\pi^h$ and $\varsigma$.

\paragraph{Incidence is exogenous, and this is essential.}
$\pi^h_j$ and $\varsigma$ depend on age alone, never on income, assets, or skill. Exogenous incidence is not an approximation of convenience but the identifying restriction that makes the exercise meaningful. Rich and poor households face the same risk of falling ill; they differ only in what they can do about it. Any health--income gradient the model produces is therefore earned through the response to health shocks rather than assumed through their arrival. Allowing incidence to depend on income would place the paper's conclusion among its premises. The same discipline governs mortality: the baseline schedule in~\eqref{eq:survival_modulation} is fixed at birth with sex and skill, so within a type the only income-varying survival object is health itself.

\paragraph{The deterministic limit.}
Setting $\bar\pi^h = 0$ makes the lottery~\eqref{eq:health_lom_shock} degenerate at $\min\{\tilde h, \bar h\}$: deficit-scaled repair without risk. Deleting in addition the deficit factor from~\eqref{eq:health_production_deficit} recovers the constant-returns rule
\begin{equation}
\label{eq:health_formation}
h' = \min\bigl\{(1 - \delta_j^h)\, h + \overline{H}_j\, m^{\zeta_h},\; \bar h\bigr\},
\end{equation}
familiar from deterministic health-capital models. The baseline is not that rule, for a reason worth stating. Under~\eqref{eq:health_formation} the return to medical spending is independent of the health stock, so restoring health from just below the ceiling is exactly as productive as restoring it from far below; when depreciation at working ages is small, any household that spends at all is pushed back to $\bar h$, working-age mass concentrates at the ceiling, and within-cohort health dispersion degenerates. A mechanism transmitting income inequality into health inequality then operates on a variable with almost no variation. The two departures are complements: deficit scaling makes the marginal product of care fall with health, so heterogeneous households stop short of the ceiling at heterogeneous points, and the shock keeps the cross-section from being reabsorbed into it. The model nests~\eqref{eq:health_formation} as the stated boundary case.


\subsection{Government}
\label{subsec:government}

The government runs two separate balance sheets: a general fiscal block of taxes, public goods and one-period debt, and a self-financed pay-as-you-go (PAYG) pension block. Pension contributions $\tau^p$ enter the general budget only indirectly, through the labor-tax wedge faced by workers; they are not a revenue item for the general government.

\subsubsection{Fiscal Policy}
\label{subsubsec:fiscal}
The government taxes consumption at rate $\tau^c$, labor income at rate $\tau^\omega$, capital income at rate $\tau^k$, and medical spending at rate $\tau^m$. It purchases public goods $G = g_y Y$ and may issue one-period debt $B$. All aggregate quantities ($C$, $K$, $L$, $M$, $Y$, $G$, $B$) are in the stationary per-cohort-mass units of Section~\ref{subsec:markets}; $C$, $L$, $M$, $Y$, $G$ are per-period flows and $K$, $B$ are stocks. The flow budget is
\begin{equation}
\label{eq:gov_budget}
\tau^c C + \tau^m M + \tau^\omega w L + \tau^k r K + B' \;=\; G + (1+r)B,
\end{equation}
where $C$ is aggregate consumption, $K$ the aggregate capital stock, $M$ aggregate health investment, and $B'$ next-period government debt.

\paragraph{Fiscal closure.}
$(\tau^c, \tau^\omega, \tau^k, \tau^m, g_y)$ are exogenous policy parameters. In the stationary equilibrium $B' = B$, so debt service $rB$ equals the primary surplus:
\begin{equation}
\label{eq:gov_stationary}
r\, B \;=\; \tau^c\, C + \tau^\omega\, w\, L + \tau^k\, r\, K + \tau^m\, M - G.
\end{equation}
$B$ is the residual that closes the budget.

\subsubsection{Pension System}
\label{subsubsec:pension}
A PAYG pension is financed by a payroll tax $\tau^p$ on workers' labor income. The system balances each period:
\begin{equation}
\label{eq:pension_balance}
\tau^p\, w\, L \;=\; \overline{pen}\, N^R,
\end{equation}
where $N^R = \sum_g \sum_\theta \sum_{j=j_R}^{J} \phi_j^{g}(\theta)$ is the mass of retirees. The pension benefit is the replacement-rate parameter $\kappa$ times average per-worker labor income:
\begin{equation}
\label{eq:pension_benefit}
\overline{pen} \;=\; \kappa \cdot \frac{w L}{N^W},
\end{equation}
where $N^W = \sum_g \sum_\theta \sum_{j=1}^{j_R - 1} \phi_j^{g}(\theta)$ is the mass of working-age individuals.

\paragraph{Pension closure.}
The replacement rate $\kappa$ is an exogenous policy parameter and the contribution rate $\tau^p$ adjusts endogenously to balance the system each period. Substituting~\eqref{eq:pension_benefit} into~\eqref{eq:pension_balance} yields
\begin{equation}
\label{eq:pension_contribution}
\tau^p \;=\; \kappa\, \frac{N^R}{N^W},
\end{equation}
so the contribution rate is the replacement rate scaled by the system dependency ratio $N^R/N^W$. This pins $\tau^p$ down once the stationary distribution is known, and it is the link through which demographic change is transmitted to the tax wedge on formal labor.


\subsection{Markets}
\label{subsec:markets}
\par Throughout, we work in \emph{stationary per-cohort-mass} units. Aggregates are integrals of individual quantities against the stationary population measure $\mu_j^{g,\theta}$, with $\sum_g \sum_\theta \sum_j \phi_j^{g}(\theta) = 1$. Population growth $n_p$ is already absorbed into the cohort-weight recursion~\eqref{eq:cohort_weights}, so the aggregates defined below do not grow with the actual population. A symbol such as $K$ in the Markets and Equilibrium sections refers to the stationary capital stock per unit of cohort mass, related to the absolute capital stock by $K^{\mathrm{abs}}_t = K \cdot N_t$, where $N_t$ is total population at $t$ growing at gross rate $1+n_p$. Stocks and per-period flows are distinguished throughout. \emph{Stocks} (units: goods, or units of agent mass): $K$, $A_{\mathrm{dom}}$, $B$, $N^W$, $N^R$. \emph{Flows} (units: goods per period, or efficiency units of labor per period): $Y$, $C$, $M$, $G$, $(n_p+\delta) K$, $L$, $\Lambda_{\mathrm{void}}$, $\overline{pen}$. Gross investment in capital is $(n_p+\delta)K$ per cohort mass per period: $\delta K$ replaces depreciation and $n_p K$ equips the next cohort, which enters at gross rate $1+n_p$. The goods-market identity~\eqref{eq:goods_market} is a flow identity, and every term in it has units of goods per period.

\par Asset markets are incomplete. Households save in a single non-contingent asset $a$ and cannot borrow, so
\begin{equation}
\label{eq:borrowing_constraint}
a' \;\geq\; 0.
\end{equation}
Both factor prices $(r, w)$ are determined endogenously by~\eqref{eq:foc_r}--\eqref{eq:foc_w}. Markets clear when the following conditions hold.

\paragraph{Labor market.}
\begin{equation}
\label{eq:labor_market}
L \;=\; \sum_g\, \sum_\theta\, \sum_{j=1}^{j_R - 1}\, \int \nu_j(h, \eta;\, g, \theta)\, \ell_j(s;\, g, \theta)\, d\mu_{j}^{g,\theta}(s).
\end{equation}
\paragraph{Capital market.}
Domestic household savings finance the entire capital stock and government debt:
\begin{equation}
\label{eq:capital_market}
K \;=\; A_{\mathrm{dom}} - B,
\qquad
A_{\mathrm{dom}} \;=\; \sum_g\,\sum_\theta\, \sum_{j=1}^{J}\, \int a\, d\mu_j^{g,\theta}(s).
\end{equation}
This condition, together with the firm's first-order condition~\eqref{eq:foc_r}, determines the equilibrium interest rate.

\paragraph{Goods market.}
\begin{equation}
\label{eq:goods_market}
Y \;=\; C + M + (n_p + \delta)\, K + G + \Lambda_{\mathrm{void}},
\end{equation}
where every term is a flow in stationary per-cohort-mass units, $(n_p + \delta) K$ is gross investment, and the voided-asset flow is
\begin{equation}
\label{eq:leakage_voided}
\Lambda_{\mathrm{void}}
\;\equiv\;
\sum_g\, \sum_\theta\, \sum_{j=1}^{J}\, \bigl(1 - \bar\psi_{j+1}^{g}(\theta)\bigr)
\int a_j'(s;\, g, \theta)\, d\mu_{j}^{g,\theta}(s).
\end{equation}
\paragraph{Government bond market.}
The government issues exactly $B$ bonds, and households' aggregate asset holdings finance both $K$ and $B$:
\begin{equation}
A_{\mathrm{dom}} = K + B,
\end{equation}
which is~\eqref{eq:capital_market} rearranged.

\subsection{Optimization Problems}
\label{subsec:optimization}

\subsubsection{Net Labor Income}
\label{subsubsec:net_income}

Net labor income at age $j$ is
\begin{equation}
\label{eq:net_labor_income}
y_j(h, \eta, \ell;\, g, \theta) \;=\; w\, \nu_j(h, \eta;\, g, \theta)\, \ell\, \bigl(1 - \tau^\omega - \tau^p\bigr),
\end{equation}
where $\tau^\omega$ is the labor income tax and $\tau^p$ the pension contribution rate. All workers in the baseline model are formal; a formalization counterfactual is left for future work.

\subsubsection{Workers' Problem}
\label{subsubsec:worker}

Individuals observe their productivity shock and then choose consumption, savings, labor supply, and health investment. Given labor supply $\ell$, household resources at age $j$ are
\begin{equation}
\label{eq:resources_worker}
x_j(a, h, \eta, \ell;\, g, \theta) \;=\; \bigl(1 + r(1 - \tau^k)\bigr) a + y_j(h, \eta, \ell;\, g, \theta).
\end{equation}
The worker's Bellman equation is
\begin{equation}
\label{eq:worker_bellman}
V_j(a, h, \eta;\, g,\theta)
\;=\;
\max_{\substack{c \geq 0,\, a' \geq 0 \\ \ell \in [0,1],\, m \geq 0}}
\Bigl\{
u(c, \ell, h)
+ \beta \sum_{k \in \{0,1\}} \Pr(k \mid j)\;
\psi_{j+1}^{g,\theta}\bigl(h'_k\bigr)\,
\E\bigl[V_{j+1}(a', h'_k, \eta';\, g, \theta) \mid \eta\bigr]
\Bigr\}
\end{equation}
subject to
\begin{align}
(1 + \tau^c)\, c + (1 + \tau^m)\, m + a' &= x_j(a, h, \eta, \ell;\, g, \theta), \label{eq:worker_budget}\\
h'_k &\ \text{given by~\eqref{eq:health_lom_shock} with } \tilde h = (1 - \delta_j^h)\, h + H_j(m, h), \label{eq:worker_health}\\
a' &\geq 0,\quad m \geq 0,\quad \ell \in [0,1]. \label{eq:worker_borrow}
\end{align}
The survival factor sits inside the sum over shock nodes, because a household struck by a shock faces both a lower continuation value and a lower probability of reaching it; collapsing the two would price survival at the wrong health level. The expectation over $\eta'$ is nested inside the expectation over health and the two are independent by construction, so health shocks carry no information about future earnings beyond their effect on the health state itself. Denote the optimal savings policy by $a_j'(a, h, \eta;\, g, \theta)$, and analogously for $c_j, \ell_j, m_j$.

\subsubsection{Retiree's Problem}
\label{subsubsec:retiree}
For retirees ($j \geq j_R$) labor income is zero, so resources are
\begin{equation}
\label{eq:resources_retiree}
x_j(a;\, g, \theta) \;=\; \bigl(1 + r(1-\tau^k)\bigr) a + \overline{pen}.
\end{equation}
The productivity shock $\eta$ does not enter retirees' resources but is retained as a state to keep the dimension of the state space uniform across ages. The Bellman equation is
\begin{equation}
\label{eq:retiree_bellman}
V_j(a, h, \eta;\, g, \theta)
\;=\;
\max_{c \geq 0,\, a' \geq 0,\, m \geq 0}
\Bigl\{
u(c, 0, h)
+ \beta \sum_{k \in \{0,1\}} \Pr(k \mid j)\;
\psi_{j+1}^{g,\theta}\bigl(h'_k\bigr)\,
\E\bigl[V_{j+1}(a', h'_k, \eta';\, g, \theta) \mid \eta\bigr]
\Bigr\}
\end{equation}
subject to
\begin{align}
(1 + \tau^c)\, c + (1 + \tau^m)\, m + a' &= x_j(a;\, g, \theta), \label{eq:retiree_budget}\\
h'_k &\ \text{given by~\eqref{eq:health_lom_shock} with } \tilde h = (1 - \delta_j^h)\, h + H_j(m, h), \label{eq:retiree_health}\\
a' &\geq 0,\quad m \geq 0. \label{eq:retiree_borrow}
\end{align}
At the terminal age $j = J$, $\psi_{J+1}^{g,\theta}(\cdot) \equiv 0$, the continuation value vanishes, and the optimum is $a'_J = 0$ and $m_J = 0$: neither assets nor health carry any continuation value once survival is impossible.

\paragraph{Optimality.}
The first-order, envelope, and Kuhn--Tucker conditions of the two problems are collected in Appendix~\ref{app:focs}; three of their properties organize the economics. First, medical spending earns two returns that the model keeps distinct: a survival return---spending raises next-period health, and with it the probability of reaching the continuation value at all---and a health-capital return---conditional on survival, the extra health that spending produces is valued at the marginal continuation value of health, which collects its productivity, amenity, and future-survival returns. The medical-spending condition~\eqref{eq:foc_m} prices both, averaged over the two nodes of the health lottery. Second, survival enters the consumption Euler equation~\eqref{eq:euler} as an effective discount factor, so households that invest in health also lengthen the horizon over which they save; and because next period's marginal utility is evaluated at the realized health node, the health shock is marginal-utility risk against which precautionary saving is the only available insurance---there is no insurance market for the shock. Third, the amenity cancels from the intratemporal consumption--leisure condition~\eqref{eq:cons_leisure}, so health raises labor supply only through productivity $\nu_j$ and the labor margin remains free of wealth effects.

\subsection{Stationary Competitive Equilibrium}
\label{subsec:equilibrium}

Let $\mu_j^{g,\theta}$ denote the time-invariant population measure of age-$j$ households of sex $g$ and skill $\theta$ over the state space $\mathcal{S} = \mathcal{A} \times \mathcal{H}_h \times \mathcal{H}_\eta$, with $\mathcal{A} = [0, \bar a]$, $\mathcal{H}_h = (0, \bar h]$, and $\mathcal{H}_\eta = \{\eta_1, \ldots, \eta_{N_\eta}\}$. The measure is \emph{unconditional}: $\int d\mu_j^{g,\theta} = \phi_j^{g}(\theta)$, so aggregates are integrals over it without an outer cohort weight. Let $\Pi^\eta(\eta, \eta')$ denote the transition kernel of the AR(1) process~\eqref{eq:ar1} discretized to the finite grid $\mathcal{H}_\eta$; the discretization is described in Section~\ref{sec:calibration}.

\subsubsection{Aggregate Objects}

\begin{align}
C &\;=\; \sum_g\,\sum_\theta\, \sum_{j=1}^{J}\, \int c_j(s;\, g, \theta)\, d\mu_j^{g,\theta}(s), \label{eq:agg_C}\\
M &\;=\; \sum_g\,\sum_\theta\, \sum_{j=1}^{J}\, \int m_j(s;\, g, \theta)\, d\mu_j^{g,\theta}(s), \label{eq:agg_M}\\
L &\;=\; \sum_g\,\sum_\theta\, \sum_{j=1}^{j_R - 1}\, \int \nu_j(h, \eta;\, g, \theta)\, \ell_j(s;\, g, \theta)\, d\mu_j^{g,\theta}(s), \label{eq:agg_L}\\
A_{\mathrm{dom}} &\;=\; \sum_g\,\sum_\theta\, \sum_{j=1}^{J}\, \int a\, d\mu_j^{g,\theta}(s), \label{eq:agg_A}\\
N^W &\;=\; \sum_g\,\sum_\theta\, \sum_{j=1}^{j_R - 1}\, \phi_{j}^{g}(\theta), \qquad N^R \;=\; \sum_g\,\sum_\theta\, \sum_{j=j_R}^{J}\, \phi_{j}^{g}(\theta), \label{eq:agg_NWR}\\
\Lambda_{\mathrm{void}} &\;=\; \sum_g\,\sum_\theta\, \sum_{j=1}^{J}\, \bigl(1 - \bar\psi_{j+1}^{g}(\theta)\bigr) \int a_j'(s;\, g, \theta)\, d\mu_j^{g,\theta}(s). \label{eq:agg_Lvoided}
\end{align}

\subsubsection{Law of Motion for the Distribution}

For each measurable set $\mathcal{B} \subseteq \mathcal{S}$, each age $j = 1, \ldots, J-1$, each sex $g$, and each skill $\theta$:
% REVIEW-COPY ONLY: display broken over two lines (\begin{split}); as written
% it overflows the page edge by ~161pt and the right end is clipped. No
% mathematical token was changed.
\begin{equation}
\label{eq:lom}
\begin{split}
\mu_{j+1}^{g,\theta}(\mathcal{B})
\;=\;
\frac{1}{1 + n_p}
\int
\sum_{k \in \{0,1\}} \Pr(k \mid j)\;
\psi_{j+1}^{g,\theta}\bigl(h'_k(s; j, g, \theta)\bigr)\,\\
\Bigl[\sum_{\eta' \in \mathcal{H}_\eta}
\Pi^\eta(\eta, \eta')\,
\mathbf{1}\bigl[\bigl(a_j'(s;\, g, \theta),\, h'_k(s; j, g, \theta),\, \eta'\bigr) \in \mathcal{B}\bigr]
\Bigr]
\, d\mu_j^{g,\theta}(s),
\end{split}
\end{equation}
where $h'_k(s; j, g, \theta)$, $k \in \{0,1\}$, are the two nodes of the health lottery~\eqref{eq:health_lom_shock} evaluated at the optimal policies. Mass flows to both nodes in proportion to the arrival probability, each node surviving at its own health level, so the health shock is a source of within-cohort dispersion in the stationary distribution and not only of individual risk.
The factor $1/(1+n_p)$ expresses the unconditional next-period measure in per-newborn-cohort units, consistent with the cohort-weight recursion~\eqref{eq:cohort_weights}. The recursion runs for $j = 1, \ldots, J-1$; no measure $\mu^{g,\theta}_{J+1}$ is defined, because $\psi^{g}_{J+1} \equiv 0$ implies that all age-$J$ agents die at the period boundary. The initial condition for newborns is
\begin{equation}
\label{eq:initial_dist}
\mu_1^{g,\theta}(\mathcal{B})
\;=\;
\phi_{1}^{g}(\theta) \cdot
\sum_{\eta_1 \in \mathcal{H}_\eta} \pi^\eta_{\mathrm{erg}}(\eta_1)\,
\mathbf{1}[(0, h_0, \eta_1) \in \mathcal{B}],
\end{equation}
where $\pi^\eta_{\mathrm{erg}}$ is the ergodic distribution of $\Pi^\eta$ and $h_0 \in (0, \bar h]$ is the initial health endowment, a parameter of the model. Newborns hold no assets and draw $\eta$ from the ergodic distribution, so all initial heterogeneity across households of a given type is in the productivity draw.

\subsubsection{Definition}

\begin{definition}[Stationary Recursive Competitive Equilibrium]
\label{def:equilibrium}
Given exogenous policy $(\tau^c, \tau^\omega, \tau^k, \tau^m, g_y, \kappa)$, birth shares $\{\pi^{g,\theta}_1\}_{g,\theta}$, population growth $n_p$, initial health endowment $h_0 \in (0, \bar h]$, and parameters of preferences, technology, mortality, productivity, and health formation, a \emph{stationary recursive competitive equilibrium} consists of
\begin{enumerate}
\item value functions $\{V_j(\cdot;\, g, \theta)\}_{j=1,\ldots,J;\, g,\theta}$ and policy functions $\{c_j, \ell_j, m_j, a_j'\}(\cdot;\, g, \theta)$ defined on $\mathcal{S}$;
\item factor prices $(r, w)$;
\item endogenous fiscal and pension variables $(B, \tau^p, \overline{pen})$;
\item aggregates $(K, L, C, M, A_{\mathrm{dom}}, N^W, N^R, \Lambda_{\mathrm{void}})$ in stationary per-cohort-mass units (Section~\ref{subsec:markets});
\item cohort weights $\{\phi_j^{g}(\theta)\}_{j, g, \theta}$ and population-average survivals $\{\bar\psi_{j+1}^{g}(\theta)\}_{j, g, \theta}$;
\item stationary distributions $\{\mu_j^{g,\theta}\}_{j, g, \theta}$;
\end{enumerate}
such that:
\begin{enumerate}
\item \emph{Worker optimality.} For each $(g, \theta)$ and $j < j_R$, $V_j$ solves the worker's Bellman problem~\eqref{eq:worker_bellman}--\eqref{eq:worker_borrow}, with policy functions $(c_j, \ell_j, m_j, a_j')$.
\item \emph{Retiree optimality.} For each $(g, \theta)$ and $j \geq j_R$, $V_j$ solves the retiree's Bellman problem~\eqref{eq:retiree_bellman}--\eqref{eq:retiree_borrow}, with policy functions $(c_j, m_j, a_j')$ and $\ell_j \equiv 0$.
\item \emph{Firm optimality.} $(r, w)$ satisfy~\eqref{eq:foc_r}--\eqref{eq:foc_w} given $(K, L)$.
\item \emph{Government budget.} $B$ satisfies~\eqref{eq:gov_stationary}.
\item \emph{Pension balance.} $\tau^p$ satisfies~\eqref{eq:pension_contribution} and $\overline{pen}$ satisfies~\eqref{eq:pension_benefit}.
\item \emph{Market clearing.} Labor~\eqref{eq:labor_market}, capital~\eqref{eq:capital_market}, and goods~\eqref{eq:goods_market}.
\item \emph{Cohort weights.} $\phi_{j}^{g}(\theta)$ satisfy~\eqref{eq:cohort_weights} with $\bar\psi_{j+1}^{g}(\theta)$ given by~\eqref{eq:avg_survival} and birth shares $\phi_1^{g}(\theta) \propto \pi^{g,\theta}_1$, normalized so that $\sum_g \sum_\theta \sum_j \phi_j^{g}(\theta) = 1$.
\item \emph{Stationarity of the distribution.} For each $(g, \theta, j)$, $\mu_j^{g,\theta}$ satisfies the law of motion~\eqref{eq:lom} with the newborn condition~\eqref{eq:initial_dist}.
\item \emph{Aggregate consistency.} $(C, M, L, A_{\mathrm{dom}}, N^W, N^R, \Lambda_{\mathrm{void}})$ are computed from~\eqref{eq:agg_C}--\eqref{eq:agg_Lvoided}, all in stationary per-cohort-mass units.
\item \emph{Demographic--distributional fixed point.} Conditions 7 and 8 are mutually dependent. The population-average survivals $\bar\psi^{g}_{j+1}(\theta)$ are functionals of the stationary distribution $\mu^{g,\theta}_j$ through~\eqref{eq:avg_survival}, while $\mu^{g,\theta}_j$ depends on $\bar\psi^{g}_{j+1}(\theta)$ through the $1/(1+n_p)$ scaling in~\eqref{eq:lom} that anchors the cohort weights. Equilibrium requires $\bigl(\{\phi^{g}_j(\theta)\}, \{\bar\psi^{g}_{j+1}(\theta)\}, \{\mu^{g,\theta}_j\}\bigr)$ to be a joint fixed point of~\eqref{eq:cohort_weights},~\eqref{eq:avg_survival}, and~\eqref{eq:lom}, evaluated at the equilibrium policy functions.
\end{enumerate}
\end{definition}

Condition 10 is what distinguishes this equilibrium from one with exogenous mortality. Because survival depends on a state that households choose, the age distribution of the population is an equilibrium object rather than a demographic input, and the same medical-spending decisions that determine health also determine the dependency ratio that prices formal labor.

\subsection{Existence and Computation}
\label{subsec:existence}

The equilibrium of Definition~\ref{def:equilibrium} is the fixed point of three nested maps. The inner map is the household problem, solved by backward induction on the finite-horizon Bellman equations~\eqref{eq:worker_bellman} and~\eqref{eq:retiree_bellman} together with the savings and medical-spending conditions~\eqref{eq:foc_a}--\eqref{eq:foc_m} of Appendix~\ref{app:focs}, with labor supply concentrated out through the intratemporal condition~\eqref{eq:cons_leisure} and constrained maximization over the remaining controls $(a', m)$. The middle map iterates the population measure~\eqref{eq:lom} jointly with the cohort weights~\eqref{eq:cohort_weights} and the health-dependent survival schedule~\eqref{eq:avg_survival} to their mutual fixed point, which is condition 10 of the definition. The outer map adjusts the capital--labor ratio until the factor-price conditions~\eqref{eq:foc_r}--\eqref{eq:foc_w} and the capital-market condition~\eqref{eq:capital_market} clear, with the contribution rate and the stationary debt level moving through~\eqref{eq:pension_contribution} and~\eqref{eq:gov_stationary}.

The endogenous grid method is not available here. Medical spending enters the continuation value non-separably and through its level, since the survival derivative in~\eqref{eq:foc_m} multiplies $\E[V_{j+1}]$ rather than only its marginal value, and next-period health is itself a control. The end-of-period state is therefore not exogenous and the Euler equation cannot be inverted in closed form.

We do not prove existence. The standard arguments for stationary equilibria in this class take the survival schedule as a primitive; here it is a functional of the equilibrium distribution through~\eqref{eq:avg_survival}, and the household objective contains the product $\psi^g_{j+1}(h')\,\E[V_{j+1}]$, whose concavity in $(a',m)$ is not implied by the restrictions placed on preferences and on $H_j$. Existence is therefore established constructively: the algorithm converges to a point satisfying every condition of Definition~\ref{def:equilibrium}, and the convergence record, the Euler-equation residuals, and the market-clearing residuals that support that claim are reported in Appendix~\ref{app:calibration}. Solved life-cycle profiles are displayed at the calibrated baseline in Section~\ref{subsec:calibrated_baseline}. One closure choice is worth restating at this point: assets of decedents are voided rather than bequeathed or annuitized, which isolates the life-cycle saving motive at the cost of a resource leakage, $\Lambda_{\mathrm{void}}$ in~\eqref{eq:goods_market}. An accidental-bequest or annuity closure is the natural alternative and is left to future work.

Two features of the closure organize everything that follows. First, the PAYG block ties the contribution rate to the system dependency ratio, $\tau^p = \kappa\,N^R/N^W$, so any force that moves the ratio of retirees to workers moves the labor-tax wedge one for one. Second, because survival is endogenous to medical spending, health investment feeds forward into the dependency ratio itself: longer lives raise $N^R/N^W$, and hence $\tau^p$. The interaction of these two links --- demographic pressure transmitted to the fiscal wedge, and health-driven longevity feeding back into that same wedge --- is the \emph{fiscal amplification loop}. It is the mechanism the experiments of Section~\ref{sec:experiments} stress from the dependency side, and the reason the net fiscal effect of better population health is ambiguous in sign.
%% ====================================================================
\section{Calibration}
\label{sec:calibration}

This section is the empirical core of the paper. We calibrate the model to M\'exico in two steps. The first step fixes the exogenous processes and the externally assigned parameters: demographics, the sex- and education-specific baseline survival schedules and the survival floor, the sex- and skill-specific age-efficiency profiles and skill intercepts estimated from the ENIGH household surveys, the productivity process, the health-depreciation schedule, and the fiscal block. The second step estimates the remaining nine parameters---preferences, the health-production technology, the health-risk process, and the health--productivity slope---by classical minimum distance on ten moments, nine of them estimated from Mexican micro-data. Section~\ref{subsec:firststep} documents the inputs, Section~\ref{subsec:secondstep_smm} the estimator and targets, Section~\ref{subsec:calibration_results} the estimates and fit, and Section~\ref{subsec:validation} the untargeted diagnostics. Identification limits are stated where they bind rather than collected in a disclaimer, and three of them govern how the results may be read: the health--productivity slope is identified only jointly with the model's internal health scale, the two health-risk parameters share their identifying moments with the rest of the health block, and the two targeted survival levels overlap a measured first-step input, so they function as disclosed consistency anchors rather than as independent identifying information.

\begin{quote}
\small\noindent\textbf{Status of the numbers in this paper.}
\label{par:superseded}
The estimation of Section~\ref{subsec:secondstep_smm} has been run, and Sections~\ref{subsec:calibration_results}--\ref{subsec:validation} report its results. Those numbers are current, were produced under the model of Section~\ref{sec:model}, and are evaluated at the prices of the model's own general-equilibrium solve at the estimates; that solve is the paper's calibrated baseline. Its aggregates are these. Investment is $22.6\%$ of output against $20.2\%$ in the data, an overshoot of $2.4$ percentage points that moves with the higher population growth rate and the longer lives of this calibration and is disclosed rather than chased. Pension spending is $3.4\%$ of GDP against a measured $4.5\%$---the right order of magnitude, with a shortfall to be expected, since the measured figure includes pensions paid to the legacy defined-benefit (Ley~73) transition generation, which a stationary model cannot carry. Public debt stands at $1.12$ of annual output against roughly $0.5$ for M\'exico---a disclosed limitation of the residual closure. The capital--output ratio is $2.72$ against $3.29$ in the data, $17\%$ below: the shortfall was $22\%$ under the preceding baseline and narrowed by about a fifth when the education-specific survival schedules entered the first step---longer high-education lives strengthen the life-cycle saving motive---which is a result of the longevity-savings channel, not a recalibration. The remainder survived a full re-estimation of the household parameters at the sourced first-step bundle, so it is a property of the calibrated configuration rather than of any particular parameter point, and the two margins that might yet move it (a discount-factor re-test at this bundle, the replacement rate within its sensitivity band) are untried by choice. We report it and do not chase it.

One archived solve, reported for provenance in Appendix~\ref{app:provenance}, was produced under an \emph{earlier specification} of both the model and its first step. This paragraph is the canonical list of how that superseded specification differs from the one this paper describes; the marked passages elsewhere refer back to it.

\emph{Preferences.} The health amenity entered the GHH composite additively, $u=[(c+s(h)-v(\ell))^{1-\gamma}-1]/(1-\gamma)$, and there was no flow value of being alive $\chi$. The change is one of model primitive, not of calibration: Section~\ref{subsec:preferences} shows the two forms differ in whether willingness to pay for health carries an income gradient at all.

\emph{Health block.} Health was deterministic, with the constant-returns law of motion that Section~\ref{subsubsec:health} now states as the nested boundary case; there was no health shock, no deficit-scaled repair, and no education dimension in survival.

\emph{Estimation block.} Six free parameters (five, with $\xi$ frozen), with the health--productivity slopes held at external anchors rather than estimated, against six literature-anchored moments---not the nine parameters and ten micro-data moments of Section~\ref{subsec:secondstep_smm}.

\emph{First step.} The labour inputs came from ENOE rather than ENIGH; the persistent-income process was $(\rho,\sigma_\varepsilon)=(0.782,0.265)$; the health-depreciation schedule ran from $0.02$ to $0.80$ per five-year period rather than the estimated $0.0011$ to $0.186$---roughly an eighteen-fold difference at working ages---and the survival floor $\varpi$ did not exist.

\emph{Aggregate calibration.} The fiscal and technology block was unsourced: statutory tax rates $(0.16,\,0.20,\,0.20)$, a public-goods share of $0.19$, a replacement rate of $0.50$, a capital share of $0.36$, capital depreciation of $8.23\%$ per year, and population growth of $1\%$ per year---against the sourced, concept-matched bundle of Table~\ref{tab:firststep_params}.

\noindent The reproducibility caveats attaching to that solve are in Section~\ref{subsec:repro}.
\end{quote}

\subsection{Units and Normalization}
\label{subsec:units}

The model is dimensionless until two normalizations fix its scale, and every quantitative statement in the paper depends on stating them once and applying them consistently. \emph{Time.} A model period is five years; $j=1$ is age 20 and retirement is at $j_R=10$ (age 65). Output, consumption, medical spending, and the wage bill are five-year flows; capital and government debt are stocks. Ratios of a stock to a flow are therefore reported at annual frequency by multiplying the model (five-year) ratio by five: the annual capital--output ratio is $5\,(K/Y)$ and the annual debt--output ratio is $5\,(B/Y)$. Rates are annualized geometrically, $r^{\mathrm{ann}}=(1+r)^{1/5}-1$ for the five-year return $r$. \emph{Goods.} One unit of the consumption good is converted to dollars by a single factor, $\text{usd\_per\_unit\_c}=15{,}000$. This factor is a \emph{reporting convention only} and is uncalibrated: it enters no moment---the value of a statistical life is reported as a scale-free model output rather than targeted (Section~\ref{subsec:vsl_structural})---and it is retained solely so that model quantities can be quoted in dollars. Dollar figures in this paper should be read accordingly, and none of the estimation depends on it. \emph{Labour.} Effective labour is measured in efficiency units normalized so that the age-efficiency of a male aged 20--24 equals one; the wage $w$ is the price of one such efficiency unit. \emph{Population.} All aggregates are in stationary per-cohort-mass units (Section~\ref{subsec:markets}); levels therefore have no independent meaning across calibrations and only ratios, rates, and within-calibration comparisons are interpreted.

\subsection{First-Step Inputs}
\label{sec:data}
\label{subsec:firststep}

Table~\ref{tab:firststep_params} collects the externally assigned and first-step parameters with their values, units, and sources. Every aggregate parameter is sourced for M\'exico at a stated concept---effective tax rates rather than statutory, government consumption net of public health, a dependency-matched population growth rate---and the remainder of this subsection records how the Mexican data discipline each block.

\begin{table}[H]
\centering
\caption{Externally assigned and first-step parameters (M\'exico calibration).}
\label{tab:firststep_params}
\small
\setlength{\tabcolsep}{4pt}
% REVIEW-COPY ONLY: table as written overflows the text width under these
% margins/font; scaled to fit so no column is clipped.
\resizebox{\textwidth}{!}{%
\begin{tabular}{@{}llcl@{}}
\toprule
 & \theadl{Description} & \thead{Value} & \theadl{Target/Source} \\
\midrule
\tgroup{4}{Demographics and technology}
                              & Period length & 5 years & model frequency \\
$J,\ j_R$                     & Periods of life; retirement period & $17,\ 10$ & max age 100; retire at 65 \\
$n_p$                         & Population growth & $3.51\%$/yr ($18.8\%$ per 5-yr) & dependency-matched: $65{+}/(20\text{--}64)=0.1414$, INEGI Censo 2020 (see text) \\
$\alpha$                      & Capital share & $0.40$ & M\'exico factor shares (see text) \\
$\delta$                      & Capital depreciation & $5\%$/yr ($22.6\%$ per 5-yr) & M\'exico-specific literature value (see text) \\
$A$                           & TFP (scale) & $1.60$ & normalization, conditional (see text) \\
\tgroup{4}{Preferences (fixed externally)}
$\gamma$                      & GHH curvature & $2.0$ & --- \\
$\nu_\ell$                    & Inverse Frisch & $2.0$ & --- \\
$\beta$                       & Discount factor, per 5-yr & $0.998^5$ & external; $K/Y$-invariant (see text) \\
$\xi$                         & Amenity curvature & $0.50$ & pinned (see Section~\ref{subsec:secondstep_smm}) \\
\tgroup{4}{Fiscal and pension}
$\tau^c,\ \tau^\omega,\ \tau^k$ & Consumption, labour, capital taxes & $0.09,\ 0.10,\ 0.12$ & effective, concept-matched rates (see text) \\
$\tau^m,\ g_y$                & Medical tax (baseline); $G/Y$ & $0,\ 0.090$ & $g_y$ ex public health (see text) \\
$\kappa$                      & Pension replacement rate & $0.40$ & flat pension $\div$ average earnings (see text) \\
\tgroup{4}{Labour (ENIGH 2016, 2018; this paper)}
$\theta_L,\ \theta_H$         & Skill intercepts, log points & $-0.020,\ +0.020$ & ages 20--24, GHH-adjusted (see text) \\
% RESOLVED (code path verified 2026-08-11): under the skill-specific
% profile staging the solver consumes the skill-specific intercepts
% (-0.0202/+0.0201, skill_params_skill.csv) via the apply_first_step!
% overwrite; the common-profile pair (-0.240/+0.239) remains the
% deterministic/parity input only. The intercepts are small BECAUSE the
% profiles carry the skill divergence beyond the normalization age --
% each skill's profile is normalized at its own (male, 20-24) cell.
$e^{g,\theta}_j$              & Age-efficiency profiles, efficiency units & profiles & sex- and skill-specific; male 20--24 $\equiv 1$ \\
$\pi^{g,\theta}_1$            & Birth shares & shares & sex-by-education composition at age 20 \\
\tgroup{4}{Productivity and health}
$\rho,\ \sigma_\varepsilon$   & AR(1) of $\eta$, per 5-yr & $0.837,\ 0.284$ & external (see text) \\
$\delta^h_j$                  & Health depreciation, per 5-yr & $0.0011\!\to\!0.186$ & age-rising; ENASEM (see text) \\
$\psi^{\circ\, g,\theta}_j$   & Baseline survival, sex- and education-specific & schedules & ENASEM + CONAPO (see text) \\
$\varpi$                      & Survival floor as $h\downarrow 0$ & $0.623$ & ENASEM \\
$h_0$                         & Initial health, health units & $\bar h = 1$ & --- \\
\tgroup{4}{Measurement}
$s$                           & Out-of-pocket share of health spending & $0.41$ & WHO GHED; see Section~\ref{subsec:secondstep_smm} \\
\bottomrule
\end{tabular}}
\tabnote{0.92}{Rates with ``/yr'' are annual; per-5-yr values are the model-period magnitudes used in computation. Tax rates are effective, concept-matched rates, not statutory rates (see text). $\tau^p$ is not listed because it is endogenous (it clears the PAYG block). The internally estimated parameters---$\Psi$, $\Xi$, $\overline{H}_0$, $h^{\mathrm{slope}}$, $\zeta_h$, $\varrho$, $\chi$, $\bar\pi^h$, $\varsigma$---are listed in Table~\ref{tab:theta_hat_md}. The $\delta^h_j$ range spans the youngest working band to the terminal bands (flat from age 90); its smallness at working ages is the fact behind the health-block design discussed in Section~\ref{subsec:health_block_evidence}.}
\end{table}

\paragraph{Technology, demographic scale, and the discount factor.}
Capital depreciates at $\delta=0.05$ per year, the M\'exico-specific value of \citet{BergoeingKehoeKehoeSoto2002}, who validate it against the UN-reported share of capital consumption in Mexican GDP and explicitly consider and reject $0.08$; the capital share $\alpha=0.40$ follows the factor-share measurement for M\'exico of \citet{GarciaVerdu2005}. The TFP level $A=1.60$ is a normalization, with one caveat stated precisely: the model is non-homothetic (bounded health, the flow value $\chi$), so output scale is not neutral, and the normalization is internally consistent only because the estimation of Section~\ref{subsec:secondstep_smm} is carried out at this value---changing it would require re-estimation. Population growth is set so that the model's stationary age distribution reproduces the 2020 census: $n_p=3.51\%$ per year delivers an old-age dependency ratio $65{+}/(20\text{--}64)$ of $0.1414$ (INEGI Censo 2020) under the model's survival schedules---a growth rate at or above the upper end of the range M\'exico actually experienced over the decades that produced the 2020 age pyramid, the price of matching a 2020 age distribution with a stationary demography. Because the required $n_p$ depends on the survival schedules, it is re-solved by bisection whenever those schedules change---the present value reflects the education-specific schedules below---so the dependency ratio is matched by construction at the reported baseline (realized $0.1414$, within $0.02\%$ of the census target) rather than approximately, and it is accordingly not listed among the validation checks. The discount factor is fixed externally at $\beta=0.998^5$. A natural objection is that $\beta$ should instead be disciplined by the capital--output ratio; the recorded diagnostic is that it cannot be: pushing $\beta$ to the top of its admissible range moved the equilibrium capital--output ratio by less than two per cent, because the general-equilibrium interest-rate response stabilizes capital and because the effective discount is $\beta$ times survival, which dominates at the ages where wealth is held. That test predates the final parameter bundle and has deliberately not been repeated at it; either way, $\beta$ is not the instrument for the capital--output shortfall disclosed at the head of this section.

\paragraph{Demographics and mortality.}
The baseline survival schedules $\psi^{\circ\, g,\theta}_j$---four of them, one per sex-by-education group---are the paper's own estimates from the ENASEM/MHAS panel: deaths counted between waves identify survival over ages 50 and above, using the same decoupled estimator as the earlier pooled-by-education schedules with a parametric education index (an education main effect and an education-by-age interaction), the pre-committed strategy for the thin high-education cells at the oldest ages. Below age 50 the schedules are spliced to official Mexican (CONAPO) life tables, pooled across education---no education split of mortality exists at those ages---so the education gradient of baseline survival enters at ages 50 and above only. The survival floor $\varpi=0.623$---the fraction of baseline survival retained as health approaches its lower bound in the modulation~\eqref{eq:survival_modulation}---is estimated once, in the pooled pass, from the mortality gradient in the measured deficit index among the 50+ population, and is read rather than re-estimated when the education split is constructed. One limitation of this input must be disclosed at its source, because the fit results inherit it: the schedules are capped at one, and measured survival at ages 50--69 is high enough that the cap clips the estimated education gradient there---entirely for women, partially for men---so the input understates the education difference in survival at exactly the ages the targeted survival moments cover. This is the known global-floor limitation of the survival estimator; the fix, an age-varying floor, is deferred (Section~\ref{subsec:roadmap}).\footnote{A mechanical consequence of the same cap: the female baseline sits essentially at its ceiling at ages 50--64, so the female half of the survival anchors is identified off the health modulation alone.} The female survival advantage embedded in these schedules is the differential the gender validation of Section~\ref{subsec:validation} exploits.

\paragraph{Labour productivity.}
The labour block is estimated from the pooled ENIGH 2016 and 2018 cross-sections---the waves contemporaneous with the ENASEM window that defines the steady state being calibrated; later waves are held out as an out-of-sample check on the measurement layer described in Section~\ref{subsec:secondstep_smm}. Taking logs of the productivity equation~\eqref{eq:labor_productivity} gives the decomposition
\begin{equation}
\label{eq:logwage_decomp}
\log w_{i,t} \;=\; \underbrace{\log e^{g,\theta}_{j(i,t)}}_{\text{age-by-sex-by-skill profile}} \;+\; \underbrace{\theta_{i}}_{\text{skill intercept}} \;+\; \underbrace{\varrho\, h_{i,t}}_{\text{health--productivity slope}} \;+\; \eta_{i,t}.
\end{equation}
From the survey we recover the age-efficiency profiles $e^{g,\theta}_j$ and the intercepts $\theta_L,\theta_H$, normalized so the male 20--24 cell equals one. Two features of this block matter for everything downstream. First, the profiles are \emph{skill-specific}: education groups differ in their measured returns to experience, and the difference is stark---the low-skill male profile peaks at ages 45--49 and falls back to roughly its age-20 level by 60--64, while the high-skill male profile roughly doubles by its peak in the fifties. The life-cycle growth of the skill premium is therefore a measured input, not an estimated outcome, a fact with a direct consequence for what the validation exercises can claim (Section~\ref{subsec:validation}). Second, the profiles and intercepts are estimated from log \emph{earnings} but consumed by the model as \emph{productivity}, and GHH labour supply re-amplifies any productivity shifter into earnings by the factor $1+1/\nu_\ell = 1.5$; the estimated log-earnings objects are accordingly divided by this factor before entering the model, so that the model reproduces the measured earnings gaps rather than amplifying them a second time. The division is the model-consistent correction; the empirical alternative---an hourly-wage profile built from reported hours, under which no amplification is implied at all---is preferable in principle and deferred.

The persistent component $\eta$ would ordinarily be recovered from the residuals of this regression, but ENIGH is a repeated cross-section and is therefore uninformative about earnings \emph{dynamics}, and no accessible Mexican household panel at the model's frequency exists. We consequently calibrate $(\rho,\sigma_\varepsilon)$ externally rather than estimating them. We take the annual persistence and innovation variance to be the average of representative estimates from the canonical restricted-income-profile literature \citep{StoreslettenTelmerYaron2004, Guvenen2009, HubbardSkinnerZeldes1994}, $\rho_{\mathrm{ann}}\approx0.965$ and $\sigma^2_{\varepsilon,\mathrm{ann}}\approx0.0185$, and time-aggregate the underlying annual AR(1) to the five-year period by skip-sampling, $\rho=\rho_{\mathrm{ann}}^5$ and $\sigma_\varepsilon^2=\sigma^2_{\varepsilon,\mathrm{ann}}\sum_{i=0}^{4}\rho_{\mathrm{ann}}^{2i}$, which gives $(\rho,\sigma_\varepsilon)=(0.837,0.284)$. The innovation scale is then re-pinned so that the model's stationary cross-sectional dispersion of log earnings matches that observed for Mexican formal-sector workers \citep{Puggioni2022}. The process is an external anchor and is treated as such: the persistence of the shock process is the dominant driver of the calibrated capital stock, so it is the input to which the aggregate results are most sensitive.

The health--productivity slope $\varrho$ in~\eqref{eq:logwage_decomp} is estimated internally within the minimum-distance block of Section~\ref{subsec:secondstep_smm} rather than fixed at an external anchor, and two features of the environment force this choice. First, working-age health cannot be measured directly: the Latin American health-and-aging surveys---Mexico's MHAS/ENASEM, Costa Rica's CRELES, and the Panama SABE study---all begin at age 50, so no health measure exists over most of the working life at which the wage slope must be identified. Hosseini, Kopecky, and Zhao (2025) exploit precisely such a panel for the United States, using PSID health deficits observed from age 25; the region has no working-age health panel of that kind. Second, even where health is observed, an external OLS of the log wage on the health stock is inconsistent with the model's own data-generating process: health is chosen through income-elastic medical investment, so the health stock is positively correlated with the persistent productivity shock $\eta$, biasing the OLS slope. Estimating $\varrho$ inside the structural model addresses both concerns, at the price of an interpretation constraint disclosed in Section~\ref{subsec:secondstep_smm}: the data identify $\varrho$ only jointly with the model's internal health scale.

\paragraph{Health depreciation.}
The health stock is a bounded quantity, $h\in(0,\bar h]$, with newborns endowed at the ceiling, $h_0=\bar h$, and $\bar h=1$ the normalization of Section~\ref{subsec:preferences}. Over the ages at which health is observed (50+), the age-rising depreciation rate $\delta^h_j$ is disciplined by a harmonized deficit index built from a \emph{reversible} subset of the ENASEM items---functional-status, depressive-symptom and self-rated measures that medical spending can plausibly restore---rather than the full frailty index of the deficit-accumulation literature, whose components are largely irreversible and therefore inconsistent with the restorable health capital of~\eqref{eq:health_production_deficit}. Health is measured on the model scale as $h = 1 - \text{deficit} \in (0,1]$.

The estimand is the \emph{within-survivor matched} contraction of the health stock: $1-\delta^h_j$ is computed on individuals observed in both adjacent waves, so the schedule is not contaminated by the composition change that mortality induces between waves. This decouples the depreciation estimate from the survival estimate, at the cost of measuring depreciation conditional on survival rather than unconditionally; the resulting profile should be read as a lower bound on the force of health decline, since the individuals who drop out are those whose health fell fastest. The wave spacing is three years and is rescaled to the model's five-year period. Working-age $\delta^h_j$, unobserved for want of a working-age panel, is set by log-linear extrapolation of the observed 50+ profile downward in age and capped at the highest reliably estimated band. The resulting schedule rises from $0.0011$ at ages 20--24 to $0.186$ from ages 90 on. \emph{The smallness of the working-age values is the binding fact for the paper's mechanism}: at a five-year depreciation of a tenth of one per cent, a constant-returns health technology returns essentially every spender to $\bar h$, and it is this degeneracy that the deficit-scaled, risky health block of Section~\ref{subsubsec:health} exists to break (the evidence is in Section~\ref{subsec:health_block_evidence}). The estimator is detailed in Appendix~\ref{app:calibration}.

\paragraph{Fiscal block.}
Every tax parameter enters at an \emph{effective}, concept-matched rate; statutory rates appear in this calibration only where statute and model concept coincide (the retirement age). The consumption tax $\tau^c=0.09$ is the effective rate implied by Mexican VAT and excise collections relative to private consumption: the OECD's VAT revenue ratio for M\'exico is about $0.35$, so the statutory $16\%$ IVA overstates the burden households actually face by roughly half. The labour tax $\tau^\omega=0.10$ is the income-tax-only wedge on the average worker (OECD \emph{Taxing Wages} puts it near $9.4\%$ in 2020, and effective-tax-rate estimates net of social contributions at $7.7$--$8.5\%$ \citep{BachasFisherPostJensenZucman2022}); social-security contributions are deliberately excluded from $\tau^\omega$ because the model levies the pension contribution $\tau^p$ separately and endogenously---it clears the PAYG block through~\eqref{eq:pension_contribution}, and at the baseline equals $5.66\%$, so the combined labour wedge is about $15.7\%$ against a measured total wedge near $20\%$, the difference being employer contributions (including the housing fund) outside the model. One disclosure attaches: the model's $\tau^\omega$ is flat, so distributional effects that run through the progressivity of the Mexican income tax are outside the model. The capital tax $\tau^k=0.12$ lies within the range of effective capital-tax estimates for M\'exico \citep[$0.108$--$0.121$;][]{BachasFisherPostJensenZucman2022}, consistent with the marginal-rate estimates of \citet{Anton2005}. The public-goods share $g_y=0.090$ is government consumption \emph{net of public health}: measured general-government consumption ($12.2\%$ of GDP in 2020) minus domestic government health spending ($3.2\%$, WHO Global Health Expenditure Database). The subtraction is the fiscal sibling of the out-of-pocket measurement wedge $s$ described in Section~\ref{subsec:secondstep_smm}: the same public financing that pays the $(1-s)$ share of medical care is removed from $G$, so it is counted exactly once. The replacement rate $\kappa=0.40$ matches the concept the model's flat pension implements---the non-contributory-style flat benefit divided by economy-wide average earnings---at the value documented for M\'exico by \citet{Bonthuis2024} and the OECD's \emph{Pensions at a Glance}. The stationary debt level $B$ is the residual that closes the general budget through~\eqref{eq:gov_stationary}; both $\tau^p$ and $B$ are equilibrium objects. The comparative fiscal and health-financing landscape for the three countries---of which the salient fact for the Mexican calibration is the high out-of-pocket share ($\approx 41\%$ of current health expenditure, against $\approx 22\%$ in Costa Rica), which both leaves households exposed to medical-expense risk and supplies the measurement wedge $s$---is collected in Appendix~\ref{app:supplementary}.

\subsection{Second-Step Estimation: Classical Minimum Distance}
\label{subsec:secondstep_smm}

The parameters not pinned by external evidence or first-step estimation are estimated jointly by classical minimum distance. The estimated vector is
\[
\Theta_{\mathrm{MD}} \;=\; \bigl(\Psi,\; \Xi,\; \overline{H}_0,\; h^{\mathrm{slope}},\; \zeta_h,\; \varrho,\; \chi,\; \bar\pi^h,\; \varsigma\bigr),
\]
nine parameters: the labour-disutility scale $\Psi$; the amenity scale $\Xi$ in $\Omega(h)$; the scale, age slope, and curvature of the health-production technology,
\begin{equation}
\label{eq:healthproduction}
\overline{H}_{j} \;=\; \overline{H}_{0}\,\exp\bigl(-h^{\mathrm{slope}}\, j\bigr), \qquad \zeta_h \in (0,1),
\end{equation}
in the deficit-scaled production function~\eqref{eq:health_production_deficit}---a positive $h^{\mathrm{slope}}$ tilts the technology downward with age, and the sign is estimated, not imposed; the health--productivity slope $\varrho$ of~\eqref{eq:labor_productivity}, internalized for the reasons given in Section~\ref{subsec:firststep}; the flow value of being alive $\chi$ of Section~\ref{subsec:preferences}; and the mean arrival rate $\bar\pi^h$ and size $\varsigma$ of the health shock~\eqref{eq:health_lom_shock}--\eqref{eq:shock_prob}. An earlier specification carried a tenth parameter, an additive education shifter inside the survival modulation, intended to absorb the non-investment education difference in mortality. It was retired for two reasons established before this baseline was adopted: its estimate proved to be a composition artifact---the objective rewarded ever-deeper negative values for removing high-education mass from the spending moments, at the direct expense of the survival moment the parameter existed for---and a sweep of its full admissible range showed the modulation's cap makes a non-negative shifter structurally incapable of producing the positive education--survival gap it was created to represent. The education mortality channel is instead measured: it now lives in the education-specific baseline schedules of the first step, and the shifter is gone from the model. Two parameters adjacent to this block are not estimated. The amenity curvature $\xi$ is pinned at an interior value: in the deterministic variant of the model, where this was measured, the objective is numerically flat in $\xi$ over its entire admissible range---the variation is at the level of solver noise---and $\xi$ has been held at the same interior value throughout the stochastic baseline; we disclose the pinned value as a convention rather than an estimate. The discount factor $\beta$ is fixed externally at $0.998^5$, since the capital--output ratio that would discipline it is an equilibrium object and estimation is carried out in partial equilibrium (below).

\paragraph{Why the estimator is minimum distance and not SMM.} Because the state space is discretized and the stationary distribution over $(j,g,\theta,a,h,\eta)$ is computed exactly by forward iteration, every model moment is an exact distribution-weighted functional rather than a Monte Carlo average over simulated panels. The objective is therefore deterministic---re-evaluating it at the same $\Theta$ returns exactly the same value---the asymptotic theory is that of classical minimum distance with no simulation-inflation term, and any non-smoothness the optimizer encounters is a property of the model, not of sampling noise. What remains is discretization: moments are kinked in $\Theta$ wherever mass crosses grid cells, which is one reason the search is derivative-free and among the reasons no standard errors are reported (below).

\paragraph{The objective.} The estimator minimizes the equally weighted sum of squared \emph{proportional} deviations,
\begin{equation}
\label{eq:smm_objective}
\widehat\Theta_{\mathrm{MD}} \;=\; \arg\min_{\Theta}\; \sum_{k=1}^{10} \left( \frac{m_k^{\mathrm{mod}}(\Theta) - m_k^{\mathrm{dat}}}{m_k^{\mathrm{dat}}} \right)^{2},
\end{equation}
that is, generalized minimum distance with $W=\operatorname{diag}\bigl(1/(m^{\mathrm{dat}})^2\bigr)$. We do not use the inverse-variance weighting $W=\operatorname{diag}(1/\mathrm{SE}^2)$, for two reasons stated plainly. First, the standard errors are not commensurable across the target set: one target is a convention with no sampling error at all (the hours normalization below), and the survey moments mix cluster-robust, delta-method, and bootstrap objects. Second, inverse-variance weighting turns the estimator into a matcher of $t$-statistics rather than of economic magnitudes: evaluated at the initial parameter vector it concentrated between $65$ and $80$ per cent of the objective on a single precisely estimated survey moment, leaving the moments that carry the paper's mechanism nearly weightless. Equal proportional weighting prices a ten-per-cent miss identically wherever it occurs. The cost of this choice is disclosed with equal plainness: it carries no efficiency claim; the system is over-identified by one, but the one over-identifying restriction is a consistency anchor rather than an economic restriction (below), so no overidentification test is reported; and one substantive property of the weighting must be on the record---proportional weights on two survival \emph{levels} near one make the sign of their small difference nearly costless, the fact that let an earlier specification's education survival shifter drift to a wrong-signed corner and that no reweighting can repair (Section~\ref{subsec:calibration_results}).

\paragraph{Targets.} Table~\ref{tab:targeted_moments} collects the ten targets. Nine are the paper's own estimates from Mexican micro-data; one---the hours target---is a conventional calibration value and is flagged as such. Standard errors are reported to size the misses in Section~\ref{subsec:calibration_results}; they do not enter the objective.

\begin{table}[H]
\centering
\caption{Targeted moments, values, and sources.}
\label{tab:targeted_moments}
\small
\setlength{\tabcolsep}{4pt}
\begin{tabular}{@{}lrrl@{}}
\toprule
\theadl{Targeted Moment} & \theadr{Target} & \theadr{SE} & \theadl{Source} \\
\midrule
\tgroup{4}{Labour}
Prime-age male hours, share of time endowment      & $0.333$  & ---      & convention (see text) \\
\tgroup{4}{Medical spending}
Income elasticity of medical spending, pooled       & $0.575$  & $0.029$  & ENIGH 2016--2018 \\
Income elasticity of medical spending, within age   & $0.573$  & $0.029$  & ENIGH 2016--2018 \\
Medical-spending share of expenditure, ages 25--34  & $0.0182$ & $0.0004$ & ENIGH 2016--2018 \\
Medical-spending share of expenditure, ages 50--64  & $0.0231$ & $0.0004$ & ENIGH 2016--2018 \\
Medical-spending share of expenditure, ages 65--79  & $0.0332$ & $0.0009$ & ENIGH 2016--2018 \\
Log-slope of mean medical spending, ages 25--79     & $0.0675$ & $0.0082$ & ENIGH 2016--2018 \\
\tgroup{4}{Earnings}
Skill premium in log earnings, ages 25--34          & $0.659$  & $0.013$  & ENIGH 2016--2018 \\
\tgroup{4}{Survival}
Five-year survival, low education, ages 50--64      & $0.951$  & $0.004$  & ENASEM \\
Five-year survival, high education, ages 50--64     & $0.972$  & $0.006$  & ENASEM \\
\bottomrule
\end{tabular}
\tabnote{0.94}{Survey standard errors are cluster-robust on primary sampling units (the log-slope's is a PSU-clustered bootstrap; the survival moments' are delta-method on clustered wave-to-wave survival). The hours target is a convention---a forty-hour week as one third of the time endowment---with no producing estimate behind it; it disciplines the labour-disutility scale and is not an estimated moment, which is why its standard-error cell is empty. No moment pairs one-to-one with one parameter; identification is discussed in the text.}
\end{table}

\paragraph{Samples, estimands, and measurement.} All ENIGH moments are computed on the pooled 2016 and 2018 waves---the Seguro Popular era contemporaneous with the ENASEM window---using survey expansion factors throughout, with age assigned by the household head's five-year band for household-level moments. Every spending moment is deflation-free by construction: shares are unit-invariant, the age slope is computed within wave, and the elasticities carry wave fixed effects. One measurement equation connects the survey to the model. ENIGH records \emph{out-of-pocket} spending, while the model's $m$ is the full resource cost of care; the model side of each spending moment therefore records $s\cdot m$ with $s=0.41$, the WHO Global Health Expenditure Database out-of-pocket share for M\'exico. This wedge is a recording convention, not a price: the household pays the full $m$ in its budget, and the public financing of the remainder is outside the model.\footnote{A skill-varying \emph{price} wedge---cheaper effective care for the insured---was measured and rejected: at the ceiling implied by insurance affiliation it closes only $15$--$18\%$ of the elasticity residual it was proposed to explain, the measured skill gradient of coverage is under $\pm10\%$ with an unidentified sign, and adopting it would open a second income$\to$health channel that manufactures part of the very gradient the paper holds out for validation.}

The two elasticities are Poisson pseudo-maximum-likelihood (PPML) regressions of medical spending \emph{in levels} on log labour income with wave fixed effects, on working-age heads with positive labour income and \emph{zeros retained}. The estimand choice is substantive. A log-log OLS conditional on positive spending---the familiar alternative---drops nearly half of the sample frame, while the model analogue drops essentially nothing: the two would be different objects sharing one name. PPML delivers the elasticity of the conditional mean with the zeros in, which is what the model computes. Income is labour income, matching the model's $w\,\nu\,\ell$, and the sample is working ages only, since retirees face a constant pension and carry no income variation to match. The within-age version repeats the regression band by band and averages the slopes with population-mass weights. The three spending shares are expansion-weighted means of the ratio of health spending to monetary expenditure---shares, with the model analogue $(1+\tau^m)m/[(1+\tau^c)c+(1+\tau^m)m]$---and the age slope is an equal-weight regression of log mean spending on the age band, computed within wave and averaged across waves. Because a slope is invariant to a common multiplicative shift, the slope constrains the \emph{shape} of the age profile while the 65--79 share pins its old-age level; the 50--64 share completes the profile.

The skill premium is the high-minus-low gap in log labour earnings among persons 25--34, from a saturated regression on age-band-by-sex cells; the model side is the exact closed-form counterpart of the same regression, so the two sides cannot drift apart in construction. Because the skill intercepts are first-step inputs, this moment functions partly as a consistency check on the first step; its information about $\varrho$ operates jointly with the health block, which determines the health gap the slope multiplies.\footnote{Both earnings moments condition on positive earnings in the data; the model has no participation margin, so its averages are taken over all states. Participation differs by education and sex, and the resulting selection is not corrected.}

The two survival moments are counted deaths: expansion-weighted three-year survival between ENASEM waves for persons 50--64, converted to the model's five-year frequency by the power $5/3$, split by completed education (below upper-secondary versus upper-secondary or more).
% TODO for Diego: confirm the raeducl 1 vs {2,3} split maps to
% "below upper-secondary vs upper-secondary or more" in the harmonized
% coding before this parenthetical is spliced.
These two moments have a status the rest of the table does not, and it is stated rather than blurred: with the education split of survival now a measured first-step input, the two levels overlap that input by construction, so they are not independent identifying information. They remain in the objective as \emph{disclosed consistency anchors}---what they verify is that the staged schedules, the health modulation, and the population composition combine to reproduce the levels the schedules were estimated from---and they are retained deliberately, because level anchors of exactly this kind have caught staging defects in the first-step pipeline that no other moment would have flagged. The overlap is not an identity: the input's schedules come from an unweighted estimator while the targets are expansion-weighted, and the two differ by about ten per cent on the 50--64 education gap, so the model does not hit the anchors mechanically. The education--survival \emph{gap} is retained as a validation diagnostic. Three further disclosures attach. First, the anchors route the health-correlated part of the education survival difference through the model's health channel, and the pooled pass anchors the model's average survival modulation to the modulation evaluated at the mean measured deficit index---the closest this calibration comes to a mapping between the model's internal health units and a measured index, routed through an input and stated here as a measurement choice. Second, the ENASEM education split and the model's stationary skill mix are different populations: the survey's high-education share at 50--64 is about $20\%$, while the model---which is stationary and cannot represent M\'exico's education expansion---carries over $50\%$ at every age. Third, the data convert three-year to five-year survival by a power map applied to group means; the Jensen wedge this induces is an order of magnitude below the targets' standard errors and partly cancels in the education difference.

\paragraph{What is deliberately absent.} No second-moment spending target appears in the table. Two candidates were constructed and measured---a cross-sectional coefficient of variation of medical spending from ENIGH, and a wave-to-wave autocorrelation of out-of-pocket spending from ENASEM---and both proved unreachable by the model over the entire admissible parameter box, for structural reasons documented with the validation diagnostics in Section~\ref{subsec:validation}. They are reported there against the model rather than targeted. The consequence for identification is stated here rather than discovered later: the two health-risk parameters $\bar\pi^h$ and $\varsigma$ have no dedicated identifying moment and are disciplined jointly by the moments that respond to them---chiefly the young and middle-aged spending shares and the survival levels.

\paragraph{Identification.} With nine parameters against ten moments the order condition holds with one over-identifying restriction---and that restriction is the consistency-anchor overlap disclosed above, not an economic restriction the theory imposes on the data, which is why no overidentification test is reported. The rank condition has been checked numerically only at the estimates of the superseded configuration from which the present estimation descends: there, a finite-difference Jacobian of the proportional-deviation map had numerical rank eight of ten, with the two near-degenerate directions loading on the survival levels and the young skill premium. That check has not been repeated at the estimates of Section~\ref{subsec:calibration_results}, and one feature of the current point keeps the local-derivative apparatus unreliable there, so we do not claim a rank condition at the reported estimates: two of the nine parameters sit at boundaries of the admissible set, where only one-sided variation exists. The broad division of labour, stable across every configuration measured, is as follows: the hours target disciplines $\Psi$; the three spending shares, the age slope, and the two elasticities jointly discipline the health technology $(\overline{H}_0, h^{\mathrm{slope}}, \zeta_h)$ and, through the level of spending, the amenity scale $\Xi$ and the flow value $\chi$, which has no assigned moment and is disciplined indirectly---estimation additionally excludes, by a smooth penalty, the region in which any household would assign negative value to survival; the skill premium disciplines $\varrho$ jointly with the health block; and the survival anchors, together with the spending profile, discipline the shock parameters. Three constraints on interpretation follow from this structure and are rules for reading every number in the next subsection. \emph{First}, earnings data identify only the product of $\varrho$ with the health gap between groups, and the scale of $h$ is pinned by the model's internal normalization $\bar h = 1$ rather than by any measurement; $\widehat\varrho$ is therefore a structural parameter meaningful relative to the model's health units, not a measurable return to health, and it is not comparable to regression coefficients in the health-economics literature. It is legitimate for counterfactuals, which is what it is for. \emph{Second}, the model's education--survival gap decomposes, as in Section~\ref{subsubsec:mortality}, into a component inherited from the measured baseline schedules and a component generated by the education gap in health investment; only the second is the mechanism, so we always report the decomposition, never the total gap alone, as evidence on it. \emph{Third}, a parameter estimate at a boundary of its admissible set is a diagnosis about the objective, not a measurement of the parameter; two of the nine estimates sit at boundaries, and Section~\ref{subsec:calibration_results} reads them as such.

\paragraph{Inference.} We report point estimates without standard errors, for two reasons that we prefer to state rather than paper over. The finite-difference Jacobian that the minimum-distance sandwich requires is not yet reliable: the derivative of the dominant moment is not sign-stable across candidate step sizes, and a standard error built on it would be confidently wrong. And two estimates sit at boundaries of the admissible set (Section~\ref{subsec:calibration_results}), where the unconstrained sandwich does not apply in any case. Resolving both is on the roadmap (Section~\ref{subsec:roadmap}); until then the fit table, not a $t$-statistic, is the measure of precision on offer.

\paragraph{Estimation is partial equilibrium; the model is not.} Parameters are obtained by solving the household problem at \emph{fixed} anchor prices $(r,w,\tau^p)$ rather than by clearing markets at every candidate $\Theta$. The device concerns how $\widehat\Theta_{\mathrm{MD}}$ is computed, not the model itself: the counterfactuals of Section~\ref{sec:experiments} are computed in general equilibrium with factor prices and the pension contribution rate responding, which is essential rather than incidental, since the fiscal amplification loop the paper is built to quantify operates through exactly those prices. After estimation, the model's full general equilibrium is solved at the estimates---including the re-bisection of the population growth rate to the dependency target under the resulting survival schedule---and every number reported in the next two subsections is evaluated at that equilibrium's own prices and matched demographics: the configuration the counterfactuals run in. The anchor at which the estimation ran and the resulting equilibrium are not identical; re-evaluating the estimates at the equilibrium's own prices improves the objective relative to the estimation anchor rather than degrading it, but the reported point is not an estimation--equilibrium fixed point to numerical tolerance, and re-estimating with the market-clearing loop inside the objective remains the natural robustness check. The outer optimizer is a derivative-free simplex search terminated by a no-further-improvement stopping rule rather than a formal convergence criterion. One reporting convention must be stated plainly: every number reported for this baseline---estimates, fit, and validation alike---is computed at the search discretization of the state space. The coarser search discretization and the finer one differ in the levels of some moments, so the reported levels carry that caveat; the polish on the finer discretization is deliberately deferred until the specification decisions of Section~\ref{subsec:roadmap} are settled, so that it is carried out once.

\subsection{Calibration Results}
\label{subsec:calibration_results}

Table~\ref{tab:theta_hat_md} reports the nine parameter values; Table~\ref{tab:moment_match_md} reports the fit of the ten targeted moments. Both are evaluated at the baseline's own general-equilibrium prices, with the population growth rate re-bisected to the dependency target, and at the search discretization (Section~\ref{subsec:secondstep_smm}). Standard errors are deliberately absent for the reasons given there. Two of the nine parameters sit at boundaries of their admissible sets; following the third interpretation rule of Section~\ref{subsec:secondstep_smm}, they are reported as corner diagnoses, not as estimates.

\begin{table}[H]
\centering
\caption{Minimum-distance parameter vector $\widehat\Theta_{\mathrm{MD}}$.}
\label{tab:theta_hat_md}
\small
\begin{tabular}{@{}llr@{}}
\toprule
 & \theadl{Description} & \theadr{Value} \\
\midrule
\tgroup{3}{Preferences and labour}
$\Psi$                & Labour-disutility scale            & $24.5$ \\
$\Xi$                 & Health-amenity scale               & $0.0017^{\dagger}$ \\
$\chi$                & Flow value of being alive          & $29.2$ \\
\tgroup{3}{Health technology}
$\overline{H}_0$      & Health-production scale            & $3.36$ \\
$h^{\mathrm{slope}}$  & Age slope of health production     & $-0.0188$ \\
$\zeta_h$             & Health-production curvature        & $0.489$ \\
\tgroup{3}{Health risk}
$\bar\pi^h$           & Mean shock arrival rate            & $0.0100^{\dagger}$ \\
$\varsigma$           & Shock size                         & $0.283$ \\
\tgroup{3}{Earnings}
$\varrho$             & Health--productivity slope         & $0.983$ \\
\midrule
$\xi$                 & Amenity curvature                  & $0.50$ (pinned) \\
\bottomrule
\end{tabular}
\tabnote{0.88}{No standard errors are reported (Section~\ref{subsec:secondstep_smm}, \emph{Inference}). $^{\dagger}$At a boundary of the admissible set: $\Xi$ and $\bar\pi^h$ at their lower bounds. Boundary values are corner diagnoses, not estimates, and are read as such in the text. $\xi$ is pinned at an interior value and disclosed as a convention: any admissible value delivers the same fit.}
\end{table}

\paragraph{The interior estimates.} The health-production scale is large while the curvature stays near one half ($\overline{H}_0=3.36$, $\zeta_h=0.489$), so medical care is productive but sharply diminishing---the configuration that lets heterogeneous households stop short of the health ceiling at heterogeneous points. The age profile of the technology is estimated to \emph{rise} with age: $\widehat h^{\mathrm{slope}}=-0.0188$ in the convention of~\eqref{eq:healthproduction} makes $\overline{H}_j$ grow by about $1.9$ per cent per five-year period, roughly a third over the life span, so the technology partially offsets the steeply rising depreciation schedule of the first step rather than compounding it---care is estimated to be somewhat more restorative at the ages where there is more to restore. The flow value of being alive is $\widehat\chi=29.2$; with the amenity scale at its floor (below), $\chi$, together with the survival and productivity returns priced in the first-order condition~\eqref{eq:foc_m}, carries the entire demand for medical care. The shock size $\widehat\varsigma=0.283$---more than a quarter of the health range when a shock lands---is identified conditional on the arrival rate, which sits at its floor, and inherits that corner's caveat. The health--productivity slope, at a boundary in every earlier pass, is now interior: $\widehat\varrho=0.983$, high in its unit interval---the earnings channel remains the dominant use of health---but an estimate rather than a corner diagnosis for the first time.

\paragraph{Two parameters at the boundary: diagnoses, not estimates.} The two boundary values are individually uninformative as point estimates, but jointly they are coherent, and that coherence is the finding: the estimator switches off both channels through which health matters other than earnings and survival. The amenity scale $\Xi$ sits at its lower bound, switching the direct utility value of health off ($\Omega(h)\to 1$); a diagnostic in which the bound was relaxed by two orders of magnitude saw the value fall a further twenty-four-fold while the objective stayed essentially flat in that direction---the corner is the data speaking, not the box. The shock arrival rate $\bar\pi^h$ sits at its floor, switching health risk off; we read this as censoring rather than as evidence that health risk is absent, since the dispersion and persistence facts that motivate the shock are second-moment objects the target set excludes for the measurement reasons of Section~\ref{subsec:secondstep_smm}, and the case for the stochastic specification rests on Section~\ref{subsec:health_block_evidence}, not on the fitted arrival rate. Both corners predate this baseline and survived the respecification of the education survival channel unchanged, so they are properties of the health block against these data, not artifacts of any one estimation pass; the specification question they pose is the binding item of Section~\ref{subsec:roadmap}. In the fitted model, health matters through what it earns and through how long it lets one live---not through how it feels, and not through the risk of losing it.

\begin{table}[H]
\centering
\caption{Moment fit at $\widehat\Theta_{\mathrm{MD}}$ (search discretization, general-equilibrium prices).}
\label{tab:moment_match_md}
\small
\begin{tabular}{@{}lrrr@{}}
\toprule
\theadl{Moment} & \theadr{Data} & \theadr{Model} & \theadr{Model/Data} \\
\midrule
\tgroup{4}{Labour}
Prime-age male hours (convention)               & $0.333$  & $0.326$  & $0.98$ \\
\tgroup{4}{Medical spending}
Income elasticity of $m$, pooled                & $0.575$  & $0.724$  & $1.26$ \\
Income elasticity of $m$, within age            & $0.573$  & $0.831$  & $1.45$ \\
Spending share, ages 25--34                     & $0.0182$ & $0.0211$ & $1.16$ \\
Spending share, ages 50--64                     & $0.0231$ & $0.0225$ & $0.97$ \\
Spending share, ages 65--79                     & $0.0332$ & $0.0203$ & $0.61$ \\
Log-slope of mean $m$, ages 25--79              & $0.0675$ & $0.0885$ & $1.31$ \\
\tgroup{4}{Earnings}
Skill premium, ages 25--34                      & $0.659$  & $0.717$  & $1.09$ \\
\tgroup{4}{Survival (consistency anchors)}
Five-year survival, low education, 50--64       & $0.951$  & $0.995$  & $1.05$ \\
Five-year survival, high education, 50--64      & $0.972$  & $0.998$  & $1.03$ \\
\bottomrule
\end{tabular}
\tabnote{0.9}{Model moments are exact functionals of the stationary distribution at $\widehat\Theta_{\mathrm{MD}}$, evaluated at the baseline's own general-equilibrium prices and dependency-matched population growth---the configuration the experiments run in---and at the search discretization. The education-survival \emph{difference} implied by the last two rows is right-signed but about one seventh of the measured gap; both facts are discussed in the text.}
\end{table}

\paragraph{Reading the fit.} Half of the moments land within ten per cent of their targets: hours ($0.98$), the 50--64 spending share ($0.97$), the young skill premium ($1.09$), and the two survival anchors in proportional terms ($1.05$, $1.03$). The misses are concentrated in the spending block, and they are one-directional in an instructive way: both income elasticities overshoot ($1.26$ pooled, $1.45$ within age), the age slope of spending overshoots ($1.31$), and the age profile of the spending \emph{shares} is tilted the wrong way---the young share overshoots by $16\%$ while the old-age share undershoots by nearly two fifths ($0.61$). Under the equal-weight proportional objective the residual is dominated by exactly these terms: the within-age elasticity and the old-age share together account for nearly two thirds of the remaining objective, while the two survival anchors carry under one per cent. One comparison must be stated candidly: the immediately preceding specification, whose education survival shifter posted a somewhat better objective and a closer spending fit, and part of that fit was purchased by the wrong-signed shifter itself---driving it deeper removed high-education mass at exactly the ages where spending is high, tilting the population composition toward the spending targets while inverting the education--survival gradient. That subsidy no longer exists by design, and the spending misfit reported here is carried honestly rather than bought back through a channel the model could not defend. One claim from an earlier stage of this project also does not survive and is corrected here: earlier configurations bracketed the elasticities between a deterministic-health undershoot and a stochastic-health overshoot, and the estimation was expected to reconcile them at the data value; at the adopted baseline they land above the data, so the bracket has narrowed but not closed---the model's income-responsiveness of medical spending remains too strong at the parameter values that fit the spending levels.

Figure~\ref{fig:lifecycle_baseline} plots the life-cycle profiles the estimated model produces, as cohort means by sex and skill at the baseline's own general-equilibrium prices. Labour supply rises from labour-market entry to a plateau across the forties and fifties, ordered by skill and by sex throughout. Assets are hump-shaped, peaking in the seventies for three of the four types and slightly later for high-skill women---highest for high-skill men, lowest for low-skill women---and drawn down thereafter. Consumption rises through working life in the same ordering and peaks late. Medical spending rises with age from entry to a peak in the seventies and falls at the oldest ages, where the remaining horizon over which repair pays is short. The health stock is the panel that most distinguishes this calibration: it stays within one per cent of its ceiling at every age and for every type, which is why panel~(f) repeats it magnified. The ordering inside that band is the expected one, high skill above low at every age, and its width is the quantitative content of the survival result above---differences in health investment do not open a health gap large enough to move mortality.

\begin{figure}[H]
\centering
\caption{Life-cycle profiles at $\widehat\Theta_{\mathrm{MD}}$ by sex and skill: (a) consumption, (b) labour supply, (c) assets, (d) medical spending, (e) health stock, (f) health stock magnified.}
\label{fig:lifecycle_baseline}
\includegraphics[width=0.95\textwidth]{Figures/lifecycle_baseline_educ_v1.pdf}
\fignote{0.92}{Cohort-mean profiles by five-year age period from the general-equilibrium solve at the estimates. Skill groups: \serone{high skill} and \sertwo{low skill}; solid lines men, dashed lines women. Labour is shown for working ages only. Panel~(f) repeats panel~(e) on a magnified vertical scale: across all four types and all ages the health stock spans $[0.990,\,1.000]$.}
\end{figure}

\subsection{Validation}
\label{subsec:validation}

The mechanism the paper isolates runs from income inequality to health inequality through endogenous investment, so the health--income gradient is an outcome the calibrated model should produce, not a moment to be targeted; targeting it and then reporting its magnitude as the strength of the mechanism would be circular. The same discipline governs everything in this subsection: no moment here entered the objective, and every comparison is between the model and an \emph{observable}---spending, earnings, counted deaths---never between the model's latent health index and a survey health measure, whose scales share no mapping.

\paragraph{Retired and held-out education diagnostics.} Table~\ref{tab:validation_educ} collects four diagnostics that were considered for the target set and deliberately kept out of it, evaluated at the estimated parameters.

\begin{table}[H]
\centering
\caption{Untargeted diagnostics at $\widehat\Theta_{\mathrm{MD}}$.}
\label{tab:validation_educ}
\footnotesize
\setlength{\tabcolsep}{4pt}
\begin{tabular}{@{}lrrl@{}}
\toprule
\theadl{Diagnostic} & \theadr{Model} & \theadr{Data} & \theadl{Reading} \\
\midrule
Education gap in 5-yr survival (H$-$L), 50--64      & $+0.003$ & $+0.021$ & right sign, one seventh; Section~\ref{subsec:calibration_results} \\
Growth of skill premium, 25--34 $\to$ 50--64        & $0.316$  & $0.319$  & carried by the measured profiles \\
CV of medical spending, singles 50--64              & $0.69$   & $5.70$   & no extensive margin; recall window \\
Autocorrelation of OOP spending, 50--64             & $0.81$   & $0.227$  & model spending is all signal \\
\bottomrule
\end{tabular}
\tabnote{0.94}{Data sources: survival gap and autocorrelation, ENASEM (the autocorrelation is the six-year correlation of log out-of-pocket spending among persons with positive spending at both dates); premium growth and coefficient of variation, ENIGH 2016--2018 (the CV is computed on single-person households, the one moment whose value depends on the aggregation unit, since household sums average individual shocks away). Model values at the search discretization, at the baseline's own general-equilibrium prices.}
\end{table}

Each row has a structural reading, and honesty about them cuts in different directions. The survival gap is right-signed for the first time---because the education split of baseline survival is now measured rather than estimated---and its shortfall in magnitude is dissected in Section~\ref{subsec:calibration_results}: the input's gradient is cap-clipped at the targeted ages and the endogenous investment channel contributes only $+0.0006$ (measured at these estimates by the pooled-schedule counterfactual). The growth of the skill premium with age is matched---but by construction rather than by estimation: it is carried by the measured skill-specific experience profiles of Section~\ref{subsec:firststep}, which is precisely why it was removed from the target set (targeting an input's own implication would be circular), and we accordingly report it without claiming it. The two second-moment spending diagnostics fail by an order of magnitude each, in opposite directions, and for reasons no parameter can repair. The cross-sectional CV is built on the survey's $57\%$ zero-spending mass, and the model produces no zeros for two compounding reasons: health production satisfies an Inada condition, so zero spending is never optimal for any household---the model lacks an extensive margin by functional form---and the survey's recall window is one quarter against a model period of five years, so ``spent nothing in three months'' and ``spent nothing in five years'' are different events ($0.57^{20}\approx 1.3\times10^{-5}$). The coefficient of variation is the one candidate moment that is not invariant to the model period, which is why it was excluded. The autocorrelation fails in the opposite direction: model spending is a deterministic policy of slow-moving states---assets, health, persistent productivity---so it over-persists relative to measured out-of-pocket spending, whose own flat lag profile indicates it is roughly $70\%$ transitory noise. Both failures say the same thing: the model's spending is the \emph{signal} component of measured spending, and second moments of spending compare model signal against data signal-plus-noise. Bringing such moments into scope requires a transitory-expenditure or measurement-error layer, which is on the roadmap; until then the boundary is disclosed rather than blurred.

\paragraph{Gender.} Nothing gender-specific is estimated: gender enters the model only through first-step inputs---the sex-specific efficiency profiles and baseline survival schedules and the sex-by-skill birth shares---and no structural parameter differs by sex. Every row of Table~\ref{tab:validation_gender} is therefore the model speaking for itself.

\begin{table}[H]
\centering
\caption{Gender validation at $\widehat\Theta_{\mathrm{MD}}$ (untargeted).}
\label{tab:validation_gender}
\small
\begin{tabular}{@{}lrr@{}}
\toprule
\theadl{Gender Object} & \theadr{Model} & \theadr{Data} \\
\midrule
\tgroup{3}{Earnings}
M$-$F log earnings gap, 25--34            & $+0.393$ & $+0.290$ \\
M$-$F log earnings gap, 35--49            & $+0.420$ & $+0.382$ \\
M$-$F log earnings gap, 50--64            & $+0.370$ & $+0.362$ \\
\tgroup{3}{Medical spending and survival}
F$-$M medical-spending share gap, 50--64  & $+0.0030$ & $+0.0018$ \\
F$-$M 5-yr survival gap, 50--64           & $+0.005$ & $+0.022$ \\
\bottomrule
\end{tabular}
\tabnote{0.94}{Data: ENIGH 2016--2018 (earnings and spending; earners, expansion-weighted) and ENASEM (survival). Model values at the search discretization, at the baseline's own general-equilibrium prices. The survival row is displayed for completeness but is \emph{not} claimed as validation: the sex-specific baseline schedules are first-step inputs from the same source, so the row is partly inherited, and even so the model delivers only a fraction of the measured gap---the same health-saturation compression that mutes the education gradient (Section~\ref{subsec:calibration_results}). No hours row appears: the survey's gender hours gap is an extensive-margin object (employment rates of $0.70$ for men against $0.40$ for women) and the model has no participation margin, so the comparison would be a category mismatch rather than a misfit.}
\end{table}

The earnings rows support the one claim this table still makes. The male--female earnings gap reproduces its age \emph{hump}---rising into 35--49 and easing at 50--64, as in the data---with levels overshooting throughout, most in the youngest band, an overshoot that most plausibly reflects earner selection embedded in the female profile input rather than the model; all three rows sit closer to the data at this baseline than at its predecessor. The spending-share row, by contrast, remains reported as a disclosure rather than a validation. Under an earlier parameter bundle the model reproduced the female--male gap in the medical-spending share at 50--64 nearly exactly; a subsequent recalibration overshot the measured gap by a factor of $2.2$, and the present baseline recovers about half of that degradation, overshooting by a factor of $1.6$. Part of the original move was definitional---the consumption tax sits in the share's denominator, and its recalibration from the statutory $0.16$ to the effective $0.09$ accounts for roughly six per cent of it---but the remainder is behavioural and undiagnosed, and no claim rests on this row until it is; it stays flagged for re-examination, with the improvement noted. The compressed survival row doubles as a marked out-of-sample test: any respecification that lets low-education health off the ceiling should reopen this gap toward its measured value without touching anything gender-specific.

\subsection{The Calibrated Baseline Equilibrium}
\label{subsec:calibrated_baseline}

Solving the full stationary general equilibrium at the estimates yields the calibrated Mexican economy of Table~\ref{tab:baseline_ge}. This is the equilibrium whose prices every number in Sections~\ref{subsec:calibration_results}--\ref{subsec:validation} is evaluated at, and it is the paper's current baseline.

\begin{table}[H]
\centering
\caption{Baseline general equilibrium at $\widehat\Theta_{\mathrm{MD}}$ (M\'exico, 2020): stationary aggregates.}
\label{tab:baseline_ge}
\small
\begin{tabular}{@{}lr@{}}
\toprule
\theadl{Quantity} & \theadr{Value} \\
\midrule
\tgroup{2}{Stationary aggregates and prices}
$K$ (capital, per cohort mass)              & $5.16$ \\
$L$ (effective labour)                      & $6.51$ \\
$Y$ (output)                                & $9.49$ \\
$r$ (annual)                                & $8.58\%$ \\
$w$ (wage)                                  & $0.875$ \\
$\tau^p$ (PAYG contribution rate)           & $5.66\%$ \\
$\overline{pen}$ (pension flow per retiree) & $0.475$ \\
$C/Y$                                       & $61.1\%$ \\
$M/Y$                                       & $3.21\%$ \\
$N^R/N^W$ (PAYG dependency)                 & $0.141$ \\
$5K/Y$ (annual capital--output ratio)       & $2.72$ \\
$5B/Y$ (annual debt--output ratio)          & $112\%$ \\
\bottomrule
\end{tabular}
\tabnote{0.92}{Stationary equilibrium at the estimates of Table~\ref{tab:theta_hat_md}, at the search discretization. $\tau^p$ clears the PAYG block and $B$ is the residual that closes the general budget, so neither is a free parameter. Lifetime utility at birth is deliberately not tabulated: its level is not interpretable, and only its \emph{change} under a counterfactual carries welfare content (Section~\ref{subsec:welfare_incidence}).}
\end{table}

Four of these deserve comment. The old-age dependency ratio comes out at $0.1414$, equal to the 2020 census value---but this is matched by construction, not validated: the population growth rate is re-bisected to the dependency target whenever the survival schedules change (Section~\ref{subsec:firststep}), so the row records a discipline of the calibration, not evidence for the model. The annual capital--output ratio is $2.72$ against a Mexican value near $3.29$---a $17\%$ shortfall. This is the disclosed limitation of the head of this section, and its movement is itself a result: the shortfall was $22\%$ before the education-specific survival schedules entered, and the longevity-savings channel---longer high-education lives strengthening life-cycle accumulation---recovered about a fifth of the gap without any recalibration aimed at it. The investment rate now overshoots, at $22.6\%$ of output against $20.2\%$ in the data, a $2.4$-percentage-point miss that moves with the higher population growth rate and the longer lives; it is disclosed rather than chased. The implied public debt, at $112\%$ of annual output against a Mexican federal debt near $50\%$, remains the residual of the closure---$B$ absorbs whatever stationary stock balances the primary account---and is worse at this baseline than at its predecessor; it is a disclosed limitation, deliberately not chased, and the natural refinement remains targeting a debt ratio and letting a tax instrument clear the budget.

The experiments of Section~\ref{sec:experiments} are computed at exactly this equilibrium. An earlier solve, generated by the superseded six-parameter $\widehat\theta$ of Appendix~\ref{app:calibration}, inherits in full the specification and first step listed at the head of this section (p.~\pageref{par:superseded})---including the earlier health-depreciation schedule and the earlier fiscal and technology block; its stationary aggregates are reported with that provenance in Appendix~\ref{app:provenance}, and are kept out of the main text so that one baseline equilibrium, the one above, is the paper's. One result of that earlier solve is reported here rather than in the appendix, because it identifies which force governs the aggregate capital stock rather than pinning the level of any aggregate; whether its magnitudes survive re-estimation has not been tested.

Aggregate capital, finally, is governed by income persistence first and population composition second, and the baseline makes this quantitative. The calibrated economy holds less than half the capital of a comparison economy that shares its calibrated preferences and health technology but carries the higher AR(1) persistence and symmetric birth shares of a literature parameterization: $K\approx 10.5$ there against $4.96$ here. Table~\ref{tab:k_decomp} traces the $-53\%$ move to its two sources through a four-corner experiment in which each cell is a separate general-equilibrium solve at the calibrated parameters, varying only the productivity process and the birth-share vector. The Shapley decomposition attributes $-3.64$ of the $-5.53$ fall to the income process and $-1.88$ to cohort composition. Persistence dominates. Lowering $\rho$ from $0.98$ to $0.782$ turns near-permanent income shocks into transitory ones, and the buffer-stock saving motive that near-unit-root risk generates collapses---even though the innovation variance rises enough to lift the unconditional variance from $0.063$ to $0.181$---because transitory risk is cheap to self-insure through consumption smoothing. Cohort composition works in the same direction but more weakly: M\'exico's empirical birth shares place $77.5\%$ of cohorts on the low-skill types, against $50\%$ under symmetry, and low-skill cohorts earn, save, and supply less labour. The aggregate capital stock of a calibrated emerging-economy life-cycle model is thus a statement about the persistence of idiosyncratic income risk before it is a statement about the skill composition of the population.

\begin{table}[H]
\centering
\caption{Aggregate capital under the income process and birth-share legs (four-corner decomposition). \emph{Superseded solve} (p.~\pageref{par:superseded}); not the current baseline of Section~\ref{subsec:calibration_results}.}
\label{tab:k_decomp}
\small
\begin{tabular}{@{}lcc@{}}
\toprule
\theadl{Aggregate capital $K$} & \thead{Symmetric $\pi_{\rm birth}$} & \thead{Asymmetric $\pi_{\rm birth}$} \\
\midrule
High persistence ($\rho=0.98$, $\sigma_\varepsilon=0.05$)   & $10.49$ & $8.21$ \\
Data persistence ($\rho=0.782$, $\sigma_\varepsilon=0.265$) & $6.45$  & \textbf{4.96} \\
\bottomrule
\end{tabular}
\tabnote{0.92}{Each cell is a stationary general-equilibrium solve at the calibrated parameters and the baseline closure, varying only the AR(1) process and the birth-share vector; the bold entry is the calibrated baseline. Total move $\Delta K = 4.96-10.49=-5.53$ ($-52.7\%$). The symmetric Shapley split attributes $-3.64$ to the income-process re-anchoring and $-1.88$ to the birth-share asymmetry, exhausting the move; the two channels are mildly sub-additive. The income-process leg is the larger of the two.}
\end{table}

This is the economy the reader should carry into Section~\ref{sec:experiments}: a calibrated M\'exico in which buffer-stock saving is weak because idiosyncratic income mean-reverts, the contribution rate clears the pension system at $5.66\%$, and the four types are separated by the ENIGH skill premium of the first step (Table~\ref{tab:firststep_params}). The experiments perturb exactly this object.


%% ====================================================================
\section{The Health Block: What the Estimated Apparatus Delivers}
\label{sec:healthblock}

\subsection{What Mexican Households Buy When They Buy Medical Care}
\label{subsec:health_evidence}

One convention governs the evidence in this section, and it is stated here once, for the record. The household expenditure survey from which the spending facts are drawn, ENIGH, records out-of-pocket medical spending and captures between 22 and 34 per cent of the out-of-pocket total recorded in the national health accounts. The survey is accordingly used throughout this paper for shapes, gradients, and composition, and never for levels: no peso figure from it is reported as a national total, and nothing below rests on the absolute size of any spending flow. The rule has been followed from the first-step estimation onward---the profiles and spending moments of Section~\ref{sec:calibration} are shares, slopes, and ratios, not totals---and writing it down here makes explicit what the calibration has practiced throughout.

The first fact is composition. Classifying every out-of-pocket expenditure item by an auditable rule into unambiguously preventive, unambiguously curative, and ambiguous yields an honest bracket rather than a point estimate: 2.0 per cent of spending is unambiguously preventive, 53.6 per cent is unambiguously curative, and the remaining 44.5 per cent---over-the-counter medicines, dental and optical care, and other items whose intent the expenditure classification cannot resolve---is reported as ambiguous rather than allocated by assumption (Figure~\ref{fig:mot_composition}). The bracket licenses one verbal summary and one only: Mexican out-of-pocket medical spending is majority curative, under any assignment of the ambiguous block. The classification cannot say how much of the ambiguous mass is prevention; it does not need to, because even the assignment most generous to prevention leaves curative spending the majority.

\begin{figure}[H]
\centering
\includegraphics[width=0.95\textwidth]{Figures/fig07_composition.pdf}
\caption{Composition of out-of-pocket health spending by income decile under the
stated curative/preventive classification. Every expenditure item is assigned to
one of three classes by an auditable rule; the ambiguous block (over-the-counter
medicines, dental, optical, chronic-disease maintenance, unqualified diagnostics,
insurance premiums) is reported rather than allocated. Assigning it entirely to
prevention leaves curative spending the majority; assigning it to treatment makes
spending 98.1 percent curative. Because the 2024 COICOP recode dropped
prescription status, the composition is comparable across deciles within a wave
but not across waves. Source: ENIGH 2024 (INEGI), authors' calculations.}
\label{fig:mot_composition}
\end{figure}

The expenditure classification is corroborated from the individual side by an instrument of entirely different construction. Where the composition bracket reads receipts, a direct survey question reads motives: asked about the last health need, households in the poorest socioeconomic tercile report preventive motives---vaccination, check-ups, prenatal control---for 11.2 per cent of needs, rising only to 15.9 per cent in the richest tercile. The two instruments differ in what they measure, in how they measure it, and in the unit of observation, and they agree on the substantive point: Mexican medical care is predominantly a response to realized illness, and prevention is a minority of care even where resources are least scarce.

The capacity to finance that response is, finally, deeply unequal. Catastrophic health expenditure reaches 3.9 per cent of households in the poorest decile, against 1.4 to 1.9 per cent in deciles four through ten. The implication for what follows is structural rather than decorative: if measured spending is predominantly a response to realized illness, and the capacity to finance that response is distributed this unequally, then a specification in which restoration is unconstrained---available to every household on the same terms whenever a draw lands---is a strong assumption rather than a neutral one.

With the evidence on the record, the remainder of the section reports what the estimated health block delivers, what it does not deliver, and why.

\subsection{The Value of a Statistical Life}
\label{subsec:vsl_structural}

The record opens on the delivery side: this subsection reports how the estimated model prices survival against an external benchmark.

The flow value of being alive $\chi$ exists because a bare GHH--CRRA felicity cannot price survival: its normalization places the flow utility of a living household at the same zero assigned to death, mechanically delivering a near-zero willingness to pay for survival that no amenity scale repairs (Section~\ref{subsec:preferences}). An earlier version of this calibration targeted the value of a statistical life to discipline $\chi$. We no longer do. The natural target is a benefit transfer whose standard error is a judgment band rather than a sampling error, and so soft a target functions as a normalization rather than a restriction; $\chi$ is instead estimated jointly with the rest of $\Theta_{\mathrm{MD}}$, disciplined indirectly through the spending moments it shifts, with a penalty excluding the region in which any household would assign negative value to survival.

The model's VSL is then reported as an \emph{output}. The estimand is a central statistic---the mass-weighted trimmed mean over the middle fifth of the population distribution of the marginal rate of substitution between mortality risk and money---expressed scale-free as a multiple of annual output per capita, the form in which the benefit-transfer literature's unit-income-elasticity recommendation makes the number comparable across countries. At the estimates, evaluated at the baseline's own general-equilibrium prices, the model delivers a VSL of $179$ times annual output per capita, against the benefit-transfer reference of $152$ \citep{viscusi_masterman_2017}: the same order of magnitude, about a sixth above. The number is computed at the search discretization and will move with the deferred discretization polish (Section~\ref{subsec:secondstep_smm}). No household carries a negative value of survival at the estimates. The welfare statements in this paper rest on consumption and the labour wedge rather than on the monetized value of survival and do not depend on this number.

\subsection{Why the Baseline Carries Health Risk}
\label{subsec:health_block_evidence}

Before the failure can be read correctly, the architecture must be on the record: this subsection reproduces the evidence that forced the baseline to carry health risk rather than deterministic depreciation alone.

Section~\ref{subsubsec:health} adopts deficit-scaled health production and the health shock as the baseline, with the constant-returns deterministic rule~\eqref{eq:health_formation} as the nested limit. The choice was forced by quantitative evidence, and the evidence belongs here rather than in the model section.

Under the deterministic constant-returns rule, with the depreciation schedule estimated in Section~\ref{subsec:firststep}, the model's working-age health distribution is degenerate: over $99\%$ of working-age mass sits exactly at the ceiling $\bar h$, because at a five-year depreciation of a tenth of one per cent any household that spends at all is returned to full health. Two failures follow, one internal and one observable. Internally, a mechanism that transmits income inequality into health inequality has no variation to operate on---within-cohort health dispersion is zero by construction over most of the life cycle. Observably, the same degeneracy destroys the level of medical spending: with nothing to repair, the young spending share collapses to roughly three per cent of its data value (a share of $0.0005$ against a target of about $0.02$).

The two departures from the deterministic rule are complements, and neither alone suffices. Deficit-scaled repair makes the marginal product of care fall with health, which restores an interior spending motive and repairs the level of spending---but without risk, households of a cohort decline together and the cross-section stays compressed. The shock creates within-cohort dispersion---but under constant-returns repair the ceiling reabsorbs it, since any struck household spends its way straight back to $\bar h$. Jointly, in a pre-estimation probe of the specification, working-age mass at the ceiling fell from over $99\%$ to about $75\%$ and the young spending share moved from three per cent of its target to within a few per cent of it. That configuration is what the estimation of Section~\ref{subsec:secondstep_smm} then disciplined; the probe's hand-set risk parameters are superseded by the values of Table~\ref{tab:theta_hat_md}, including the boundary value of the arrival rate discussed there.

\subsection{What the Investment Block Cannot Generate}
\label{subsec:capacity_result}

For the first time in any configuration of this model, the education gap in five-year survival at ages 50--64 comes out right-signed: $+0.30$ percentage points against $+2.09$ in the data---the correct ordering at about one seventh of the measured magnitude. The sign is earned by measurement, not by estimation: the education difference in baseline survival is now a first-step input, and the decomposition of Section~\ref{subsubsec:mortality}---computed at these estimates by holding every policy and the population distribution fixed and substituting the pooled schedule---attributes $+0.24$ percentage points of the $+0.30$ to the measured schedules. The endogenous investment$\to$longevity channel---the paper's mechanism---contributes $+0.06$ of a percentage point, in line with the $+0.05$ measured at every predecessor parameter vector; it is structural, and it is disclosed as such. The shortfall in magnitude also has a located cause rather than an open diagnosis: the input's education gradient at ages 50--69 is clipped by the survival estimator's ceiling cap (Section~\ref{subsec:firststep})---entirely for women---and the near-ceiling modulation leaves no endogenous margin to replace what the cap removes, so the measured gradient passes through only partially at exactly the targeted ages. The age-varying floor that would un-clip the input is the deferred fix. What the model delivers beyond the sign is the integrated ordering: expected age at death is higher for the high-education group by $2.5$ years for men and $1.9$ years for women, and that ordering---not any education-mortality \emph{level}, which this paper does not claim---is the object the pension-incidence results use.

The retirement of the education survival shifter rests on a capacity result worth recording. A sweep of the shifter over its admissible range showed that the modulation's cap binds, for any non-negative shifter, across some ninety-nine per cent of the high-education mass at the relevant ages---whose health sits near the ceiling---so the education gap the shifter could produce saturates at about four per cent of the data gap; meanwhile the low-education survival level is pinned at its baseline-schedule ceiling of about $0.996$, against $0.951$ in the data, at every parameter vector ever evaluated. A parameter structurally incapable of representing the positive channel it was created for, and rewarded by the objective only for composition effects in the spending moments, is not an education-mortality parameter, and it was removed. The ceiling itself is behaviorally defended rather than an artifact of any single estimate: a sweep of the health technology's age slope over its full box cannot break it---making repair costlier with age induces households to spend several-fold more to stay at the ceiling rather than let health fall, and the spending data reject those configurations by an order of magnitude in the objective. Producing the data's low-education survival level would require substantial 50--64 mass to sit well below the health ceiling, which is a re-specification of the health block, not a choice of moments or bounds; that decision is the binding item of Section~\ref{subsec:roadmap}.

Read constructively, the estimation delivers a measured statement about capacity: at parameter values that fit the spending, elasticity, earnings, and survival facts, the investment mechanism as currently specified generates essentially none of the observed education--survival gradient---working-age depreciation is too small for investment differences to open a health gap before the ages at which mortality operates, and the survival cap removes the longevity return to investment for the healthy. No counterfactual in this paper attributes the education--survival gap to the investment channel.

\subsection{A Mismatch of Objects}
\label{subsec:diagnosis}

The capacity result has a diagnosis, and it is a statement about objects rather than about parameters. Medical spending in the model is committed before the health shock is realized---the timing convention of Section~\ref{sec:model}---and is therefore a preventive object by construction: it builds the stock in advance of illness, on a margin every household operates every period. Measured Mexican spending, on the evidence that opened this section, is predominantly restoration conditioned on a realized draw: care sought after illness arrives, financed out of pocket by households whose capacity to finance it differs sharply. These are not two descriptions of one object. They differ in timing---one is chosen before the draw, the other because of it; in exposure to expenditure risk---committed spending is smooth where conditioned spending is lumpy; and in how they connect spending gradients to survival gradients---a preventive gradient compounds over decades of stock-building, while a restorative gradient operates only when and where shocks land.

The estimation then does what estimation does: it fits the moments it is given with the objects it has. A block containing only the preventive object, disciplined on spending moments generated largely by the restorative one, must reproduce the level and age profile of spending with an instrument that operates on a different margin, and it resolves the tension by holding the health stock near its ceiling for every type---the configuration in which committed spending is high enough to match the data while health variation is too small to matter. The spending fit and the absent survival gradient of Section~\ref{subsec:capacity_result} are one outcome, not two: the ceiling is how a preventive block wears restorative clothes.

The result is a measurement, not a defect. Nothing in it says the Grossman-descended investment block is wrong as theory, and nothing in it says the implementation failed: at the estimated parameters the block does what its structure permits with the data it faces. What it says is that the investment block, taken to emerging-economy micro-data and disciplined on the spending, elasticity, earnings, and survival facts at once, cannot generate the observed education--survival gradient, because the spending it is asked to rationalize is by majority not the kind of spending it models. The result could not have been obtained from a partially calibrated model: with any moment released, the estimator would have found a configuration that hides the tension. Only the full discipline forces the collision, which is why the paper reports the estimation and the capacity result as a single contribution.

\subsection{Refinements Reserved for the Next Stage}
\label{subsec:roadmap}

What follows from the diagnosis is a decision about specification rather than a repair, and this closing subsection records where it binds.

Two items the earlier draft reserved have since closed. The general-equilibrium solve at the estimates has been carried out and is now the paper's own baseline, at whose prices every reported number is evaluated (p.~\pageref{par:superseded}). And the education--survival design question has been resolved by measurement rather than by targeting: reweighting the survival levels was refuted empirically, a gap target was shown unreachable from the baseline region, and the survival shifter that was meant to carry the channel proved structurally incapable of it and was retired in favour of the education-specific baseline schedules of the first step (Section~\ref{subsec:calibration_results}); the two survival levels remain in the objective as disclosed consistency anchors.

What remains open is of three kinds. \emph{Computation.} Every number in this calibration is reported at the search discretization; the polish on the finer discretization and the mechanical adoption of the sourced first-step constants are deliberately deferred until the specification decision below is settled, so that each is done once; the experiments of Section~\ref{sec:experiments} are computed at the current baseline and inherit the same discretization caveat. \emph{Inference.} Standard errors await a finite-difference step-size study and a treatment of the two boundary estimates; until both are settled we report none. \emph{Specification.} The binding question is the health block. The capacity result of Section~\ref{subsec:calibration_results} says the investment mechanism as specified cannot generate the education--survival gradient, and the two boundary values say the objective is switching off the amenity and risk channels, leaving health to matter only through earnings and survival. The options are to relocate the mechanism to old-age health deterioration, where the estimated depreciation is large; to re-specify the health block so that substantial 50--64 mass sits below the health ceiling; or to reframe the paper's contribution around the channels the estimation does support. Behind that decision, four data-side refinements retain their place: an age-varying survival floor to un-clip the measured education gradient of baseline survival at ages 50--69 (Section~\ref{subsec:firststep}); a transitory-expenditure or measurement-error layer to bring second-moment spending facts into the model's scope (Section~\ref{subsec:validation}); an hourly-wage first step to replace the GHH-inversion correction of Section~\ref{subsec:firststep}; and the held-out ENIGH waves (2022--2024) as an out-of-sample check on the out-of-pocket measurement wedge, whose value is anchored to the Seguro Popular era. The capital--output shortfall is deliberately absent from this list: it is a disclosed limitation of the calibrated configuration (p.~\pageref{par:superseded}), not future work.

%% ====================================================================
%% SECTION 6 — EXPERIMENTS (at the calibrated baseline)
%% ====================================================================
\section{Experiments}
\label{sec:experiments}
\label{sec:policy}

M\'exico's old-age dependency ratio is projected to more than double between 2020 and 2050, from $0.141$ to $0.307$ \citep{CONAPO2023}. In the model this projection has an immediate fiscal translation, because the PAYG identity~\eqref{eq:pension_contribution} ties the contribution rate to the dependency ratio one for one: at the baseline replacement rate $\kappa=0.40$, a dependency ratio of $0.307$ means a contribution rate of $12.3\%$, against $5.66\%$ today. The question this section asks is what that arithmetic costs, who pays it, and what the two classical pension margins---cutting benefits and raising the retirement age---buy relative to doing nothing.

We answer it with three general-equilibrium experiments on the calibrated baseline of Section~\ref{subsec:calibrated_baseline}. All three impose the same 2050 demographics. The first, \emph{inaction} (CF0), changes no policy and lets the contribution rate absorb the entire demographic shift. The second, the \emph{benefit margin} (CF1), cuts the replacement rate until the contribution rate returns to its 2020 value. The third, the \emph{retirement margin} (CF2), raises the retirement age from 65 to 70 and leaves benefits untouched. The PAYG identity organizes the whole section: since $\tau^p = \kappa\, N^R/N^W$, a government facing a doubling dependency ratio must give up one of three things---the contribution rate, the replacement rate, or the boundary between $N^W$ and $N^R$---and the three experiments move exactly one each. Inaction is not the absence of a policy; it is the choice of the contribution rate as the adjusting margin, and we treat it as the tax arm of the menu.

Every experiment is a stationary equilibrium computed on the unmodified baseline code path: the estimated parameter vector $\widehat\Theta_{\mathrm{MD}}$ and every externally sourced parameter are held fixed, and the contribution rate, factor prices, the pension benefit, and the stationary debt stock all clear endogenously in each economy. At the baseline inputs the counterfactual driver reproduces the baseline equilibrium of Table~\ref{tab:baseline_ge} exactly, field for field, so every difference reported below is attributable to the imposed change alone. All comparisons are between stationary equilibria---no transition path is computed---and all equilibria are solved at the search discretization of Section~\ref{subsec:secondstep_smm}; both qualifications matter for how the results may be read, and we return to them where they bind.

\subsection{Design}
\label{subsec:cf_design}

\paragraph{The aged economy.} Aging enters through the population growth rate alone. The baseline disciplines $n_p$ so that the model's stationary age distribution reproduces the 2020 census dependency ratio (Section~\ref{subsec:firststep}); the aged economy re-solves the same bisection against CONAPO's projected 2050 ratio of $0.3065$, which lowers $n_p$ from $3.51\%$ to $1.48\%$ per year. The counterfactual $n_p$ is therefore not a fertility forecast: it is the growth rate at which the model's stationary demography matches the projected 2050 population structure, exactly as the baseline $n_p$ matches the 2020 structure, and both ends of the comparison are disciplined by the same source (Appendix~\ref{app:demog2050}).

We deliberately do \emph{not} impose projected 2050 survival schedules, for a reason specific to this model: survival is partly endogenous. The measured schedules $\psi^{\circ\,g,\theta}_j$ are modulated by the health stock through~\eqref{eq:survival_modulation}, and the health stock responds to prices and income. A ``2050 survival schedule'' imposed from outside would therefore bake the behavioral responses of some other equilibrium into what is meant to be an exogenous environment change, and the counterfactual would no longer be clean. Fertility decline is in any case the dominant driver of M\'exico's projected aging, so the fertility channel carries the experiment; the survival schedules, and the birth composition $\pi^{g,\theta}_1$, stay at their 2020 values throughout, and we disclose both choices wherever the results depend on them.

\paragraph{The two reforms.} Both reforms share the aged economy's demographics, so each is read against inaction, not against 2020. In CF1 the replacement rate is set analytically from the PAYG identity---the value that returns the equilibrium contribution rate to its baseline $5.658\%$---and confirmed in equilibrium; the resulting $\widehat\kappa$ is itself a result---it measures how much benefit must be given up to hold the contribution rate fixed against a doubling dependency ratio. In CF2 the retirement period moves from $j_R=10$ to $j_R=11$, so households work through ages 65--69, and $\kappa$ stays at $0.40$; the contribution rate is left free, and where it lands measures how much of the fiscal cost of aging longer working lives recover. One input must be added for CF2: the age-efficiency profiles of Section~\ref{subsec:firststep} end at ages 60--64, because the ENIGH estimation sample ends at the baseline retirement age. We extend each of the four profiles to ages 65--69 with the same regression re-estimated on ENIGH earners aged 20--69, taking only the within-regression 65--69 to 60--64 growth ratio per type, so that the staged profiles at all other ages are untouched. The estimated ratios lie between $0.87$ and $0.96$---late-life efficiency declines $4$--$13\%$ from its 60--64 level, continuing the measured trend---but they come from the roughly one in five persons of those ages who still work ($19.2\%$ employment at 65--69, against $28.5\%$ at 60--64), a selected group. The extension is therefore an upper bound on the productivity of the marginal worker the reform forces into employment, a bias we return to when reading CF2's incidence.

\paragraph{Why these three, and not others.} The menu is short because the identity that generates it is short: $\tau^p = \kappa\, N^R/N^W$ has three arguments and no others, so within the pension system the three experiments are a closed taxonomy of responses, not a sample of interesting policies. What lies outside it does so for stated reasons. Financing the pension gap from general revenue moves no argument of the identity; it relocates the obligation to a different budget constraint and raises a fiscal-sustainability question the paper's closure does not address (Section~\ref{subsec:exp_caveats}). Restructuring the benefit rule---linking benefits to individual contribution histories, or changing indexation---changes the object rather than the margin: the model's benefit is flat and wage-indexed, as Section~\ref{sec:model} discloses, and a contributory link would be a different model. Coverage expansion through formalization moves the contribution base itself and is designed and left as an extension (Section~\ref{subsec:lit_informality}). Combinations, finally, are omitted deliberately: their fiscal arithmetic is additive and can be read off the identity, but their incidence is not---each reform changes the prices through which the others operate---and studying the margins one at a time is what makes each one's incidence attributable.

\paragraph{What the household sees.} One structural fact disciplines the interpretation of the first two experiments: neither $n_p$ nor $\kappa$ appears anywhere in the household problem of Section~\ref{subsec:optimization}. The household sees its environment only through the price-and-fiscal bundle $(r, w, \tau^p, \overline{pen})$. Holding that bundle at baseline values, the direct effect of aging and of the benefit cut on any household is identically zero; the entire incidence of CF0 and CF1 is general-equilibrium, transmitted through factor prices and the pension system. Only the retirement-age reform reaches the household problem directly, by changing the ages at which it can earn. Aging, in this economy, hurts nobody except through prices and the fiscal system---which is precisely why the experiments must be solved in general equilibrium rather than priced at baseline.

Before any experiment was solved, the PAYG identity fixed three predictions: $\tau^p(\mathrm{CF0}) \approx 0.40 \times 0.3065 = 12.3\%$, $\widehat\kappa(\mathrm{CF1}) \approx \tau^p_{2020}/0.3065 \approx 0.185$, and, with band 10 moved from retirees to workers, $\tau^p(\mathrm{CF2}) \approx 8.8\%$. The equilibria confirm all three to the reported digits. The identity, however, is silent on prices, quantities, and welfare---everything in the next four subsections is what the general equilibrium adds to the arithmetic.

\subsection{The Welfare Metric}
\label{subsec:welfare_incidence}

Every welfare statement in this section answers one question: how much consumption would a newborn give up to be born into one economy rather than the other? The answer, $\lambda$, is the constant share of consumption---surrendered at every age and in every state of life---that leaves the better economy's newborn exactly as well off as the other's; a $\lambda$ of five per cent means giving up one twentieth of consumption at every age, for life. It is a lifetime object, not a period one, computed along the newborn's own equilibrium path rather than a hypothetical one. Formally, welfare comparisons across stationary economies are stated as consumption equivalents in the compensating direction of \citet{DeNardiPashchenkoPorapakkarm2017}. For two economies $A$ and $B$ and a newborn of type $(g,\theta)$ whose expected lifetime utility is higher in $A$, we report the constant fraction $\lambda_{g,\theta}$ of consumption that the newborn of the \emph{better} economy would give up, at every age and state along its own equilibrium allocation, to be exactly as well off as the newborn of the other economy:
\begin{equation}
\label{eq:cev_newborn}
\E_1\!\left[\sum_{j=1}^{J} \beta^{j-1}\!\left(\prod_{i=2}^{j}\psi_{i}^{g,\theta}(h'_{i-1})\right) u\bigl((1-\lambda_{g,\theta})\,c_j^{A},\,\ell_j^{A},\,h_j^{A}\bigr)\right] \;=\; \mathcal{W}_1^{B}(g,\theta),
\end{equation}
where $(c^A,\ell^A,h^A)$ is the equilibrium allocation of economy $A$ held fixed. The direction is not a stylistic choice. With $\gamma>1$ the felicity~\eqref{eq:utility} is bounded above in consumption---as $c$ grows the bracket converges to $\chi + 1/(\gamma-1)$---but unbounded below as consumption falls toward the labor-disutility floor, so a proportional \emph{reduction} on the better side always reaches the target value, while the familiar increment on the worse side may not: there are gaps no amount of extra consumption can close, because part of what the better economy delivers is survival, valued through $\chi$, which extra consumption cannot buy. Equation~\eqref{eq:cev_newborn} is finite by construction and counts the survival value in full. We report $\lambda$ per type and aggregated with the birth shares $\pi^{g,\theta}_1$, signed so that $\lambda>0$ means the first-named economy is the better one.

The same object is computed by age. For each five-year band $j$, $\lambda_j$ scales the remaining-lifetime consumption of economy $A$'s age-$j$ residents, averaged over that economy's own stationary distribution at age $j$, until their remaining-lifetime value equals that of the other economy's same-age residents. Two rules govern how these profiles may be read, and we apply both throughout. First, $\lambda_j$ compares \emph{residents of different steady states} at the same age; it is not the transition incidence on any cohort alive at a reform date, because no transition is computed. In plain terms, these profiles price destinations, not journeys: no person alive at a reform date experiences the numbers reported here. Second, at old ages consumption is a weak welfare instrument: over a short remaining horizon a large share of remaining value is the flow value of being alive, which consumption scaling does not touch, so closing a given value gap requires a mechanically large $\lambda_j$. Put plainly: near the end of life much of what remains valuable is simply being alive, which consumption cannot buy, so closing even a small value gap takes a large consumption adjustment. The \emph{age gradient} of the profile is the reportable object; the old-age \emph{levels} carry this amplification and we do not lean on them.

One principle organizes every welfare number that follows. A pay-as-you-go pillar pays an implicit return equal to the growth rate of the contribution base, while private saving pays the market return \citep{Aaron1966}. Wherever the market return exceeds the economy's growth rate---and it does, in every economy on this menu---mandatory participation in the unfunded pillar transfers resources away from the participant; the size of the gap between the two returns is what every welfare comparison below ultimately measures. Reforms that shrink the mandatory pillar move newborns toward the higher return; forces that move factor prices move the gap itself.

\subsection{The Cost of Inaction}
\label{sec:aging}

Table~\ref{tab:cf_aggregates} reports the three counterfactual equilibria against the baseline; this subsection reads the second column.

\begin{table}[H]
\centering
\caption{The counterfactual equilibria: 2020 baseline, inaction, and the two reforms at 2050 demographics.}
\label{tab:cf_aggregates}
\small
\setlength{\tabcolsep}{5pt}
\begin{tabular}{@{}lcccc@{}}
\toprule
 & \thead{Baseline} & \thead{CF0} & \thead{CF1} & \thead{CF2} \\
 & \thead{2020} & \thead{Inaction} & \thead{Benefit cut} & \thead{Retire at 70} \\
\midrule
\tgroup{5}{Imposed}
$n_p$ (annual, \%)            & $3.51$  & $1.48$  & $1.48$  & $1.48$  \\
$\kappa$ (replacement rate)   & $0.40$  & $0.40$  & $0.185$ & $0.40$  \\
Retirement age                & $65$    & $65$    & $65$    & $70$    \\
\tgroup{5}{Pension block (endogenous)}
$N^R/N^W$                     & $0.1414$ & $0.3064$ & $0.3065$ & $0.2195$ \\
$\tau^p$ (\%)                 & $5.66$  & $12.26$ & $5.66$  & $8.78$  \\
$\overline{pen}$              & $0.475$ & $0.569$ & $0.303$ & $0.580$ \\
\tgroup{5}{Prices and aggregates (endogenous)}
$r$ (annual, \%)              & $8.58$  & $6.82$  & $5.89$  & $6.82$  \\
$w$                           & $0.875$ & $0.984$ & $1.053$ & $0.984$ \\
$K$                           & $5.16$  & $10.04$ & $12.82$ & $10.96$ \\
$L$                           & $6.51$  & $9.44$  & $10.18$ & $10.31$ \\
$Y$                           & $9.49$  & $15.48$ & $17.87$ & $16.91$ \\
$K/L$                         & $0.79$  & $1.06$  & $1.26$  & $1.06$  \\
$5K/Y$                        & $2.72$  & $3.24$  & $3.59$  & $3.24$  \\
$C/Y$ (\%)                    & $61.1$  & $63.1$  & $61.0$  & $63.1$  \\
$M/Y$ (\%)                    & $3.21$  & $3.16$  & $2.92$  & $3.04$  \\
$5B/Y$ (\%)                   & $112$   & $107$   & $124$   & $107$   \\
\bottomrule
\end{tabular}
\tabnote{0.94}{Each column is a converged stationary general equilibrium at $\widehat\Theta_{\mathrm{MD}}$ under the baseline closure ($\tau^p$ clears the PAYG block, $B$ residually closes the general budget), at the search discretization; levels carry that caveat, and cross-column comparisons are the reportable objects. $r$ is the annualized pre-tax return; $\tau^p$ the five-year-period contribution rate; $5K/Y$ and $5B/Y$ annual ratios. The imposed rows are exhaustive: everything below them adjusts. In CF1, $\kappa=0.1846$ is set analytically from the PAYG identity and confirmed in equilibrium: $\tau^p$ returns to its baseline value in a single solve (realized $5.660\%$ against the target $5.658\%$). The difference between the CF0 and CF1 dependency ratios---about $1\times10^{-4}$---is the endogenous-survival footprint of the benefit cut. $5B/Y$ is a model-internal outcome of the residual debt closure, disclosed at Table~\ref{tab:baseline_ge}; it is reported in every column and targeted in none.}
\end{table}

\paragraph{Aggregates.} Under inaction the contribution rate more than doubles, from $5.66\%$ to $12.26\%$, and the equilibrium confirms the PAYG identity to the last reported digit: $12.257\% = 0.40 \times 0.3064$. But the same demographic force that doubles the payroll tax also deepens capital. Slower population growth tilts the stationary age distribution toward older, asset-rich cohorts, so the capital--labor ratio rises from $0.79$ to $1.06$; the marginal product of capital falls---the annual return drops from $8.58\%$ to $6.82\%$---and the wage rises $12.4\%$. The two movements nearly cancel in the worker's budget: the net-of-tax wage $w(1-\tau^\omega-\tau^p)$ \emph{rises} about $3.6\%$ despite the doubled contribution rate, and because the flat benefit~\eqref{eq:pension_benefit} is indexed to the wage, the pension rises almost $20\%$. Aging under inaction, in general equilibrium, is not a story of impoverished workers financing unaffordable pensions; it is a story of a doubled tax wedge financed, and slightly more than financed, by capital deepening.

\paragraph{Newborn welfare: almost exactly nothing.} The welfare consequence for a newborn is the sharpest statement of that offset. Table~\ref{tab:cf_welfare} reports the consumption equivalents: a newborn drawn from the birth distribution of the aged economy would give up $0.145\%$ of lifetime consumption before being as badly off as a baseline newborn---an effect economically indistinguishable from zero, with mixed signs across types no larger than $0.43\%$ in absolute value. A newborn holds no assets, so the lower return costs it little; it sells labor at a wage $12.4\%$ higher; and it can re-optimize its entire life cycle around the new prices from the start. Decomposing through an intermediate arm that imposes the aged economy's factor prices at the baseline fiscal bundle shows the anatomy: a substantial price-channel gain to every type, a comparable fiscal loss (the doubled $\tau^p$ net of the higher pension) to every type, and the small residual the table reports as the net; the split is ordering-dependent---prices before fiscal---though neither channel's sign is. What would break the offset is what the model's closure assumes: with an open-economy interest rate, or benefits not indexed to the wage, the price channel would shrink while the fiscal channel stayed, and inaction would cost newborns in full.

The economics beneath the cancellation is plain. When the population grows more slowly, labor becomes scarce relative to capital: the wage rises and the return falls. A newborn faces a doubled contribution rate but sells its labor at the higher wage for a full working life, and because the benefit is indexed to the contemporaneous wage, the pension it will eventually collect rises as well; both sides of its lifetime budget move up together and nearly cancel. The newborn is not protected by policy; it is protected by the factor-price movement that the same demographic change produces. Part of the capital deepening behind that movement is not a savings response at all. Gross investment absorbs $(n_p+\delta)K$ each period, so slower population growth frees resources that were previously equipping new workers---a widening dividend that is a consequence of the stationary-demography device of Section~\ref{subsec:cf_design}, and part of why the aged economy is richer per worker.

\begin{table}[H]
\centering
\caption{Newborn welfare across the menu: consumption equivalents $\lambda$ (\% of lifetime consumption).}
\label{tab:cf_welfare}
\small
\begin{tabular}{@{}lrrrr@{}}
\toprule
 & \thead{CF0 vs.\ 2020} & \thead{CF1 vs.\ CF0} & \thead{CF2 vs.\ CF0} & \thead{CF2 vs.\ CF1} \\
\midrule
Men, low skill     & $-0.31$ & $+8.13$ & $+3.00$ & $-5.39$ \\
Men, high skill    & $+0.23$ & $+9.30$ & $+4.43$ & $-5.28$ \\
Women, low skill   & $+0.43$ & $+7.82$ & $+2.19$ & $-5.82$ \\
Women, high skill  & $-0.01$ & $+8.58$ & $+3.96$ & $-4.97$ \\
\midrule
Aggregate ($\pi_1$-weighted) & $+0.14$ & $+8.29$ & $+3.11$ & $-5.45$ \\
\bottomrule
\end{tabular}
\tabnote{0.92}{Signed De Nardi-direction consumption equivalents~\eqref{eq:cev_newborn}: $\lambda>0$ means the first-named economy's newborn is better off, and $|\lambda|$ is the share of lifetime consumption the better economy's newborn would surrender to be as well off as the other's. Steady-state comparisons of newborns; transition generations are not modeled (see the benefit-margin discussion). Against the 2020 baseline directly, the benefit cut delivers $+8.41\%$ in the aggregate. The aggregate CF0 entry is $0.145$; its numerical bracketing band is $[0.1448,\,0.1450]$.}
\end{table}

\paragraph{Where the cost of inaction actually sits.} Figure~\ref{fig:cf_lambda_cf0} prices the same comparison for residents of every age, and the near-zero newborn number gives way to a steep age gradient: same-age residents of the aged economy are worse off than their baseline counterparts by one to two per cent of remaining consumption at ages 25--29, six to eight per cent at 40--44, and twelve to twenty-one per cent at 65--69. The engine is the intertemporal price, not health: across all ages and types the health stock differs by less than $0.002$ between the two economies---the ceiling result of Section~\ref{subsec:calibration_results} carries over---while the lower return rotates the entire optimal consumption profile toward youth. Residents of the aged economy consume about $7\%$ \emph{more} at ages 20--24 and progressively less with age---roughly $8\%$ less in their fifties, $20\%$ less at 65--69---with asset holdings lower commensurately, compounding at the lower return. The newborn and the age profile thus tell consistent stories: the tilt is the path a newborn \emph{chooses} ex ante at the new prices, so it costs the newborn almost nothing, while a cross-steady-state comparison at age 70 lands on the far end of the chosen tilt and prices it as a large gap. The gradient, not the old-age level, is the finding; the levels at the oldest ages additionally carry the flow-value amplification discussed above.

\begin{figure}[H]
\centering
\caption{The cost of inaction by age: consumption equivalents $\lambda_j$ for same-age steady-state residents, aged economy (CF0) vs.\ the 2020 baseline.}
\label{fig:cf_lambda_cf0}
\includegraphics[width=0.85\textwidth]{Figures/cf_lambda_age_cf0.pdf}
\fignote{0.9}{Signed per-age consumption equivalents: $\lambda_j<0$ means the aged economy's age-$j$ resident is worse off than the baseline's, by $|\lambda_j|$ per cent of remaining consumption. Skill groups: \serone{high skill} and \sertwo{low skill}; solid lines men, dashed lines women. Cross-steady-state resident comparisons, not transition incidence; old-age levels are amplified by the flow-value term (Section~\ref{subsec:welfare_incidence}).}
\end{figure}

\paragraph{Why the incidence differs by type.} The age profiles fan out in a specific order: at every age past fifty the proportional loss is largest for the two high-education types (over $20\%$ at 65--69), intermediate for low-education men ($17\%$), and smallest for low-education women ($12\%$). The ordering is the return channel read through portfolio composition: inaction punishes claims on capital and rewards claims on the pension system (the wage-indexed benefit rises $20\%$), so a resident's exposure is the share of old-age consumption financed from private assets rather than the flat pension. High-education households, the economy's largest savers, hold the most capital-intensive old ages and bear the steepest proportional losses; low-education women, with the fewest private assets and the greatest pension reliance, are cushioned most. The flat benefit thus acts as its progressivity implies: it partially insures precisely the households least able to self-insure through wealth.

The newborn sign pattern in the first column of Table~\ref{tab:cf_welfare}---low-education women the largest gainers, low-education men the only meaningful losers---traces to two measured asymmetries, and we flag its fragility explicitly. First, the pension is flat, so its $20\%$ rise is worth most relative to lifetime resources to the lowest earners---which pushes low-education women to the top of the column. Second, the education--survival gradient enters the staged schedules with different timing by sex. The male gradient opens in the mid-fifties, whereas the female gradient is clipped to zero over ages 50--69 by the survival-cap limitation disclosed in Section~\ref{subsec:firststep} and appears only from age 70 onward, where it is at least as large as the male one. At the estimates the model accordingly delivers educated men an expected age at death about $2.5$ years above their less-educated counterparts, and educated women about $1.9$ years---untargeted model outcomes, not schedule inputs. What matters for the sign pattern is the timing: because their gradient opens roughly two decades earlier, low-education men are differentiated by mortality over the whole pension-relevant range---they pay the doubled contribution for a full working life, collect the raised pension for the fewest years, and lose---whereas female mortality does not differentiate by education until age 70, so the flat benefit's progressivity decides the female ranking alone. The sign pattern is thus the education-mortality input showing through the welfare numbers---and it is a fragile residual: each entry is a fraction of a per cent left over from price and fiscal channels an order of magnitude larger, and the deferred age-varying survival-floor refinement, which would un-clip the female 50--69 gradient, could reorder the female column. We report the pattern; we do not build on it.

\subsection{The Benefit Margin}
\label{subsec:exp_kappa_taum}

Holding the 2050 contribution rate at its 2020 value requires cutting the replacement rate from $0.40$ to $\widehat\kappa=0.1846$---a $54\%$ benefit cut, and the analytic value $\tau^p_{2020}/0.3064$ restores the contribution rate in a single step (realized $5.660\%$ against the $5.658\%$ target). The number itself is the first result: the identity offers no soft option, and any politically feasible partial cut buys proportionally partial relief.

\paragraph{Aggregates and the Aaron logic.} The benefit cut is the strongest capital-deepening force on the menu. With the public pension less than half as generous ($\overline{pen}$ falls $47\%$ relative to inaction), households self-insure old age through private saving: capital rises from $10.04$ to $12.82$ across the two 2050 economies, the wage gains a further $7\%$ on top of inaction's rise, and the annual return falls to $5.89\%$. For a newborn of \emph{every} type this trade is decisively favorable: $8.3\%$ of lifetime consumption against inaction ($7.8$--$9.3\%$ across types; $8.4\%$ against the 2020 baseline). The economics is the classic Aaron condition \citep{Aaron1966}. A stationary PAYG pillar pays an implicit return equal to population growth---$1.5\%$ per year in the 2050 economy---while the market pays $5.9\%$; forcing households to route retirement claims through the low-return pillar is a standing tax on their life-cycle saving, and shrinking the pillar refunds it. General equilibrium amplifies rather than dampens the direct incentive, since the induced capital raises the wage newborns sell into. The result would flip where the condition fails---an economy near dynamic inefficiency ($r$ near $n$), or one where binding borrowing constraints keep the young from doing the substitute saving; neither obtains at the estimates.

\paragraph{The caveat that must not be buried.} These gains are steady-state gains, and this experiment is where the restriction to steady states bites hardest. The $8.3\%$ accrues to newborns \emph{of the reformed steady state}, who never contributed at $\kappa=0.40$. The transition generations---workers who contributed for decades at the old rate and would retire under $\kappa=0.185$---appear nowhere in this comparison, and their losses are exactly what a transition analysis of an unfunded system's wind-down prices. The steady-state ranking says what the destination is worth, not what the journey costs, and the two must not be conflated: nothing in Table~\ref{tab:cf_welfare} licenses the claim that cutting benefits is a Pareto improvement.

\paragraph{Age incidence: a sign flip ordered by pension dependence.} The closest steady-state object to that transition intuition is the age profile, panel~(a) of Figure~\ref{fig:cf_lambda_menu}. Support for the benefit cut falls monotonically with age, crosses zero in mid-life, and turns severely negative for retirees: newborns and young workers gain $7$--$9\%$, and 65--69-year-olds are worse off by $4$--$21\%$ of remaining consumption, the losses deepening through the oldest ages. A resident's stake is the ratio of already-accumulated pension claims to the years of low-$\tau^p$ wages still ahead, so the crossing age orders by pension dependence---and the cross-type ordering confirms the mechanism exactly. Low-education women, the most pension-reliant type, flip first (by ages 45--49) and lose the most at 65--69 ($-21\%$); low-education men flip around 50--54 and lose $13\%$; high-education women flip by 55--59; and high-education men, the economy's best-buffered savers, flip last (60--64) and lose least ($-3.7\%$). The mirror image of inaction's incidence is exact: inaction taxes private wealth through the return, so the asset-rich lose most; the benefit cut taxes pension claims, so the pension-reliant lose most. The two policies do not differ mainly in how much welfare they deliver---they differ in \emph{whose} old age bears the demographic bill.

\begin{figure}[H]
\centering
\caption{The reform menu by age: consumption equivalents $\lambda_j$ for same-age steady-state residents, each reform vs.\ inaction (CF0) at 2050 demographics.}
\label{fig:cf_lambda_menu}
\includegraphics[width=\textwidth]{Figures/cf_lambda_age_menu.pdf}
\fignote{0.94}{Signed per-age consumption equivalents vs.\ inaction: $\lambda_j>0$ means the reform economy's age-$j$ resident is better off. Panel (a): the benefit cut ($\kappa\colon 0.40\to0.1846$). Panel (b): retirement at 70. Skill groups: \serone{high skill} and \sertwo{low skill}; solid lines men, dashed lines women. Cross-steady-state resident comparisons, not transition incidence; old-age levels carry the flow-value amplification (Section~\ref{subsec:welfare_incidence}). In panel (b), the 65--69 losses of the low-skill types are likely understated: the late-life efficiency input is estimated from the selected minority still working at those ages.}
\end{figure}

\subsection{The Retirement Margin}
\label{subsec:cf_retirement}

Raising the retirement age to 70 moves the 65--69 band from beneficiaries to contributors, and the PAYG identity again does the fiscal arithmetic in advance: the dependency ratio falls to $0.219$ and the contribution rate lands at $8.78\%$---almost exactly the predicted $0.40\times0.219$---recovering just over half of the fiscal cost of aging while leaving the benefit rule untouched. The pension per retiree actually edges \emph{up} relative to inaction (more contributors per beneficiary at an unchanged replacement rate on a higher wage base).

\paragraph{A scale expansion, not a price shock.} What distinguishes this reform in general equilibrium is what does not move. Against inaction, the return is unchanged to three digits ($6.819$ vs.\ $6.820\%$) and the wage to within five parts in $10^{5}$. Extending working lives raises effective labor by $9\%$, and the extra labor income induces almost exactly the capital that keeps the capital--labor ratio---and hence both factor prices---where inaction put them: the economy simply scales up ($Y$ rises from $15.48$ to $16.91$) and the fiscal burden falls. Where inaction and the benefit cut operate through prices, the retirement margin operates through quantities plus the budget: fiscal relief at constant prices.

\paragraph{Newborn ranking.} For newborns the reform is worth $+3.1\%$ of lifetime consumption relative to inaction---but $-5.5\%$ relative to the benefit cut, and the ranking $\mathrm{CF1}\succ\mathrm{CF2}\succ\mathrm{CF0}$ holds for every one of the four types (Table~\ref{tab:cf_welfare}). The reason it forfeits most of the benefit cut's gain is the Aaron logic again: working longer \emph{preserves} the unfunded pillar at $\kappa=0.40$, so the newborn keeps routing a large claim through the low-return system, and pays five more years of contributions and work disutility for the privilege of the smaller wedge. The two reforms are not close substitutes for a newborn; they are a large benefit-side gain versus a moderate quantity-side gain.

\paragraph{Incidence: the extended work band pays.} Panel~(b) of Figure~\ref{fig:cf_lambda_menu} shows an incidence pattern unlike either other experiment: not a monotone age gradient but a dip, bottoming exactly at the extended work band. Low-education women lose $9.5\%$ of remaining consumption at 65--69 and are in losses from ages 50--54 onward---the approaching work requirement is already priced into their remaining-lifetime value---while low-education men dip to $-1.3\%$ at 65--69 only. High-education residents of \emph{every} age gain ($+4$ to $+8\%$), and retirees 70 and older all gain (a slightly higher pension, no work requirement). The heterogeneity is the late-life efficiency input read through the wage: at 65--69 the estimated efficiency of high-education men and women ($1.80$ and $1.39$ efficiency units) still commands substantial earnings, so their extended years are well paid and the fiscal relief dominates; low-education women, at $0.64$---the lowest cell in the profile---are made to work five more years at wages that barely compensate the disutility, and low-education men ($0.88$) sit in between. The burden of a higher retirement age falls on exactly the workers whose late-life market productivity is lowest.

Two facts qualify these losses, in opposite directions. They are steady-state numbers, so the disruption to cohorts near 65 at a reform date is not in them. And they are likely \emph{understated}: the 65--69 efficiency input comes from the selected fifth of that age group that chooses to work, whereas the reform compels the other four fifths---including those who stopped precisely because their market productivity or health made work unattractive---so the marginal worker's efficiency is below the estimate, and the true 65--69 losses, especially the low-education women's $-9.5\%$, are a lower bound in magnitude.

\subsection{Behavior and Equilibrium: A Decomposition}
\label{subsec:cf_decomposition}

Each experiment's welfare effect can be split into what households' own re-optimization delivers at unchanged aggregates and what arrives only through equilibrium. We re-solve the household problem under each experiment's direct changes while holding the entire price-and-fiscal bundle $(r, w, \tau^p, \overline{pen})$ at baseline values---a partial-equilibrium arm in which no market clears, disclosed as such---and then let an intermediate arm impose the experiment's factor prices while keeping the baseline fiscal bundle. The three differences that result attribute the total newborn-welfare change to a direct behavioral component, a factor-price channel, and a fiscal channel; the finer split prices the price channel before the fiscal one and is ordering-dependent, though no sign in Table~\ref{tab:cf_decomposition} is. The decomposition is exact in units of lifetime value, in which the arms are additive; consumption equivalents are not additive across arms, so the table reports value units and the reader should take from it signs and proportions rather than magnitudes.

\begin{table}[H]
\centering
\caption{Decomposition of the newborn welfare effect: behavior versus equilibrium.}
\label{tab:cf_decomposition}
\begin{tabular}{lcccc}
\toprule
 & Direct (behavioral) & Factor prices & Fiscal & Total \\
\midrule
CF0 (inaction)        & $0$ (structural) & $+3.247$ & $-3.180$ & $+0.067$ \\
CF1 (benefit cut)     & $0$ (structural) & $+5.017$ & $-1.035$ & $+3.982$ \\
CF2 (retirement age)  & $-0.432$         & $+3.310$ & $-1.348$ & $+1.530$ \\
\bottomrule
\end{tabular}
\vspace{0.5em}
\begin{minipage}{0.92\textwidth}
\footnotesize \textit{Notes.} Birth-weighted newborn lifetime value, all arms measured against the 2020 baseline (level $340.60$); rows sum by construction. Direct arm: the experiment's non-price changes at the baseline $(r, w, \tau^p, \overline{pen})$; factor-price channel: the experiment's $(r,w)$ at baseline fiscal variables, net of the direct arm; fiscal channel: the remainder. Budgets do not clear inside the arms. The split is ordering-dependent (prices before fiscal); no entry's sign is.
\end{minipage}
\end{table}

The first two rows put numbers on the structural fact of the design (Section~\ref{subsec:cf_design}): because neither $n_p$ nor $\kappa$ enters the household problem, the direct effect of aging and of the benefit cut is identically zero, and everything those experiments do to anyone is done by markets and the pension system. For inaction the table restates the newborn-indifference result as channel arithmetic---a price channel worth $+3.25$ against a fiscal channel worth $-3.18$, cancelling to the $+0.067$ residual whose smallness Section~\ref{sec:aging} priced at $0.145\%$ of consumption---and it locates the benefit cut's entire $+3.98$ in the same two channels: the crowd-in of private saving deepens capital beyond what demography alone achieves ($+5.02$ through prices, against $-1.03$ through a 2050 fiscal bundle whose contribution rate matches the baseline by construction but whose pension is less than half). The Aaron gain of Section~\ref{subsec:exp_kappa_taum} arrives through the price system, not through the fiscal ledger.

The third row is the reason the decomposition earns its place. The retirement-age reform is the only experiment that reaches households directly---and its direct effect is a welfare \emph{loss}. Compelled to work through ages 65--69 at unchanged prices and fiscal variables, the newborn is worse off by $0.43$: five additional years of work disutility, sold at late-life efficiencies of $0.64$--$1.80$, are not worth their wages when nothing else improves. Every unit of the reform's value is general-equilibrium: the scale expansion holds factor prices at inaction's values while the added labor income widens the tax base and shrinks the wedge, and the equilibrium channels ($+3.31$ through prices, $-1.35$ through the 2050 fiscal bundle) must first repay the direct loss before the reform nets its $+1.53$. A planner who read the retirement margin as a household-level instruction---work longer, be better off---would have its sign wrong; the reform works only because everyone works longer at once, and the incidence of Section~\ref{subsec:cf_menu}'s third regularity follows: the band that supplies the direct loss, low-education workers in their late sixties, is the band the equilibrium gains fail to compensate.

\subsection{The Policy Menu}
\label{subsec:cf_menu}

The three experiments assemble into a menu that one identity organizes. At 2050 demographics the government chooses which variable in $\tau^p=\kappa\,N^R/N^W$ gives: the contribution rate ($12.26\%$), the replacement rate ($\widehat\kappa=0.185$, a $54\%$ cut), or the retirement boundary (age 70, $\tau^p=8.78\%$). Inaction is itself the tax arm---financing aging with a payroll-tax increase is a policy, chosen by default---and the menu prices all three.

Three regularities summarize the section. First, \emph{the fiscal arithmetic is exact and the welfare arithmetic is not}. The PAYG identity predicted every contribution rate and the required benefit cut to within rounding before any equilibrium was solved; nothing in the identity predicted that inaction would leave newborns unharmed ($+0.14\%$), that the benefit cut would be worth $8.3\%$ to them, or that the retirement margin would deliver little more than a third of that---those numbers come from the general-equilibrium price responses, above all from capital deepening. Second, \emph{for newborns the ranking is unambiguous}: the benefit cut dominates the retirement-age rise, which dominates inaction, for every type---the Aaron condition, applied to a demographic environment where the market return exceeds population growth by more than four percentage points, rewards whichever margin most shrinks the unfunded pillar. Third, \emph{the reforms differ most in who pays}. Inaction taxes wealth-holders through the return: its losses rise with age and are steepest for the high-education savers. The benefit cut taxes pension claims: its losses concentrate on retirees, ordered by pension dependence, worst for low-education women. The retirement-age rise taxes the extended work band: its losses concentrate on low-education workers in their late sixties, again worst for low-education women, while every other group gains. A planner reading this menu is not choosing a welfare level so much as choosing an incidence---and for the two groups that recur at the bottom of every incidence table, low-education workers and pension-dependent retirees, the ranking of the two reforms is reversed.

The second and third regularities are one fact read at two ages. Every margin on the menu reprices a stock that somebody already holds and cannot rebuild. Inaction lowers the return, which reprices the remaining lifetime of anyone holding assets, so it falls hardest on older savers. The benefit cut lowers the pension, which reprices the claims of those already retired and dependent on them. The retirement-age rise extends the work band, falling on those whose late-life productivity is lowest and who must now work at ages when they would not have. Newborns rank the three margins identically because a newborn holds no stock---no assets, no accrued pension claim---and can adjust to any reform from the start; the rankings diverge with age precisely because older residents cannot. The choice among the three is therefore a choice of whose old age pays, not a choice of how much welfare to have.

\subsection{Caveats}
\label{subsec:exp_caveats}

Five caveats bound what this section claims. First, every comparison is between stationary equilibria. The transition path---and with it the losses of the generations who contributed under the old rules, the object that dominates the political economy of pension reform---is not computed, and the benefit-margin discussion states where this bites hardest. General-revenue financing of the pension gap (letting $\tau^\omega$ or debt absorb the shortfall rather than $\tau^p$) would require a different budget architecture than the self-contained PAYG block of Section~\ref{subsec:government} and is left as a model extension. Second, all equilibria are computed at the search discretization; levels carry that caveat, and the cross-economy comparisons the section is built on are the reportable objects. Third, the stationary debt ratio is a model-internal residual of the closure in every column ($107$--$124\%$ of annual output against roughly $50\%$ for M\'exico); it is disclosed at the baseline, reported here, and targeted nowhere. Fourth, the demographic experiment moves fertility only: survival schedules and the birth composition are held at 2020, for the endogeneity reason given in the design, so the aging imposed here understates any 2050 mortality improvements; and within the staged schedules the female education--survival gradient is clipped to zero over ages 50--69 (it survives at older ages), which the inaction heterogeneity discussion flags where it matters. Fifth, the slate contains no health-policy experiment: at the estimates the health stock sits against its behaviorally defended ceiling for nearly every household (Section~\ref{subsec:calibration_results}), so a health-side instrument acts only through the withdrawal direction, and pricing that experiment fairly awaits the health-block specification decision of Section~\ref{subsec:roadmap}.


% ======================================================================


%% ====================================================================
%% SECTION 8 — CONCLUSION
%% ====================================================================
\section{Conclusion}
\label{sec:conclusion}

This paper builds a stochastic overlapping-generations model with continuous endogenous health capital, heterogeneous agents who differ by sex and skill, endogenous factor prices, and a pay-as-you-go pension system whose contribution rate clears endogenously, and calibrates it to M\'exico. The sex-by-age labour-efficiency profile and the skill premium are estimated from the pooled ENIGH cross-sections, the health-depreciation schedule and the sex- and education-specific survival schedules and floor from ENASEM, and the remaining nine preference, health-technology, health-risk, and health--productivity parameters by classical minimum distance on ten moments, nine of them estimated from Mexican micro-data. That estimation has been run, and the model's own general-equilibrium solve at the estimates is the paper's current baseline, at which the policy experiments reported here are computed. One moment an earlier solve could not match---the value of a statistical life---exposed that a felicity whose zero coincides with the value of death cannot price survival, and the specification now carries a flow value of being alive, estimated rather than targeted. That result, and the model's VSL reported as an output against an external benchmark, now have their own home in Section~\ref{subsec:vsl_structural} and are not retold here.

The paper's health contribution is an apparatus and a measurement made with it. The apparatus is a general-equilibrium economy in which a single stock of health capital simultaneously drives survival, labour productivity, and the amenity value of consumption, with prices, the contribution rate, and the pension benefit clearing endogenously around it, and with the health technology disciplined by classical minimum distance against Mexican micro-moments rather than assumed or imported. It solves and clears: the baseline of Section~\ref{subsec:calibrated_baseline} is its own converged equilibrium at the estimates. The measurement is what that discipline reveals (Section~\ref{subsec:capacity_result}). At parameter values that fit the spending profile, both income elasticities of medical spending, the skill premium, and the survival anchors at once, the investment block generates essentially none of the observed education--survival gradient: the endogenous investment-to-longevity channel contributes a small, disclosed fraction of the education gap in mid-life survival, the measured schedules carry the rest, and the estimator pushes the auxiliary channels---the amenity and risk margins---to the corners of their admissible ranges. This is a finding about the apparatus, not a shortfall of the calibration. It could not have been obtained from a partially calibrated model: with any moment released, the estimator would have found a configuration that hides the tension, and only the full discipline forces the collision.

The measurement has a diagnosis, and it is about objects rather than parameters (Section~\ref{sec:healthblock}). Medical spending in the model is committed before the health shock is realized, and is therefore a preventive object by construction: it builds the stock in advance of illness, on a margin every household operates every period. Measured Mexican out-of-pocket spending, on the evidence assembled in the same section, is predominantly restoration conditioned on a realized draw---care sought after illness arrives. A block containing only the first object, disciplined on spending moments generated largely by the second, resolves the tension the only way its structure permits: it holds the health stock near its behaviorally defended ceiling for nearly every household, where committed spending is high enough to match the data while health variation is too small for a spending-driven survival gradient to open. The contribution is the limit itself: the Grossman-descended investment block, taken to emerging-economy micro-data and disciplined on the spending, elasticity, earnings, and survival facts at once, meets a boundary that only a fully disciplined estimation exposes---the spending it is asked to rationalize is by majority not the kind of spending it models.

Three substantive results stand out, all computed at the estimated baseline. The first is that demographic aging under unchanged policy is a redistribution across cohorts rather than an aggregate welfare event: the equilibrium contribution rate more than doubles, from $5.66\%$ to $12.26\%$, and yet newborn welfare is essentially unchanged, because slower population growth deepens capital and the induced wage and pension gains offset the wedge; the cost is borne by cohorts already holding assets, whose remaining lifetimes are repriced at a lower return. The second is that the aggregate capital stock of the calibrated economy is governed first by the persistence of idiosyncratic income risk and only second by skill composition---the empirical income process mean-reverts, and buffer-stock saving is correspondingly weak. The third is that the two pension margins differ less in aggregate welfare than in incidence: holding the 2050 contribution rate at its 2020 value requires cutting the replacement rate from $0.40$ to $0.185$ and is worth $8.3\%$ of lifetime consumption to newborns of the reformed steady state, with losses concentrated on pension-dependent retirees, while raising the retirement age to seventy recovers about half the fiscal cost at essentially unchanged factor prices, with losses concentrated on low-education workers in their late sixties. The steady-state comparisons price destinations only: the transition generations of a benefit cut are not modeled, and nothing here licenses a Pareto claim.

The experimental program is organized by one identity. Because the pay-as-you-go budget ties the contribution rate to the product of the replacement rate and the dependency ratio, a government facing a doubling dependency ratio must give up one of three things---the contribution rate, the replacement rate, or the boundary between contributors and beneficiaries---and the three experiments of Section~\ref{sec:experiments} move exactly one each. What the identity cannot deliver, and general equilibrium does, is the welfare and incidence content: which margin is worth most to a newborn, and whose old age pays for it. That incidence content is where the health apparatus earns its place in a paper about pension margins, even though the welfare engine of the experiments is the intertemporal price rather than health: the incidence results are structured throughout by the estimated survival and efficiency profiles. Low-education men emerge as the only meaningful newborn losers under inaction because their education--survival gradient opens roughly two decades earlier than women's, so they pay the doubled contribution over a full working life and collect the raised, wage-indexed pension for the fewest years---a sign pattern the paper reports with its fragility explicitly flagged. Low-education women bear the retirement-age reform because their estimated late-life efficiency is the lowest of the four types, so the extended work band compels them to work at wages that barely compensate the disutility. Neither statement is available in a model without sex- and education-specific survival and efficiency schedules; the apparatus that cannot yet generate the survival gradient endogenously is the same apparatus whose measured inputs decide who pays for pension reform.

For the fiscal authority this paper was written for, four statements follow from the results and go no further. First, the menu is closed within the pension system: the identity has three arguments and no others, and everything outside them either relocates the obligation to a different budget constraint---general-revenue financing, which the paper's closure does not price---or changes the structure of the benefit rule, which is a different model. Second, the three margins differ far more in whose old age pays than in how much welfare the economy has; the choice among them is a distributional decision presented as a fiscal one, and this paper prices the distribution. Third, under a flat benefit indexed to the contemporaneous wage and a fixed retirement age, improvements in population health raise the contribution rate; they do not lower it. The identity contains no productivity term through which better health could pay for itself, so this is a property of the benefit rule, not of health---and it means the fiscal case for health investment cannot be made through the pension block as currently designed. Fourth, the model's dependency ratio is demographic rather than contributor-based and the baseline is calibrated to the formal-sector population, so the fiscal pressure reported here is the favourable case: the contribution-base margin that informality opens is outside this arithmetic. The paper ranks the margins for newborns and shows the ranking reverses with age; it recommends none of them.

What the paper does not establish is stated with equal care. No transition path is computed, so nothing here is incidence on any cohort alive at a reform date: every welfare number compares residents of different stationary economies, prices destinations rather than journeys, and licenses no Pareto claim about any actual reform. The economy has no informality margin---all baseline workers are formal---and no annuity market: decedents' assets are voided rather than bequeathed or annuitized, a closure disclosed where it matters. All equilibria are computed at the search discretization, so cross-economy comparisons, not levels, are the reportable objects. And the health block's inability to generate the education--survival gradient is a measured capacity result awaiting a specification decision, not a defect that qualifies the fiscal findings: the incidence of aging and reform operates through factor prices, the contribution rate, and the measured survival and efficiency schedules, the endogenous health response is negligible across the compared economies, and the fiscal results do not depend on the health channel being active.

The path forward is recorded in Section~\ref{subsec:roadmap} and compresses to five items: the specification decision on the health block---whether to relocate the mechanism to old-age deterioration, re-specify repair so that mid-life health mass sits below the ceiling, or reframe the contribution around the channels the estimation supports---which is the binding item; the transition analysis, which would price the journeys the steady-state comparisons cannot; the formalization programme, which would open the contribution-base margin held fixed here; the health-side instruments, which await the specification decision; and the extension to Costa Rica and Panam\'a. One earlier reservation has since been settled into the baseline itself. Under the deterministic law of motion the model's working-age health distribution is nearly degenerate, so the income-to-health mechanism the paper is built around would operate on a variable with almost no dispersion; stochastic health with deficit-scaled repair (Section~\ref{subsubsec:health_risk}) is now the baseline the estimation disciplines, for the evidence given in Section~\ref{subsec:health_block_evidence}.

The principle that organizes the analysis is straightforward: the question facing Latin American health systems is not whether to spend more, but when to spend, through which institutional channels, and with what consequences for the contributory architecture on which fiscal sustainability depends---and the measurement at the center of this paper is that principle made operational, because when the spending occurs, before illness or in response to it, is exactly what separates the model's object from the data's.


%% ====================================================================
%% APPENDICES
%% ====================================================================
\begin{appendices}

\section{Model Appendix: Derivations and Optimality Conditions}
\label{app:model}

\subsection{The Amenity: Multiplicative versus Additive}
\label{app:amenity}

Section~\ref{subsec:preferences} states that the multiplicative amenity creates an income gradient in the willingness to pay for health that an additive amenity lacks by construction. This appendix records the derivations behind that claim.

Scaling the GHH aggregator by $\Omega(h)$ makes the marginal utility of consumption depend on health,
\begin{equation}
\label{eq:u_ch}
u_{ch} \;=\; \Omega'(h)\,\bigl(c-v(\ell)\bigr)^{-\gamma} \;\geq\; 0,
\end{equation}
with strict inequality whenever $\Xi > 0$. The sign holds at every state, since $\Omega' \geq 0$ and $(c - v(\ell))^{-\gamma} > 0$ irrespective of the bracket in~\eqref{eq:utility}. Health is a state rather than a within-period choice, so this complementarity reaches the medical-spending decision indirectly: through $u_h$ in the health envelope condition~\eqref{eq:env_h}, which supplies the marginal continuation value in~\eqref{eq:foc_m}, and through the dependence of next period's $u_c$ on $h'$ in the consumption Euler equation~\eqref{eq:euler}.

Had the amenity entered the composite additively, as $u = \bigl[(c + s(h) - v(\ell))^{1-\gamma} - 1\bigr]/(1-\gamma)$ with $s' > 0$, the same cross-partial would carry the opposite sign,
\[
u_{ch} \;=\; -\gamma\, s'(h)\, \bigl(c + s(h) - v(\ell)\bigr)^{-\gamma-1} \;<\; 0 .
\]
The sign of $u_{ch}$ is not, however, the object that governs behavior. Medical spending responds to the marginal rate of substitution between health and consumption, which is what enters the medical-spending condition~\eqref{eq:foc_m}. Under the additive form that ratio is
\[
\frac{u_h}{u_c} \;=\; \frac{s'(h)\,\bigl(c+s(h)-v(\ell)\bigr)^{-\gamma}}{\bigl(c+s(h)-v(\ell)\bigr)^{-\gamma}} \;=\; s'(h),
\]
independent of consumption. Health and consumption are perfect substitutes inside the composite, so the willingness to pay for health carries no income gradient; the negative $u_{ch}$ above is the concavity of the outer aggregator and by itself moves nothing. Under the multiplicative form the same ratio is
\[
\frac{u_h}{u_c} \;=\; \frac{\Omega'(h)}{\Omega(h)}\left[\frac{z^{1-\gamma}-1}{1-\gamma} + \chi\right] z^{\gamma}, \qquad z \equiv c - v(\ell),
\]
which is increasing in $c$ whenever $1 + \gamma\bigl[(z^{1-\gamma}-1)/(1-\gamma) + \chi\bigr]z^{\gamma-1} > 0$. The condition holds at every state at which the bracket in~\eqref{eq:utility} is nonnegative, and can fail only inside the region where that bracket is negative, that is, at states the household ranks below not being alive.

The two forms also differ at the boundary of the feasible set. Because the amenity multiplies the aggregator, $u_c = \Omega(h)\,(c-v(\ell))^{-\gamma} \to \infty$ as $c \downarrow v(\ell)$, so consumption satisfies an Inada condition. Under an additive amenity the singularity sits instead at $c = v(\ell) - s(h)$, which is strictly negative whenever $s(h) > v(\ell)$. The Inada point is then unreachable, $u_c$ stays finite as $c \to 0$, and a household can be close to indifferent between consuming and not.

\subsection{Household Optimality Conditions}
\label{app:focs}

This appendix states the first-order, envelope, and Kuhn--Tucker conditions of the worker's problem~\eqref{eq:worker_bellman} and the retiree's problem~\eqref{eq:retiree_bellman}. Let $\lambda_j$ denote the multiplier on the budget constraint, and recall that the two nodes $h'_k$ of the health lottery~\eqref{eq:health_lom_shock} are functions of $(m, h)$ through $\tilde h$. At an optimum at which the $\min$ and $\max$ operators in~\eqref{eq:health_lom_shock} are slack at both nodes, $\partial h'_k / \partial m = \partial H_j / \partial m$ and $\partial h'_k / \partial h = \partial \tilde h / \partial h$ for both $k$, with the derivatives given in~\eqref{eq:dh_dh}. The interior first-order conditions are
\begin{align}
\text{Consumption:} \quad & u_c(c, \ell, h) \;=\; \lambda_j (1 + \tau^c), \label{eq:foc_c}\\
\text{Labor:} \quad & -u_\ell(c, \ell, h) \;=\; \lambda_j\, w\, \nu_j(h, \eta;\, g, \theta)\,(1-\tau^\omega-\tau^p), \label{eq:foc_l}\\
\text{Savings:} \quad & \lambda_j \;=\; \beta \sum_{k \in \{0,1\}} \Pr(k \mid j)\; \psi_{j+1}^{g,\theta}(h'_k)\, \E\bigl[V_{a, j+1}(a', h'_k, \eta';\, g, \theta) \mid \eta\bigr], \label{eq:foc_a}\\
\text{Medical spending:} \quad & \lambda_j (1 + \tau^m)
\;=\; \beta\, \frac{\partial H_j(m,h)}{\partial m} \sum_{k \in \{0,1\}} \Pr(k \mid j)\Bigl\{
\bigl(\psi_{j+1}^{g,\theta}\bigr)'(h'_k)\, \E[V_{j+1}(a', h'_k, \eta';\, g, \theta) \mid \eta] \notag\\
&\qquad\qquad\qquad\qquad +\, \psi_{j+1}^{g,\theta}(h'_k)\, \E[V_{h, j+1}(a', h'_k, \eta';\, g, \theta) \mid \eta]
\Bigr\}. \label{eq:foc_m}
\end{align}
Condition~\eqref{eq:foc_m} states the two returns to medical spending that the model keeps distinct---a survival return, through $(\psi^{g,\theta}_{j+1})'$ multiplying the level of the continuation value, and a health-capital return, through $\psi^{g,\theta}_{j+1}$ multiplying its marginal value in health---each averaged over the two nodes of the health lottery. Because $\varsigma$ shifts the struck node down by a constant, the marginal product of spending is common to the two nodes and factors out; what differs across nodes is the point at which the survival schedule and the continuation value are evaluated. If the floor $\underline h$ binds at the struck node, that node's health no longer responds to $m$ and the $k=1$ term drops from the right-hand side of~\eqref{eq:foc_m}.

\paragraph{Envelope conditions.} For the worker's problem ($j < j_R$),
\begin{align}
V_{a, j}(a, h, \eta;\, g, \theta) &\;=\; \lambda_j \cdot \bigl(1 + r(1-\tau^k)\bigr), \label{eq:env_a}\\
V_{h, j}(a, h, \eta;\, g, \theta) &\;=\; u_h(c, \ell, h) + \lambda_j\, w\, \nu_j(h, \eta;\, g, \theta)\, \varrho\, \ell\, (1-\tau^\omega-\tau^p) \notag\\
&\qquad +\, \beta\, \frac{\partial \tilde h}{\partial h} \sum_{k \in \{0,1\}} \Pr(k \mid j)\Bigl\{ \bigl(\psi_{j+1}^{g,\theta}\bigr)'(h'_k)\, \E[V_{j+1}(a', h'_k, \eta';\, g, \theta) \mid \eta] \notag\\
&\qquad\qquad\qquad\qquad\qquad +\, \psi_{j+1}^{g,\theta}(h'_k)\, \E[V_{h, j+1}(a', h'_k, \eta';\, g, \theta) \mid \eta] \Bigr\}, \label{eq:env_h}
\end{align}
with $\partial \tilde h / \partial h$ from~\eqref{eq:dh_dh}. The second term is the wage return to health. Its slope $\varrho$ is common to the two skill groups, so any cross-group difference in the earnings value of health operates through the wage bill $w\,\nu_j\,\ell$---through the intercept $\theta$ and the profile $e^{g,\theta}_j$---rather than through a differential return: health is worth more, in earnings, to workers whose time is worth more. For the retiree's problem ($j \geq j_R$), $\ell = 0$ and that term vanishes:
\begin{align}
V_{a, j}^{R}(a, h, \eta;\, g, \theta) &\;=\; \lambda_j \cdot \bigl(1 + r(1-\tau^k)\bigr), \label{eq:env_a_ret}\\
V_{h, j}^{R}(a, h, \eta;\, g, \theta) &\;=\; u_h(c, 0, h) \notag\\
&\qquad +\, \beta\, \frac{\partial \tilde h}{\partial h} \sum_{k \in \{0,1\}} \Pr(k \mid j)\Bigl\{ \bigl(\psi_{j+1}^{g,\theta}\bigr)'(h'_k)\, \E[V_{j+1}(a', h'_k, \eta';\, g, \theta) \mid \eta] \notag\\
&\qquad\qquad\qquad\qquad\qquad +\, \psi_{j+1}^{g,\theta}(h'_k)\, \E[V_{h, j+1}(a', h'_k, \eta';\, g, \theta) \mid \eta] \Bigr\}. \label{eq:env_h_ret}
\end{align}
At the terminal age $j = J$, $\psi^{g,\theta}_{J+1} \equiv 0$ kills the continuation terms; the envelopes remain well defined since $a_J' = m_J = 0$.

\paragraph{Euler equation.} Combining~\eqref{eq:foc_a} with~\eqref{eq:env_a} one period ahead and~\eqref{eq:foc_c},
\begin{equation}
\label{eq:euler}
u_c(c, \ell, h) \;=\; \beta\, \bigl(1 + r(1-\tau^k)\bigr) \sum_{k \in \{0,1\}} \Pr(k \mid j)\; \psi_{j+1}^{g,\theta}(h'_k)\, \E\bigl[u_c(c'_k, \ell'_k, h'_k) \mid \eta\bigr],
\end{equation}
where $(c'_k, \ell'_k)$ are next-period policies evaluated at the node-$k$ health state.

\paragraph{Consumption--leisure condition.} Combining~\eqref{eq:foc_c} and~\eqref{eq:foc_l},
\begin{equation}
\label{eq:cons_leisure}
\frac{-u_\ell(c, \ell, h)}{u_c(c, \ell, h)}
\;=\;
\frac{w\, \nu_j(h, \eta; g, \theta)(1-\tau^\omega-\tau^p)}{1 + \tau^c},
\end{equation}
where the amenity cancels from the left-hand side, since $\Omega(h)$ scales $u$ uniformly: $-u_\ell/u_c = v'(\ell)$.

\paragraph{Kuhn--Tucker for the borrowing constraint.}
The constraint $a' \geq 0$ generates a non-negative multiplier and turns the savings condition~\eqref{eq:foc_a} into the weak inequality
\begin{equation}
\label{eq:kt_borrowing}
\lambda_j \;\geq\; \beta \sum_{k \in \{0,1\}} \Pr(k \mid j)\; \psi_{j+1}^{g,\theta}(h'_k)\, \E\bigl[V_{a, j+1}(a', h'_k, \eta';\, g, \theta) \mid \eta\bigr], \qquad a' \geq 0,
\end{equation}
with complementary slackness: the product of the gap and $a'$ is zero.

\paragraph{Kuhn--Tucker for medical spending.}
The constraint $m \geq 0$ turns~\eqref{eq:foc_m} into the weak inequality
\begin{equation}
\label{eq:kt_medical}
\lambda_j (1 + \tau^m)
\;\geq\;
\beta\, \frac{\partial H_j(m,h)}{\partial m} \sum_{k \in \{0,1\}} \Pr(k \mid j)\Bigl\{
\bigl(\psi_{j+1}^{g,\theta}\bigr)'(h'_k)\, \E[V_{j+1}(\cdot) \mid \eta]
+ \psi_{j+1}^{g,\theta}(h'_k)\, \E[V_{h, j+1}(\cdot) \mid \eta]
\Bigr\},
\end{equation}
for $m \geq 0$, with complementary slackness. The cap $h'_k \leq \bar h$ in~\eqref{eq:health_lom_shock} is slack wherever $\overline{H}_j m^{\zeta_h} < \bar h$, by the asymptote property of~\eqref{eq:health_production_deficit}; where it binds it introduces a non-negative multiplier $\mu_{\bar h} \geq 0$ with complementary slackness
\begin{equation}
\label{eq:kt_healthcap}
\mu_{\bar h}\,(\bar h - h'_0) \;=\; 0, \qquad h'_0 \leq \bar h, \qquad \mu_{\bar h} \geq 0,
\end{equation}
entering the medical-spending condition as a subtracted term $\mu_{\bar h}\, \partial H_j/\partial m$ on the right-hand side of~\eqref{eq:kt_medical}. At an interior optimum with $m > 0$ and both nodes strictly inside $(\underline h, \bar h)$, all inequalities hold with equality and~\eqref{eq:foc_m} is recovered.

\section{Fiscal Gap Methodology}
\label{app:fiscal_gap}

The fiscal gap for country $c$ in year $t$ is defined as:
\begin{equation}
\label{eq:fiscal_gap}
G_{c,t} = H_{c,t}^{\text{est}} - H_{c,t}^{\text{actual}}
\end{equation}
where $H_{c,t}^{\text{est}}$ denotes estimated health expenditure requirements and $H_{c,t}^{\text{actual}}$ denotes actual public health expenditure, both expressed as shares of GDP. Estimated requirements are constructed by applying a base-year age-expenditure profile $h(j)$ to the projected population structure:
\begin{equation}
\label{eq:fiscal_gap_est}
H_{c,t}^{\text{est}} = \frac{\sum_j h(j) \cdot N_{j,c,t}}{Y_{c,t}}
\end{equation}
where $N_{j,c,t}$ is the population of age group $j$ in country $c$ at time $t$, and $Y_{c,t}$ is nominal GDP. The base-year expenditure profile $h(j)$ is estimated from national health accounts microdata for the reference year.


\section{2050 Demographic Inputs: Derivation and Sources}
\label{app:demog2050}

\paragraph{Source.} The demographic target for the 2050 experiments of Section~\ref{sec:experiments} is the old-age dependency ratio $65{+}/(20\text{--}64)$ from CONAPO's official population projections for M\'exico \citep[Proyecciones de la Poblaci\'on de M\'exico y de las Entidades Federativas, 2023 revision;][]{CONAPO2023}: $0.3065$ in 2050, against the 2020 census value of $0.1414$ that disciplines the baseline (Section~\ref{subsec:firststep}). The bracketing projections ($0.1746$ in 2030, $0.5064$ in 2070) locate 2050 on a still-steepening path. The 2050 ratio comes from the same CONAPO release as the life tables used in the first step, so the model's demographic inputs are single-sourced. The UN World Population Prospects is the natural robustness alternative and changes only the target number.

\paragraph{Method.} The experiment imposes the 2050 population \emph{structure}, not a demographic forecast path, and it does so through the single parameter the baseline already uses for the same purpose. In the stationary population, cohort weights satisfy $\phi_{j+1}/\phi_j = \bar\psi_j/(1+n_p)^5$, where $\bar\psi_j$ is the distribution-weighted average survival of band $j$---an equilibrium object, since survival is health-modulated. We extract the baseline's realized $\bar\psi_j$ schedule from its recorded stationary distribution (the extraction reproduces the baseline's recorded dependency ratio to $10^{-12}$ as a precondition), hold it fixed, and bisect $n_p$ until the implied stationary dependency ratio equals the CONAPO 2050 target. This yields $n^{*}_{2050}=1.483\%$ per year, against the baseline's $3.51\%$. Because $\bar\psi_j$ is invariant to $n_p$ only at fixed household policies, the bisected value is a candidate, not the number of record: the confirming general-equilibrium solve at $n^{*}_{2050}$ realizes a dependency ratio of $0.30643$, within $0.024\%$ of the target (tolerance $0.5\%$), so no re-extraction round was needed. The realized ratio is the number of record, and the PAYG identity evaluated at it reproduces the reported contribution rate exactly.

\paragraph{What is deliberately not moved.} Survival schedules stay at their 2020 estimates. Model survival is partly endogenous---the measured baseline schedules are modulated by the health stock through~\eqref{eq:survival_modulation}---so an imposed ``2050 survival schedule'' would import the behavioral responses embedded in the projection into what the experiment treats as an exogenous environment change; a survival-shift arm (the natural robustness exercise) would first require a defensible way to strip the projection to its exogenous component. Fertility decline is the dominant driver of M\'exico's projected aging, so the fertility channel carries the experiment. The birth composition $\pi^{g,\theta}_1$ likewise stays at 2020: a 2050 education-composition arm would conflate aging with skill upgrading and needs an external projection of educational attainment. Both choices are disclosed in Section~\ref{sec:experiments} where results depend on them.


\section{Calibration and Computational Appendix}
\label{app:calibration}

This appendix documents the apparatus behind the Mexican-data calibration of Section~\ref{sec:calibration}: the objective and weighting, the identification analysis, the multistart convergence record, the parity and Euler diagnostics, the override mechanics that let the calibration run without editing the protected solver, and the first-step provenance.

\subsection{Objective and Weighting}
\label{app:smm_objective}

The estimator minimizes the quadratic form~\eqref{eq:smm_objective}, $Q(\Theta)=[m^{\rm mod}(\Theta)-m^{\rm dat}]'\,W\,[m^{\rm mod}(\Theta)-m^{\rm dat}]$. Under $W=\operatorname{diag}(1/\mathrm{SE}^2)$ each moment contributes its squared $t$-statistic and $Q(\widehat\Theta)=\sum_k t_k^2$, which is the form used in the superseded six-moment estimation reported in Appendix~\ref{app:provenance} and in Table~\ref{tab:moment_match}. Section~\ref{subsec:secondstep_smm} explains why that weighting is treated as a specification choice rather than a default under the current ten-moment target set. At each candidate $\Theta$ the model is solved as a partial equilibrium at prices anchored to a stationary general-equilibrium solve, and the model moments are computed \emph{exactly} from the stationary distribution $\Phi$ over $(j,g,\theta,a,h,\eta)$ rather than by simulating panels---which is what makes the estimator classical minimum distance and removes the $(1+1/S)$ term from the variance formula. Asymptotic standard errors on $\widehat\Theta$ are sandwich estimates from a numerical Jacobian of $m^{\rm mod}$ at $\widehat\Theta$; because $\Xi$ enters the amenity $\Omega(h)$ that scales the GHH aggregator, its identification interacts with the externally fixed $\gamma$, which is treated as conditioning information.

\paragraph{A caveat on the Jacobian.} The finite-difference step size is not innocuous here. The derivative of the moment that currently dominates the objective---the age log-slope of medical spending---is not sign-stable across candidate step sizes, so a Jacobian computed at a default step can invert the sign of the dominant column and with it every standard error derived from the sandwich formula. A step-size sensitivity study is a prerequisite to reporting inference, not a robustness check on it. The underlying reason is the discretization noted in Section~\ref{subsec:secondstep_smm}: every moment is kinked where a household changes grid cell, and the two income-elasticity moments are additionally discontinuous where a cell crosses the spending floor at which it leaves the regression. Note that the moment whose derivative misbehaves here---the age log-slope---is \emph{not} one of the two, so the step-size problem is not a consequence of that discontinuity and is not resolved by removing it.

\subsection{Identification}
\label{app:identification}

\emph{This analysis was performed on the superseded six-parameter, six-moment system of Appendix~\ref{app:provenance}. It does not speak to the current nine-parameter, ten-moment system, whose identification is discussed in Section~\ref{subsec:secondstep_smm} and whose rank condition has not been re-analysed at the reported estimates.} The pre-optimization Jacobian at the starting values (Figure~\ref{fig:jacobian_elast}, reported as target-relative elasticities) shows that every free parameter has strong sensitivity on at least one moment, with no globally flat direction. Two near-flat directions require handling. The first is the $\Xi\leftrightarrow\xi$ ridge on the VSL target: both the amenity level and its curvature move \texttt{vsl\_usd} with the same sign and comparable magnitude, so they cannot be separately identified from that single moment. It was resolved there by freezing $\xi=0.5$ during estimation. The second is the $\overline{H}_0\leftrightarrow\zeta_h$ ridge on the medical-spending moments: the scale and curvature of the health-production function compete with opposite signs, $\zeta_h$ roughly three times the stronger, so their point estimates are within literature ranges but imprecise (the wide standard errors in Table~\ref{tab:theta_hat}). This ridge is acknowledged rather than resolved; it does not prevent identification of $\overline{H}_0$ to within range. One benign off-diagonal sign-flip (the within-age elasticity with respect to $h^{\rm slope}$) appears at the boundary $h^{\rm slope}=0$, where the sign is fragile to the income profile; it is not a new identification problem.

\begin{figure}[H]
\centering
% TODO: ensure 02b_jacobian_elasticity.png is uploaded to the Overleaf project root
% (source: Calibration/outputs/diagnostics/02b_jacobian_elasticity.png; 02_jacobian.png is the alternative).
\includegraphics[width=0.85\textwidth]{02b_jacobian_elasticity.png}
\caption{Pre-optimization Jacobian: target-relative elasticities of the six moments with respect to the six parameters. The $\Xi\leftrightarrow\xi$ ridge on \texttt{vsl\_usd} (resolved by freezing $\xi$) and the strong $\zeta_h$ column on the medical-spending moments (the acknowledged $\overline{H}_0\leftrightarrow\zeta_h$ ridge) are the two identification features behind the estimation restrictions of Section~\ref{subsec:secondstep_smm}.}
\label{fig:jacobian_elast}
\end{figure}

\subsection{Multistart Protocol and Convergence}
\label{app:multistart}

The objective is minimized by a multistart Nelder--Mead simplex: four starts of at most forty iterations each, run under four-way contention at roughly twelve minutes per partial-equilibrium evaluation. The smoke test at the starting values returns $Q(\theta_{\rm init})=73.11$ with all six moments finite. The best of the four starts is start~4 at $Q=13.28$; the remaining starts converge to $Q\approx 20.5$, so the reported optimum is well separated from the runner-up basins. Figure~\ref{fig:obj_trace} traces the running-best objective across evaluations on a log scale.

\begin{figure}[H]
\centering
% TODO: regenerate 03_objective_trace.png for the Mexican-data multistart (best Q = 13.28)
% and upload to the Overleaf project root
% (source: Calibration/outputs/diagnostics/03_objective_trace.png reflects the intermediate run).
\includegraphics[width=0.9\textwidth]{03_objective_trace.png}
\caption{Multistart progress: running-best objective $Q$ versus evaluation index, on a log scale. The descent terminates at the best-of-four optimum $Q=13.28$.}
\label{fig:obj_trace}
\end{figure}

\subsection{Parity and Euler Diagnostics}
\label{app:parity_euler}

The calibration scaffold reproduces the production \texttt{Household-Gender} solver to machine precision: the maximum absolute cohort-mean difference between the two solvers is below $10^{-10}$ (achieved at $7\times10^{-11}$), so the calibration and the experiments share an identical household solution algorithm. At the calibrated baseline (Run~1$'$) the Euler-equation residual has a maximum of $\log_{10}=-5.37$ and a mean of $-7.42$, with the mass-weighted bulk of the distribution comfortably below the convergence gate; the two-sex welfare-gap diagnostic also passes. The calibrated experiments (Runs~2--4) each converged within twenty-six outer iterations and pass the same internal gates.

\subsection{Override Mechanics}
\label{app:overrides}

The calibration and the experiments ran without editing the protected solver \texttt{ge\_model\_gender.jl}. Three override patterns were used, each respecting Julia's compilation semantics. The AR(1) persistence and innovation scale $(\rho,\sigma_\varepsilon)$ are \texttt{const} scalars inlined into the Rouwenhorst discretization, so they were passed as \emph{literals} into the model initializer rather than redefined (a \texttt{const} redefinition would silently fail to propagate into already-compiled methods). The birth-share vector $\pi_{\rm birth}$ is a \texttt{const} array, so in-place mutation propagates, as do the ENOE age-efficiency and skill-intercept overrides. The 2050 demographics are installed on typed globals (population growth and the sex-specific survival schedules) through a setter rather than by inlining. A frozen-parameter bug in the diagnostics module was worked around through a shim rather than a source edit. These patterns are recorded for replicability.

\subsection{First-Step Provenance of the Superseded Solve}
\label{app:provenance}

This subsection documents the first step behind the superseded estimates of Table~\ref{tab:theta_hat}, not the current one. The current first-step inputs are the ENIGH and ENASEM estimates of Section~\ref{subsec:firststep}.

The survival schedule is the WHO Global Health Observatory life table for M\'exico (2019), aggregated to five-year bands and closed out above age~85 by a Gompertz fit, consistent with UN WPP 2024 medium-variant magnitudes. The ENOE age-efficiency profile $e^g_j$ is constructed, for each sex and five-year band, as the skill-averaged log labour income net of the skill intercept, normalised so the male 20--24 cell equals one; the skill intercepts follow from the prime-age (35--49) cross-skill log gap of $0.745$. The hours target is the equal-weighted mean of the ENOE \texttt{hrsocup} share over male prime-age (30--44) skill cells. The ENOE-based inputs described here have been superseded and the file that documented them is no longer part of the project.

The stationary equilibrium that earlier solve generates---the economy the experiments of Section~\ref{sec:experiments} perturb---is reported in Table~\ref{tab:calibrated_baseline}. It is recorded here, rather than beside the current baseline of Table~\ref{tab:baseline_ge}, so that the two are not mistaken for one another.

At that superseded parameterization the economy supplies effective labour $L=4.99$ and holds capital $K=4.96$ per unit of cohort mass, producing output $Y=7.96$ at an annual interest rate of $4.21\%$. The pay-as-you-go pension system clears at a contribution rate of $\tau^p=9.29\%$ against a system dependency ratio of $0.186$, and households devote $48\%$ of output to consumption and $2.7\%$ to medical care.

\begin{table}[H]
\centering
\caption{Calibrated baseline equilibrium (M\'exico, 2020): stationary aggregates. \emph{Superseded solve} (p.~\pageref{par:superseded}); not the current baseline of Section~\ref{subsec:calibrated_baseline}.}
\label{tab:calibrated_baseline}
\small
\begin{tabular}{@{}lr@{}}
\toprule
\theadl{Quantity} & \theadr{Value} \\
\midrule
\tgroup{2}{Stationary aggregates and prices}
$K$ (capital, per cohort mass)              & $4.96$ \\
$L$ (effective labour)                      & $4.99$ \\
$Y$ (output)                                & $7.96$ \\
$r$ (annual)                                & $4.21\%$ \\
$w$ (wage)                                  & $1.022$ \\
$\tau^p$ (PAYG contribution rate)           & $9.29\%$ \\
$\overline{pen}$ (pension flow per retiree) & $0.365$ \\
$C/Y$                                       & $48.1\%$ \\
$M/Y$                                       & $2.73\%$ \\
$N^R/N^W$ (PAYG dependency)                 & $0.186$ \\
$5K/Y$ (annual capital--output ratio)       & $3.11$ \\
$5B/Y$ (annual debt--output ratio)          & $165\%$ \\
\bottomrule
\end{tabular}
\tabnote{0.92}{Stationary per-cohort-mass units (Section~\ref{subsec:markets}). Stock--flow ratios are annualized by the factor of five of Section~\ref{subsec:units}: $5K/Y$ and $5B/Y$ convert the five-year-flow denominator to an annual one, and $r$ is annualized geometrically. Closure: $\tau^p$ clears the PAYG block, and the stationary debt $B=2.62$ residually closes the general budget through~\eqref{eq:gov_stationary} (so $B/Y=0.329$ at the five-year frequency, $5B/Y=165\%$ annual). Welfare $\mathcal{W}_1(g,\theta)$ is expected lifetime utility at birth ($\gamma=2$); because GHH--CRRA felicity is not homogeneous in consumption, only the ordering across types is interpreted, never the levels. Euler-equation residuals: $\max\log_{10}=-5.37$, mean $-7.42$.}
\end{table}

The macroeconomic ratios in the table constitute a genuine, if modest, validation \emph{of that superseded configuration}, because none of them is targeted: the discount factor, capital share, and depreciation rate are set externally, and the estimated moments concern hours and medical spending. The annual capital--output ratio of $5K/Y=3.11$ is in line with Mexican data, the annual return of $4.2\%$ is in the empirically plausible range, and a consumption share of $48\%$ is reasonable for an economy with a $19\%$ public-goods share and life-cycle investment. One untargeted ratio is off, and we flag it rather than suppress it: the implied public debt, at $165\%$ of annual GDP, far exceeds M\'exico's federal debt (near $50\%$). The excess is a consequence of the residual closure---$B$ absorbs whatever stationary stock balances the primary account at the calibrated tax-and-spending configuration---rather than of any targeted object, and disciplining it directly, by targeting a debt ratio and letting a tax instrument clear the budget, is a refinement for the next stage. The number matters beyond bookkeeping, because the debt level enters the infeasibility result of Section~\ref{subsec:c2_infeasibility}. None of this transfers to the current baseline, whose sourced capital share, depreciation rate, and tax block differ and whose corresponding ratios are those of Table~\ref{tab:baseline_ge}. The ratios above are provenance for the superseded solve, not evidence for the model as now calibrated.

\begin{table}[H]
\centering
\caption{Superseded estimates $\widehat\Theta$ under the earlier specification and first step (best of four starts, $Q=13.28$).}
\label{tab:theta_hat}
\begin{tabular}{lrrrr}
\toprule
Parameter & Init & $\widehat\theta$ & Rel.\ change & SE \\
\midrule
$\Psi$ (labor disutility)                  & 14.00  & 13.45    & $-3.9\%$  & 0.232 \\
$\Xi$ (health amenity)                     & 0.500  & 0.275    & $-45.1\%$ & 0.171 \\
$\xi$ (amenity curvature)                  & 0.500  & 0.500    & frozen    & ---   \\
$\overline{H}_0$ (health-production scale) & 0.300  & 0.247    & $-17.5\%$ & 0.949 \\
$h^{\mathrm{slope}}$ (age-decline)         & 0.000  & $-0.0358$ & ---      & 0.138 \\
$\zeta_h$ (health-production curvature)    & 0.500  & 0.579    & $+15.8\%$ & 0.895 \\
\bottomrule
\end{tabular}

\vspace{0.5em}
\begin{minipage}{0.92\textwidth}
\footnotesize \textit{Notes.} $\widehat\theta$ is the best of four multistart Nelder--Mead runs ($Q=13.28$ at start~4; the other starts converge to $Q\approx 20.5$). $\xi$ is held fixed at $0.5$ to resolve the $\Xi\leftrightarrow\xi$ ridge. Standard errors are sandwich estimates from a numerical Jacobian at $\widehat\theta$; the wide $\overline{H}_0$ and $\zeta_h$ intervals reflect the medical-spending ridge documented in Appendix~\ref{app:calibration}.
\end{minipage}
\end{table}

\begin{table}[H]
\centering
\caption{Moment match at $\widehat\theta$.}
\label{tab:moment_match}
\begin{tabular}{lrrlrr}
\toprule
Moment                         & Data         & SE        & Source                & Model     & $t$-stat \\
\midrule
Prime-age male hours           & 0.2686       & 0.000946  & ENOE (this paper)     & 0.2687    & $+0.07$  \\
VSL (2022 USD)                 & 2{,}000{,}000 & 700{,}000 & OECD transfer (Lit.)  & 27{,}698  & $-2.82$  \\
Cross-section $\eta_m$         & 1.000        & 0.300     & Hall--Jones (Lit.)    & 0.471     & $-1.76$  \\
Mean $m$, ages 25--35          & 0.040        & 0.020     & Hall--Jones; ENIGH    & 0.0131    & $-1.35$  \\
$d\log\bar m/dj$, $j\in[2,12]$ & 0.120        & 0.050     & Literature            & 0.136     & $+0.33$  \\
Within-age $\eta_m$            & 0.400        & 0.150     & Hall--Jones; RAND     & 0.483     & $+0.56$  \\
\bottomrule
\end{tabular}

\vspace{0.5em}
\begin{minipage}{0.92\textwidth}
\footnotesize \textit{Notes.} Source column distinguishes the ENOE-this-paper moment from the retained literature/external anchors. Five of six moments are within $2$ SE (three within $1$ SE); the VSL is the structural divergence of Section~\ref{subsec:vsl_structural}. $\sum_k t_k^2 = 13.28 = Q(\widehat\theta)$.
\end{minipage}
\end{table}

\label{subsec:repro}
Two facts about the superseded results should be on the record rather than left to be discovered. First, the superseded estimates of Table~\ref{tab:theta_hat}, the moment fit of Table~\ref{tab:moment_match} and the baseline equilibrium of Table~\ref{tab:calibrated_baseline} have no producing artifact in the project repository: they are internally consistent---the reported objective equals the sum of squared $t$-statistics---but they cannot be regenerated from what is archived, and the general-equilibrium solve that is archived differs materially from them. Second, the fixed price anchor at which the estimation is carried out is taken from that archived general-equilibrium run, and is therefore \emph{not} the equilibrium reported in Table~\ref{tab:calibrated_baseline}: the anchor carries a contribution rate of $14.5\%$ and an annual interest rate of $4.80\%$ against the $9.29\%$ and $4.21\%$ reported there. The claim in Section~\ref{subsec:secondstep_smm} that the anchor lies near the general-equilibrium point is therefore true of one equilibrium of this model and not of the one the paper reports, and the bound it places on the partial-equilibrium approximation error should be read accordingly. Both are reasons to treat the superseded numbers as provenance for the present draft rather than as results, and both are resolved by the re-estimation and the general-equilibrium solve at its estimates that this draft now reports; what remains is the regeneration of the baseline exhibits and experiments from it (Section~\ref{subsec:roadmap}).

\section{Supplementary Tables and Data}
\label{app:supplementary}

This appendix collects the three-country comparative landscape that motivates the project. Only M\'exico is calibrated in the present draft; these indicators discipline none of the parameters in Section~\ref{sec:calibration} and are reported for context.

\subsection{Demographic and Health-Status Indicators}
\label{app:demog_table}

\begin{table}[H]
\centering
\caption{Demographic and Health Status Indicators}
\label{tab:demog}
\begin{tabular}{lccc}
\hline\hline
\textbf{Indicator} & \textbf{Mexico} & \textbf{Costa Rica} & \textbf{Panama} \\
\hline
Life expectancy at birth, total (2023)           & 75.1 & 81.0 & 78.1 \\
Life expectancy, male (2023)                     & 72.1 & 78.6 & 75.4 \\
Life expectancy, female (2023)                   & 78.3 & 83.5 & 80.9 \\
Fertility rate (births per woman, 2023)          & 1.82 & 1.51 & 2.30 \\
Infant mortality (per 1{,}000 live births, 2022) & 12.4 & 8.7  & 14.5 \\
Population aged 65+ (\% of total, 2023)          & 8.1  & 10.4 & 8.9  \\
Diabetes prevalence (\% of adults)               & 16.9 & 9.8  & 11.0 \\
\hline\hline
\end{tabular}
\begin{minipage}{\linewidth}
\smallskip
\footnotesize
\textit{Sources:} \citet{PAHOHIA2024}; \citet{WHOGHO2024};
\citet{WorldBankWDI2024}; \citet{OECDHaG2023}.
\end{minipage}
\end{table}

\subsection{Health Expenditure Indicators}
\label{app:health_table}

\begin{table}[H]
\centering
\caption{Health Expenditure Indicators (Reference Year: 2022)}
\label{tab:health_exp}
\begin{tabular}{lccc}
\hline\hline
\textbf{Indicator} & \textbf{Mexico} & \textbf{Costa Rica} & \textbf{Panama} \\
\hline
Current health expenditure (\% of GDP)       & 5.5          & 7.1   & 7.8   \\
Public health expenditure (\% of CHE)        & $\approx$51  & $\approx$69 & $\approx$72 \\
Out-of-pocket expenditure (\% of CHE)        & $\approx$41$^{a}$ & 22.4 & $\approx$26 \\
OOP per capita, USD current                  & $\approx$253 & 220   & $\approx$290 \\
Total health exp.\ per capita, USD PPP       & $\approx$1{,}100 & 1{,}658 & $\approx$1{,}450 \\
OOP as \% of household final consumption     & $\approx$5.0 & 2.6   & $\approx$3.5 \\
\hline\hline
\end{tabular}
\begin{minipage}{\linewidth}
\smallskip
\footnotesize
\textit{Sources:} \citet{WHOGHED2024}; \citet{OECDHaG2023};
\citet{WHOGHO2024}. CHE = Current Health Expenditure; OOP =
Out-of-pocket; PPP = Purchasing Power Parity. Values marked
($\approx$) are approximations. $^{a}$~Mexico's OOP share
corresponds to the 2021 data vintage, the most recent available.
\end{minipage}
\end{table}

\subsection{Fiscal Indicators}
\label{app:fiscal_table}

\begin{table}[H]
\centering
\caption{Fiscal Indicators}
\label{tab:fiscal}
\begin{tabular}{lccc}
\hline\hline
\textbf{Indicator} & \textbf{Mexico} & \textbf{Costa Rica} & \textbf{Panama} \\
\hline
Standard VAT rate                                  & 16\% (IVA) & 13\% (IVA) & 7\% (ITBMS) \\
Top personal income tax rate (2023)                & 35\%  & 25\%  & 25\%  \\
Social security contributions (employer + employee) & $\approx$30\% & $\approx$37\% & $\approx$19\% \\
General government revenue (\% of GDP, 2022)       & $\approx$23 & $\approx$24 & $\approx$21 \\
Tax revenue (\% of GDP, 2022)                      & $\approx$17 & $\approx$14 & $\approx$12 \\
\hline\hline
\end{tabular}
\begin{minipage}{\linewidth}
\smallskip
\footnotesize
\textit{Sources:} \citet{PwCWWTS2024}; \citet{OECDRevStats2024};
\citet{WorldBankWDI2024}; \citet{CCSS2023}.
\end{minipage}
\end{table}

\subsection{Labour Market Indicators}
\label{app:labour_table}

\begin{table}[H]
\centering
\caption{Labour Market Indicators (Reference Year: 2022)}
\label{tab:labour}
\begin{tabular}{lccc}
\hline\hline
\textbf{Indicator} & \textbf{Mexico} & \textbf{Costa Rica} & \textbf{Panama} \\
\hline
Male labour force participation rate (\%)    & 76.7        & 72.4        & 77.6        \\
Female labour force participation rate (\%)  & 45.4        & 50.7        & 51.8        \\
Gender participation gap (male/female ratio) & 1.69        & 1.43        & 1.50        \\
Gender wage gap (\% of male wage)            & $\approx$14 & $\approx$18 & $\approx$19 \\
Informality rate (\% of employed)            & $\approx$55 & $\approx$43 & $\approx$47 \\
Unemployment rate (\%)                       & 3.3         & 11.2        & 7.7         \\
Gini coefficient                             & 0.43        & 0.48        & 0.49        \\
\hline\hline
\end{tabular}
\begin{minipage}{\linewidth}
\smallskip
\footnotesize
\textit{Sources:} \citet{ILOILOSTAT2023}; \citet{ECLACILO2023};
\citet{IMF2023}; \citet{WEFGGR2024}; \citet{WorldBankWDI2024};
\citet{WorldBankPovCalNet2022}.
\end{minipage}
\end{table}

\subsection{Life Table for Mexico}
\label{app:survival}

The following life table corresponds to Mexico and is used to
calibrate the age-specific survival probabilities $\psi_{j}$ in the
model. Data are drawn from \citet{UISP2019}.

\begin{table}[H]
\centering
\caption{Life Table for Mexico: Mortality and Survival Probabilities}
\label{tab:survival_mx}
\begin{tabular}{lcc}
\hline\hline
\textbf{Age group} & \textbf{Probability of death} & \textbf{Survival probability} \\
\hline
$<$ 1 year   & 0.011576 & 0.9884 \\
1--4         & 0.002282 & 0.9977 \\
5--9         & 0.001132 & 0.9989 \\
10--14       & 0.001470 & 0.9985 \\
15--19       & 0.003858 & 0.9961 \\
20--24       & 0.005897 & 0.9941 \\
25--29       & 0.006550 & 0.9935 \\
30--34       & 0.007966 & 0.9920 \\
35--39       & 0.011217 & 0.9888 \\
40--44       & 0.013710 & 0.9863 \\
45--49       & 0.019216 & 0.9808 \\
50--54       & 0.028413 & 0.9716 \\
55--59       & 0.042568 & 0.9574 \\
60--64       & 0.063162 & 0.9368 \\
65--69       & 0.093160 & 0.9068 \\
70--74       & 0.137294 & 0.8627 \\
75--79       & 0.208926 & 0.7911 \\
80--84       & 0.312093 & 0.6879 \\
85+          & 1.000000 & 0.0000 \\
\hline\hline
\end{tabular}
\begin{minipage}{\linewidth}
\smallskip
\footnotesize
\textit{Source:} \citet{UISP2019}. Net probability of death and
survival probability by age group for the total population of Mexico.
\end{minipage}
\end{table}

\end{appendices}


%% ====================================================================
%% BIBLIOGRAPHY
%% ====================================================================
\bibliography{references}

\end{document}
